% %%%%%%%%%%%%%%%%%%%%%%%%%%%%%%%%%%%%%%%%%%%%%%%%%%%%%%%%%%%%%%%%%%%%%%%%%%%%%%%%%%%%%%%%%%%%%%%%%%%%%%%%%%%%%%%%%
% %% Kapitel 6 - CaseStudy %%%%%%%%%%%%%%%%%%%%%%%%%%%%%%%%%%%%%%%%%%%%%%%%%%%%%%%%%%%%%%%%%%%%%%%%%%%%%%%%%%%%%%%%
% %%%%%%%%%%%%%%%%%%%%%%%%%%%%%%%%%%%%%%%%%%%%%%%%%%%%%%%%%%%%%%%%%%%%%%%%%%%%%%%%%%%%%%%%%%%%%%%%%%%%%%%%%%%%%%%%%

\chapter{Industrielle Fallstudie: Empirische Validierung der Extrapolationsabsicherung} \label{chap:casestudy}

\section{Von der stochastischen Simulation zur physischen Lebensdauererprobung} \label{sec:case_intro}

Die Ausführungen in Kapitel~\ref{chap:entwurf} erbringen auf Basis umfangreicher \acp{MCS} den theoretischen und algorithmischen Nachweis, dass das \ac{EmALT}-Design im Rahmen der multivariaten, effizienten Lebensdauererprobung signifikante informationstheoretische Vorteile gegenüber klassisch orthogonalen Versuchsplänen bietet.
Für den finalen Transfer dieser Methodik in die industrielle Zuverlässigkeitstechnik dient das vorliegende Kapitel der ergänzenden empirischen Validierung an einem real vermessenen Hardwaredatensatz.
Diese Fallstudie ist dabei bewusst als praxisnahe Ergänzung konzipiert und untermauert die numerischen Ergebnisse des vorherigen Kapitels durch die Betrachtung realer Bauteilversuche an Zahnriemen für Lenksysteme in automobilen Anwendungen \cite{Arndt.2026}.

Der methodische Übergang von der Simulation zur physischen Erprobung impliziert zwei wesentliche Verschärfungen der Randbedingungen:
Erstens weicht die idealisierte Gesetzmäßigkeit unendlicher Simulationsiterationen der harten statistischen Realität stark limitierter Stichprobenumfänge.
Die nachfolgende Untersuchung basiert auf einem realen Versuchsdatensatz mit lediglich zweifacher Replikation und damit $\sym{N} = 26$ Prüflingen.
Zweitens unterliegen experimentelle \ac{EoL}-Versuche zwingenden physikalischen und normativen Restriktionen.
In der Praxis können hochgradig beschleunigte Versuchspunkte (wie die extremen \acp{SP} im nordöstlichen Parameterraum) oftmals nicht wie theoretisch geplant abgefahren werden.
Zu hohe Stressniveaus induzieren hierbei typischerweise eine Überlastung der Prüftechnik, provozieren unrepräsentative Ausfallmechanismen (bspw. Zahnabriss anstatt abrasiven Verschleißes) oder sind physikalisch schlichtweg nicht umsetzbar.

Anstatt diese technologischen Restriktionen als Mangel zu betrachten, generieren sie aus Sicht der adaptiven Versuchsplanung ein zwingendes \textbf{geometrisches Gestaltungsbudget} nach Abschnitt~\ref{sec:adaptive_planung}.
Die nicht realisierbaren Hochlast-Versuche werden aus der Design-Matrix gestrichen.
Anstatt den Versuchsplan dadurch jedoch lediglich passiv auszudünnen und eine asymmetrische Varianzaufweitung in Kauf zu nehmen, wird das freigewordene Budget proaktiv umgewidmet: Die Versuchspunkte werden in die direkte Peripherie des anvisierten Feldbetriebspunktes im Niedrigstressbereich re-allokiert.
Entscheidend ist hierbei, dass die Reallokation algorithmisch optimiert nach den in Kapitel~\ref{chap:entwurf} entwickelten \ac{EmALT}-Kriterien in Kombination mit realen Prüfstandsbeschränkungen erfolgen kann.
Tabelle~\ref{tab:case_ieee_dataset} fasst die experimentell ermittelten Ausfallzeiten des untersuchten Lenksystem-Zahnriemens für die jeweiligen Faktorkombinationen übersichtlich zusammen.

\input{tables/tab_ieee_dataset.tex}

Da diese \acp{AP} im Gegensatz zur theoretischen Simulation nun mit realen Ausfallzeiten vorliegen, ermöglicht dieser Datensatz einen einzigartigen, zweiseitigen Modellvergleich: Die Parameterschätzung und Lebensdauerprognose erfolgt separat für das klassische Basis-\ac{CCD} (ohne \acp{AP}, isoliert aus den Untersuchungen von \textcite{Arndt.2026}) sowie für das re-allokierte \ac{EmALT}-Design inklusive der Ankerpunkte.
Dies erlaubt nicht nur die finale Evaluation der realen Modellanpassung, der \ac{SPV} und der verursachten Versuchskosten, sondern wirft vor allem eine fundamentale Validierungsfrage auf:
Unterscheidet sich das durch die Ankerpunkte gestützte Zuverlässigkeitsmodell lediglich graduell in seiner Prognosevarianz vom Basis-Modell, oder erzwingen die Realdaten auf dem Extrapolationspfad ein signifikant abweichendes, neues Degradationsgesetz?

\section{Empirische Modellschätzung und Anpassungsgüte} \label{sec:case_evaluation}

Die \ac{MLE}-basierte Auswertung der Modelparameter nach Gleichung~\ref{eq:log_likelihood_explicit} offenbart für beide Design-Varianten zunächst signifikante Diskrepanzen in der Gewichtung der physikalischen Effekte.
Abbildung~\ref{fig:ma_abb6.1_parameter_estimates} veranschaulicht die Gegenüberstellung der geschätzten Regressionskoeffizienten in $\hat{\sym{beta}}$ sowie des korrespondierenden Weibull-Formparameters $\hat{\sym{b}}$.

\begin{figure}[htbp]
    \centering
    \setlength\figurewidth{0.9\textwidth}
    \setlength\figureheight{5.5cm}
    % This file was created by matlab2tikz.
%
\definecolor{mycolor1}{rgb}{0.12941,0.12941,0.12941}%
%
\begin{tikzpicture}

\begin{axis}[%
width=0.951\figurewidth,
height=\figureheight,
at={(0\figurewidth,0\figureheight)},
scale only axis,
bar shift auto,
xmin=0.51,
xmax=7.49,
xtick={1.00,2.00,3.00,4.00,5.00,6.00,7.00},
xticklabels={{$\hat{\sym{b}}$},{$\hat{\sym{beta}}_0$},{$\hat{\sym{beta}}_1$},{$\hat{\sym{beta}}_2$},{$\hat{\sym{beta}}_{12}$},{$\hat{\sym{beta}}_{11}$},{$\hat{\sym{beta}}_{22}$}},
ymin=-5.00,
ymax=10.38,
ylabel style={font=\color{mycolor1}},
ylabel={Schätzwert},
axis background/.style={fill=white},
xmajorgrids,
ymajorgrids,
legend style={legend cell align=left, align=left},
grid=major,
grid style={dotted, gray!50},
y tick label style={/pgf/number format/precision=0}
]
\addplot[ybar, bar width=0.229, fill=white!60!black, draw=mycolor1, area legend] table[row sep=crcr] {%
1.00	2.39\\
2.00	9.94\\
3.00	-2.31\\
4.00	0.69\\
5.00	0.22\\
6.00	-0.01\\
7.00	0.53\\
};
\addplot[forget plot, color=mycolor1] table[row sep=crcr] {%
0.51	0.00\\
7.49	0.00\\
};
\addlegendentry{Basis-\ac{CCD}}

\addplot[ybar, bar width=0.229, fill=black, draw=mycolor1, area legend] table[row sep=crcr] {%
1.00	2.11\\
2.00	10.38\\
3.00	-3.31\\
4.00	1.84\\
5.00	1.39\\
6.00	-1.13\\
7.00	-0.12\\
};
\addplot[forget plot, color=mycolor1] table[row sep=crcr] {%
0.51	0.00\\
7.49	0.00\\
};
\addlegendentry{\ac{EmALT}}

\addplot [color=black, only marks, forget plot]
 plot [error bars/.cd, y dir=both, y explicit, error bar style={line width=0.5pt}, error mark options={line width=0.5pt, mark size=6.0pt, rotate=90}]
 table[row sep=crcr, y error plus index=2, y error minus index=3]{%
0.86	2.39	0.90	0.90\\
1.86	9.94	0.30	0.30\\
2.86	-2.31	0.25	0.25\\
3.86	0.69	0.25	0.25\\
4.86	0.22	0.29	0.29\\
5.86	-0.01	0.29	0.29\\
6.86	0.53	0.29	0.29\\
};
\addplot [color=black, only marks, forget plot]
 plot [error bars/.cd, y dir=both, y explicit, error bar style={line width=0.5pt}, error mark options={line width=0.5pt, mark size=6.0pt, rotate=90}]
 table[row sep=crcr, y error plus index=2, y error minus index=3]{%
1.14	2.11	0.00	0.00\\
2.14	10.38	0.00	0.00\\
3.14	-3.31	0.00	0.00\\
4.14	1.84	0.73	0.73\\
5.14	1.39	0.73	0.73\\
6.14	-1.13	0.52	0.52\\
7.14	-0.12	0.30	0.30\\
};
\end{axis}
\end{tikzpicture}%
    \caption{Vergleich der \ac{MLE}-Parameterschätzungen für das klassische Basis-\ac{CCD} und das \ac{EmALT}-Design am realen Datensatz. Die Integration der Ankerpunkte führt zu einer Neukalibrierung der dominanten Haupt- und Krümmungseffekte.}
    \label{fig:ma_abb6.1_parameter_estimates}
\end{figure}

Es wird ersichtlich, dass die Integration der Ankerpunkte das Regressionsmodell spürbar korrigiert. Während das Basis-\ac{CCD} aufgrund seiner massiven Konzentration im Hochlastbereich zu einer leichten Überschätzung isolierter Stressgradienten neigt, "erden" die realen Laufzeiten der \acp{AP} die Koeffizienten in der Peripherie.
Diese numerische Verschiebung der Parameter manifestiert sich visuell prägnant in der rekonstruierten Antwortfläche (vgl. Abbildung~\ref{fig:ma_abb6.1_response_surface}).

\begin{figure}[htbp]
    \centering
    \setlength\figurewidth{0.9\textwidth}
    \setlength\figureheight{6.5cm}
    % This file was created by matlab2tikz.
%
\definecolor{mycolor1}{rgb}{0.18400,0.74500,0.93700}%
\definecolor{mycolor2}{rgb}{0.81900,0.01500,0.54500}%
\definecolor{mycolor3}{rgb}{0.92900,0.69400,0.12500}%
\definecolor{mycolor4}{rgb}{0.52100,0.08600,0.81900}%
\definecolor{mycolor5}{rgb}{0.12941,0.12941,0.12941}%
%
\begin{tikzpicture}

\begin{axis}[%
width=\figurewidth,
height=\figureheight,
at={(0\figurewidth,0\figureheight)},
scale only axis,
unbounded coords=jump,
xmin=-4.00,
xmax=2.00,
tick align=outside,
xlabel style={font=\color{mycolor5}},
xlabel={$x_1$},
ymin=-2.80,
ymax=2.00,
ylabel style={font=\color{mycolor5}},
ylabel={$x_2$},
zmin=0.00,
zmax=4050000.00,
zlabel style={font=\color{mycolor5}},
zlabel={$\hat{T}$ [h]},
view={3.00}{50.00},
axis background/.style={fill=white},
xmajorgrids,
ymajorgrids,
zmajorgrids,
grid style={dotted, white!30!black, opacity=0.4},
legend style={at={(0.97,0.5)}, anchor=east, legend cell align=left, align=left, draw=white!80!black}
]

\addplot3[%
surf,
fill opacity=0.55, draw opacity=0.80, fill=white!95!black, faceted color=white!40!black, z buffer=sort, colormap={mymap}{[1pt] rgb(0pt)=(0.2422,0.1504,0.6603); rgb(1pt)=(0.2444,0.1534,0.6728); rgb(2pt)=(0.2464,0.1569,0.6847); rgb(3pt)=(0.2484,0.1607,0.6961); rgb(4pt)=(0.2503,0.1648,0.7071); rgb(5pt)=(0.2522,0.1689,0.7179); rgb(6pt)=(0.254,0.1732,0.7286); rgb(7pt)=(0.2558,0.1773,0.7393); rgb(8pt)=(0.2576,0.1814,0.7501); rgb(9pt)=(0.2594,0.1854,0.761); rgb(11pt)=(0.2628,0.1932,0.7828); rgb(12pt)=(0.2645,0.1972,0.7937); rgb(13pt)=(0.2661,0.2011,0.8043); rgb(14pt)=(0.2676,0.2052,0.8148); rgb(15pt)=(0.2691,0.2094,0.8249); rgb(16pt)=(0.2704,0.2138,0.8346); rgb(17pt)=(0.2717,0.2184,0.8439); rgb(18pt)=(0.2729,0.2231,0.8528); rgb(19pt)=(0.274,0.228,0.8612); rgb(20pt)=(0.2749,0.233,0.8692); rgb(21pt)=(0.2758,0.2382,0.8767); rgb(22pt)=(0.2766,0.2435,0.884); rgb(23pt)=(0.2774,0.2489,0.8908); rgb(24pt)=(0.2781,0.2543,0.8973); rgb(25pt)=(0.2788,0.2598,0.9035); rgb(26pt)=(0.2794,0.2653,0.9094); rgb(27pt)=(0.2798,0.2708,0.915); rgb(28pt)=(0.2802,0.2764,0.9204); rgb(29pt)=(0.2806,0.2819,0.9255); rgb(30pt)=(0.2809,0.2875,0.9305); rgb(31pt)=(0.2811,0.293,0.9352); rgb(32pt)=(0.2813,0.2985,0.9397); rgb(33pt)=(0.2814,0.304,0.9441); rgb(34pt)=(0.2814,0.3095,0.9483); rgb(35pt)=(0.2813,0.315,0.9524); rgb(36pt)=(0.2811,0.3204,0.9563); rgb(37pt)=(0.2809,0.3259,0.96); rgb(38pt)=(0.2807,0.3313,0.9636); rgb(39pt)=(0.2803,0.3367,0.967); rgb(40pt)=(0.2798,0.3421,0.9702); rgb(41pt)=(0.2791,0.3475,0.9733); rgb(42pt)=(0.2784,0.3529,0.9763); rgb(43pt)=(0.2776,0.3583,0.9791); rgb(44pt)=(0.2766,0.3638,0.9817); rgb(45pt)=(0.2754,0.3693,0.984); rgb(46pt)=(0.2741,0.3748,0.9862); rgb(47pt)=(0.2726,0.3804,0.9881); rgb(48pt)=(0.271,0.386,0.9898); rgb(49pt)=(0.2691,0.3916,0.9912); rgb(50pt)=(0.267,0.3973,0.9924); rgb(51pt)=(0.2647,0.403,0.9935); rgb(52pt)=(0.2621,0.4088,0.9946); rgb(53pt)=(0.2591,0.4145,0.9955); rgb(54pt)=(0.2556,0.4203,0.9965); rgb(55pt)=(0.2517,0.4261,0.9974); rgb(56pt)=(0.2473,0.4319,0.9983); rgb(57pt)=(0.2424,0.4378,0.9991); rgb(58pt)=(0.2369,0.4437,0.9996); rgb(59pt)=(0.2311,0.4497,0.9995); rgb(60pt)=(0.225,0.4559,0.9985); rgb(61pt)=(0.2189,0.462,0.9968); rgb(62pt)=(0.2128,0.4682,0.9948); rgb(63pt)=(0.2066,0.4743,0.9926); rgb(64pt)=(0.2006,0.4803,0.9906); rgb(65pt)=(0.195,0.4861,0.9887); rgb(66pt)=(0.1903,0.4919,0.9867); rgb(67pt)=(0.1869,0.4975,0.9844); rgb(68pt)=(0.1847,0.503,0.9819); rgb(69pt)=(0.1831,0.5084,0.9793); rgb(70pt)=(0.1818,0.5138,0.9766); rgb(71pt)=(0.1806,0.5191,0.9738); rgb(72pt)=(0.1795,0.5244,0.9709); rgb(73pt)=(0.1785,0.5296,0.9677); rgb(74pt)=(0.1778,0.5349,0.9641); rgb(75pt)=(0.1773,0.5401,0.9602); rgb(76pt)=(0.1768,0.5452,0.956); rgb(77pt)=(0.1764,0.5504,0.9516); rgb(78pt)=(0.1755,0.5554,0.9473); rgb(79pt)=(0.174,0.5605,0.9432); rgb(80pt)=(0.1716,0.5655,0.9393); rgb(81pt)=(0.1686,0.5705,0.9357); rgb(82pt)=(0.1649,0.5755,0.9323); rgb(83pt)=(0.161,0.5805,0.9289); rgb(84pt)=(0.1573,0.5854,0.9254); rgb(85pt)=(0.154,0.5902,0.9218); rgb(86pt)=(0.1513,0.595,0.9182); rgb(87pt)=(0.1492,0.5997,0.9147); rgb(88pt)=(0.1475,0.6043,0.9113); rgb(89pt)=(0.1461,0.6089,0.908); rgb(90pt)=(0.1446,0.6135,0.905); rgb(91pt)=(0.1429,0.618,0.9022); rgb(92pt)=(0.1408,0.6226,0.8998); rgb(93pt)=(0.1383,0.6272,0.8975); rgb(94pt)=(0.1354,0.6317,0.8953); rgb(95pt)=(0.1321,0.6363,0.8932); rgb(96pt)=(0.1288,0.6408,0.891); rgb(97pt)=(0.1253,0.6453,0.8887); rgb(98pt)=(0.1219,0.6497,0.8862); rgb(99pt)=(0.1185,0.6541,0.8834); rgb(100pt)=(0.1152,0.6584,0.8804); rgb(101pt)=(0.1119,0.6627,0.877); rgb(102pt)=(0.1085,0.6669,0.8734); rgb(103pt)=(0.1048,0.671,0.8695); rgb(104pt)=(0.1009,0.675,0.8653); rgb(105pt)=(0.0964,0.6789,0.8609); rgb(106pt)=(0.0914,0.6828,0.8562); rgb(107pt)=(0.0855,0.6865,0.8513); rgb(108pt)=(0.0789,0.6902,0.8462); rgb(109pt)=(0.0713,0.6938,0.8409); rgb(110pt)=(0.0628,0.6972,0.8355); rgb(111pt)=(0.0535,0.7006,0.8299); rgb(112pt)=(0.0433,0.7039,0.8242); rgb(113pt)=(0.0328,0.7071,0.8183); rgb(114pt)=(0.0234,0.7103,0.8124); rgb(115pt)=(0.0155,0.7133,0.8064); rgb(116pt)=(0.0091,0.7163,0.8003); rgb(117pt)=(0.0046,0.7192,0.7941); rgb(118pt)=(0.0019,0.722,0.7878); rgb(119pt)=(0.0009,0.7248,0.7815); rgb(120pt)=(0.0018,0.7275,0.7752); rgb(121pt)=(0.0046,0.7301,0.7688); rgb(122pt)=(0.0094,0.7327,0.7623); rgb(123pt)=(0.0162,0.7352,0.7558); rgb(124pt)=(0.0253,0.7376,0.7492); rgb(125pt)=(0.0369,0.74,0.7426); rgb(126pt)=(0.0504,0.7423,0.7359); rgb(127pt)=(0.0638,0.7446,0.7292); rgb(128pt)=(0.077,0.7468,0.7224); rgb(129pt)=(0.0899,0.7489,0.7156); rgb(130pt)=(0.1023,0.751,0.7088); rgb(131pt)=(0.1141,0.7531,0.7019); rgb(132pt)=(0.1252,0.7552,0.695); rgb(133pt)=(0.1354,0.7572,0.6881); rgb(134pt)=(0.1448,0.7593,0.6812); rgb(135pt)=(0.1532,0.7614,0.6741); rgb(136pt)=(0.1609,0.7635,0.6671); rgb(137pt)=(0.1678,0.7656,0.6599); rgb(138pt)=(0.1741,0.7678,0.6527); rgb(139pt)=(0.1799,0.7699,0.6454); rgb(140pt)=(0.1853,0.7721,0.6379); rgb(141pt)=(0.1905,0.7743,0.6303); rgb(142pt)=(0.1954,0.7765,0.6225); rgb(143pt)=(0.2003,0.7787,0.6146); rgb(144pt)=(0.2061,0.7808,0.6065); rgb(145pt)=(0.2118,0.7828,0.5983); rgb(146pt)=(0.2178,0.7849,0.5899); rgb(147pt)=(0.2244,0.7869,0.5813); rgb(148pt)=(0.2318,0.7887,0.5725); rgb(149pt)=(0.2401,0.7905,0.5636); rgb(150pt)=(0.2491,0.7922,0.5546); rgb(151pt)=(0.2589,0.7937,0.5454); rgb(152pt)=(0.2695,0.7951,0.536); rgb(153pt)=(0.2809,0.7964,0.5266); rgb(154pt)=(0.2929,0.7975,0.517); rgb(155pt)=(0.3052,0.7985,0.5074); rgb(156pt)=(0.3176,0.7994,0.4975); rgb(157pt)=(0.3301,0.8002,0.4876); rgb(158pt)=(0.3424,0.8009,0.4774); rgb(159pt)=(0.3548,0.8016,0.4669); rgb(160pt)=(0.3671,0.8021,0.4563); rgb(161pt)=(0.3795,0.8026,0.4454); rgb(162pt)=(0.3921,0.8029,0.4344); rgb(163pt)=(0.405,0.8031,0.4233); rgb(164pt)=(0.4184,0.803,0.4122); rgb(165pt)=(0.4322,0.8028,0.4013); rgb(166pt)=(0.4463,0.8024,0.3904); rgb(167pt)=(0.4608,0.8018,0.3797); rgb(168pt)=(0.4753,0.8011,0.3691); rgb(169pt)=(0.4899,0.8002,0.3586); rgb(170pt)=(0.5044,0.7993,0.348); rgb(171pt)=(0.5187,0.7982,0.3374); rgb(172pt)=(0.5329,0.797,0.3267); rgb(173pt)=(0.547,0.7957,0.3159); rgb(175pt)=(0.5748,0.7929,0.2941); rgb(176pt)=(0.5886,0.7913,0.2833); rgb(177pt)=(0.6024,0.7896,0.2726); rgb(178pt)=(0.6161,0.7878,0.2622); rgb(179pt)=(0.6297,0.7859,0.2521); rgb(180pt)=(0.6433,0.7839,0.2423); rgb(181pt)=(0.6567,0.7818,0.2329); rgb(182pt)=(0.6701,0.7796,0.2239); rgb(183pt)=(0.6833,0.7773,0.2155); rgb(184pt)=(0.6963,0.775,0.2075); rgb(185pt)=(0.7091,0.7727,0.1998); rgb(186pt)=(0.7218,0.7703,0.1924); rgb(187pt)=(0.7344,0.7679,0.1852); rgb(188pt)=(0.7468,0.7654,0.1782); rgb(189pt)=(0.759,0.7629,0.1717); rgb(190pt)=(0.771,0.7604,0.1658); rgb(191pt)=(0.7829,0.7579,0.1608); rgb(192pt)=(0.7945,0.7554,0.157); rgb(193pt)=(0.806,0.7529,0.1546); rgb(194pt)=(0.8172,0.7505,0.1535); rgb(195pt)=(0.8281,0.7481,0.1536); rgb(196pt)=(0.8389,0.7457,0.1546); rgb(197pt)=(0.8495,0.7435,0.1564); rgb(198pt)=(0.86,0.7413,0.1587); rgb(199pt)=(0.8703,0.7392,0.1615); rgb(200pt)=(0.8804,0.7372,0.165); rgb(201pt)=(0.8903,0.7353,0.1695); rgb(202pt)=(0.9,0.7336,0.1749); rgb(203pt)=(0.9093,0.7321,0.1815); rgb(204pt)=(0.9184,0.7308,0.189); rgb(205pt)=(0.9272,0.7298,0.1973); rgb(206pt)=(0.9357,0.729,0.2061); rgb(207pt)=(0.944,0.7285,0.2151); rgb(208pt)=(0.9523,0.7284,0.2237); rgb(209pt)=(0.9606,0.7285,0.2312); rgb(210pt)=(0.9689,0.7292,0.2373); rgb(211pt)=(0.977,0.7304,0.2418); rgb(212pt)=(0.9842,0.733,0.2446); rgb(213pt)=(0.99,0.7365,0.2429); rgb(214pt)=(0.9946,0.7407,0.2394); rgb(215pt)=(0.9966,0.7458,0.2351); rgb(216pt)=(0.9971,0.7513,0.2309); rgb(217pt)=(0.9972,0.7569,0.2267); rgb(218pt)=(0.9971,0.7626,0.2224); rgb(219pt)=(0.9969,0.7683,0.2181); rgb(220pt)=(0.9966,0.774,0.2138); rgb(221pt)=(0.9962,0.7798,0.2095); rgb(222pt)=(0.9957,0.7856,0.2053); rgb(223pt)=(0.9949,0.7915,0.2012); rgb(224pt)=(0.9938,0.7974,0.1974); rgb(225pt)=(0.9923,0.8034,0.1939); rgb(226pt)=(0.9906,0.8095,0.1906); rgb(227pt)=(0.9885,0.8156,0.1875); rgb(228pt)=(0.9861,0.8218,0.1846); rgb(229pt)=(0.9835,0.828,0.1817); rgb(230pt)=(0.9807,0.8342,0.1787); rgb(231pt)=(0.9778,0.8404,0.1757); rgb(232pt)=(0.9748,0.8467,0.1726); rgb(233pt)=(0.972,0.8529,0.1695); rgb(234pt)=(0.9694,0.8591,0.1665); rgb(235pt)=(0.9671,0.8654,0.1636); rgb(236pt)=(0.9651,0.8716,0.1608); rgb(237pt)=(0.9634,0.8778,0.1582); rgb(238pt)=(0.9619,0.884,0.1557); rgb(239pt)=(0.9608,0.8902,0.1532); rgb(240pt)=(0.9601,0.8963,0.1507); rgb(241pt)=(0.9596,0.9023,0.148); rgb(242pt)=(0.9595,0.9084,0.145); rgb(243pt)=(0.9597,0.9143,0.1418); rgb(244pt)=(0.9601,0.9203,0.1382); rgb(245pt)=(0.9608,0.9262,0.1344); rgb(246pt)=(0.9618,0.932,0.1304); rgb(247pt)=(0.9629,0.9379,0.1261); rgb(248pt)=(0.9642,0.9437,0.1216); rgb(249pt)=(0.9657,0.9494,0.1168); rgb(250pt)=(0.9674,0.9552,0.1116); rgb(251pt)=(0.9692,0.9609,0.1061); rgb(252pt)=(0.9711,0.9667,0.1001); rgb(253pt)=(0.973,0.9724,0.0938); rgb(254pt)=(0.9749,0.9782,0.0872); rgb(255pt)=(0.9769,0.9839,0.0805)}, mesh/rows=37]
table[row sep=crcr, point meta=\thisrow{c}] {%
%
x	y	z	c\\
-3.40	-2.80	26514744709.54	26514744709.54\\
-3.40	-2.65	16111762065.41	16111762065.41\\
-3.40	-2.50	9987042316.04	9987042316.04\\
-3.40	-2.35	6314936690.53	6314936690.53\\
-3.40	-2.20	4073234114.89	4073234114.89\\
-3.40	-2.05	2680081522.93	2680081522.93\\
-3.40	-1.90	1798849779.93	1798849779.93\\
-3.40	-1.75	1231629409.68	1231629409.68\\
-3.40	-1.60	860207975.20	860207975.20\\
-3.40	-1.45	612865450.50	612865450.50\\
-3.40	-1.30	445415242.01	445415242.01\\
-3.40	-1.15	330219921.78	330219921.78\\
-3.40	-1.00	249735180.68	249735180.68\\
-3.40	-0.85	192661280.99	192661280.99\\
-3.40	-0.70	151616833.12	151616833.12\\
-3.40	-0.55	121713474.25	121713474.25\\
-3.40	-0.40	99670849.12	99670849.12\\
-3.40	-0.25	83259906.25	83259906.25\\
-3.40	-0.10	70948290.82	70948290.82\\
-3.40	0.05	61671742.33	61671742.33\\
-3.40	0.20	54685068.36	54685068.36\\
-3.40	0.35	49464037.32	49464037.32\\
-3.40	0.50	45640314.00	45640314.00\\
-3.40	0.65	42958187.24	42958187.24\\
-3.40	0.80	41245970.83	41245970.83\\
-3.40	0.95	40397582.31	40397582.31\\
-3.40	1.10	40361516.85	40361516.85\\
-3.40	1.25	41135600.81	41135600.81\\
-3.40	1.40	42766771.96	42766771.96\\
-3.40	1.55	45355854.93	45355854.93\\
-3.40	1.70	49068016.75	49068016.75\\
-3.40	1.85	54150430.38	54150430.38\\
-3.40	2.00	60959806.70	60959806.70\\
-3.25	-2.80	13458783092.05	13458783092.05\\
-3.25	-2.65	8288645513.91	8288645513.91\\
-3.25	-2.50	5207144135.38	5207144135.38\\
-3.25	-2.35	3336982221.30	3336982221.30\\
-3.25	-2.20	2181455946.47	2181455946.47\\
-3.25	-2.05	1454712805.43	1454712805.43\\
-3.25	-1.90	989569582.03	989569582.03\\
-3.25	-1.75	686678853.30	686678853.30\\
-3.25	-1.60	486070506.32	486070506.32\\
-3.25	-1.45	350980612.12	350980612.12\\
-3.25	-1.30	258526606.11	258526606.11\\
-3.25	-1.15	194252055.62	194252055.62\\
-3.25	-1.00	148889565.44	148889565.44\\
-3.25	-0.85	116412914.41	116412914.41\\
-3.25	-0.70	92848805.78	92848805.78\\
-3.25	-0.55	75542219.92	75542219.92\\
-3.25	-0.40	62696229.65	62696229.65\\
-3.25	-0.25	53080055.29	53080055.29\\
-3.25	-0.10	45841579.09	45841579.09\\
-3.25	0.05	40385552.86	40385552.86\\
-3.25	0.20	36293659.16	36293659.16\\
-3.25	0.35	33271604.19	33271604.19\\
-3.25	0.50	31113938.29	31113938.29\\
-3.25	0.65	29680724.08	29680724.08\\
-3.25	0.80	28882332.14	28882332.14\\
-3.25	0.95	28670039.07	28670039.07\\
-3.25	1.10	29031038.54	29031038.54\\
-3.25	1.25	29987145.05	29987145.05\\
-3.25	1.40	31597005.75	31597005.75\\
-3.25	1.55	33962136.13	33962136.13\\
-3.25	1.70	37237655.38	37237655.38\\
-3.25	1.85	41649320.36	41649320.36\\
-3.25	2.00	47519489.90	47519489.90\\
-3.10	-2.80	6901995273.65	6901995273.65\\
-3.10	-2.65	4307989266.73	4307989266.73\\
-3.10	-2.50	2742918023.60	2742918023.60\\
-3.10	-2.35	1781514266.12	1781514266.12\\
-3.10	-2.20	1180331619.15	1180331619.15\\
-3.10	-2.05	797732046.29	797732046.29\\
-3.10	-1.90	549981773.11	549981773.11\\
-3.10	-1.75	386792296.11	386792296.11\\
-3.10	-1.60	277488886.85	277488886.85\\
-3.10	-1.45	203072742.79	203072742.79\\
-3.10	-1.30	151598859.87	151598859.87\\
-3.10	-1.15	115445895.32	115445895.32\\
-3.10	-1.00	89680767.82	89680767.82\\
-3.10	-0.85	71065431.96	71065431.96\\
-3.10	-0.70	57445464.18	57445464.18\\
-3.10	-0.55	47368686.98	47368686.98\\
-3.10	-0.40	39844209.11	39844209.11\\
-3.10	-0.25	34188286.55	34188286.55\\
-3.10	-0.10	29924556.37	29924556.37\\
-3.10	0.05	26718764.16	26718764.16\\
-3.10	0.20	24335668.07	24335668.07\\
-3.10	0.35	22610410.14	22610410.14\\
-3.10	0.50	21429491.61	21429491.61\\
-3.10	0.65	20718273.24	20718273.24\\
-3.10	0.80	20433064.47	20433064.47\\
-3.10	0.95	20556620.34	20556620.34\\
-3.10	1.10	21096391.92	21096391.92\\
-3.10	1.25	22085280.31	22085280.31\\
-3.10	1.40	23585001.44	23585001.44\\
-3.10	1.55	25692546.97	25692546.97\\
-3.10	1.70	28550694.16	28550694.16\\
-3.10	1.85	32364167.85	32364167.85\\
-3.10	2.00	37424023.63	37424023.63\\
-2.95	-2.80	3575971435.05	3575971435.05\\
-2.95	-2.65	2262122808.53	2262122808.53\\
-2.95	-2.50	1459743590.34	1459743590.34\\
-2.95	-2.35	960893456.34	960893456.34\\
-2.95	-2.20	645226439.74	645226439.74\\
-2.95	-2.05	441964447.73	441964447.73\\
-2.95	-1.90	308816698.84	308816698.84\\
-2.95	-1.75	220116438.61	220116438.61\\
-2.95	-1.60	160045121.10	160045121.10\\
-2.95	-1.45	118705450.88	118705450.88\\
-2.95	-1.30	89812574.37	89812574.37\\
-2.95	-1.15	69317340.03	69317340.03\\
-2.95	-1.00	54573887.74	54573887.74\\
-2.95	-0.85	43829463.85	43829463.85\\
-2.95	-0.70	35907542.54	35907542.54\\
-2.95	-0.55	30008443.51	30008443.51\\
-2.95	-0.40	25582295.82	25582295.82\\
-2.95	-0.25	22247121.36	22247121.36\\
-2.95	-0.10	19735420.85	19735420.85\\
-2.95	0.05	17859003.42	17859003.42\\
-2.95	0.20	16485658.90	16485658.90\\
-2.95	0.35	15523643.46	15523643.46\\
-2.95	0.50	14911429.17	14911429.17\\
-2.95	0.65	14611107.66	14611107.66\\
-2.95	0.80	14604452.23	14604452.23\\
-2.95	0.95	14891061.76	14891061.76\\
-2.95	1.10	15488320.17	15488320.17\\
-2.95	1.25	16433165.57	16433165.57\\
-2.95	1.40	17785922.88	17785922.88\\
-2.95	1.55	19636760.33	19636760.33\\
-2.95	1.70	22115743.02	22115743.02\\
-2.95	1.85	25408059.09	25408059.09\\
-2.95	2.00	29776915.41	29776915.41\\
-2.80	-2.80	1871819400.26	1871819400.26\\
-2.80	-2.65	1200074657.07	1200074657.07\\
-2.80	-2.50	784857559.20	784857559.20\\
-2.80	-2.35	523614527.03	523614527.03\\
-2.80	-2.20	356345098.24	356345098.24\\
-2.80	-2.05	247382031.67	247382031.67\\
-2.80	-1.90	175187766.31	175187766.31\\
-2.80	-1.75	126554519.51	126554519.51\\
-2.80	-1.60	93258795.91	93258795.91\\
-2.80	-1.45	70103581.98	70103581.98\\
-2.80	-1.30	53756238.84	53756238.84\\
-2.80	-1.15	42049012.51	42049012.51\\
-2.80	-1.00	33552198.94	33552198.94\\
-2.80	-0.85	27310172.03	27310172.03\\
-2.80	-0.70	22675984.11	22675984.11\\
-2.80	-0.55	19206406.34	19206406.34\\
-2.80	-0.40	16594506.53	16594506.53\\
-2.80	-0.25	14625840.48	14625840.48\\
-2.80	-0.10	13149692.07	13149692.07\\
-2.80	0.05	12060035.72	12060035.72\\
-2.80	0.20	11282877.32	11282877.32\\
-2.80	0.35	10767859.95	10767859.95\\
-2.80	0.50	10482797.38	10482797.38\\
-2.80	0.65	10410300.00	10410300.00\\
-2.80	0.80	10545994.96	10545994.96\\
-2.80	0.95	10898083.58	10898083.58\\
-2.80	1.10	11488173.07	11488173.07\\
-2.80	1.25	12353501.32	12353501.32\\
-2.80	1.40	13550877.60	13550877.60\\
-2.80	1.55	15162926.85	15162926.85\\
-2.80	1.70	17307602.72	17307602.72\\
-2.80	1.85	20152505.86	20152505.86\\
-2.80	2.00	23936434.26	23936434.26\\
-2.65	-2.80	989884142.47	989884142.47\\
-2.65	-2.65	643207159.83	643207159.83\\
-2.65	-2.50	426339561.02	426339561.02\\
-2.65	-2.35	288269475.49	288269475.49\\
-2.65	-2.20	198829098.19	198829098.19\\
-2.65	-2.05	139894128.84	139894128.84\\
-2.65	-1.90	100405451.67	100405451.67\\
-2.65	-1.75	73511174.36	73511174.36\\
-2.65	-1.60	54901939.65	54901939.65\\
-2.65	-1.45	41827342.82	41827342.82\\
-2.65	-1.30	32506568.35	32506568.35\\
-2.65	-1.15	25770345.04	25770345.04\\
-2.65	-1.00	20840474.51	20840474.51\\
-2.65	-0.85	17192270.41	17192270.41\\
-2.65	-0.70	14467621.03	14467621.03\\
-2.65	-0.55	12419361.70	12419361.70\\
-2.65	-0.40	10875260.98	10875260.98\\
-2.65	-0.25	9714453.20	9714453.20\\
-2.65	-0.10	8851875.43	8851875.43\\
-2.65	0.05	8227928.09	8227928.09\\
-2.65	0.20	7801604.71	7801604.71\\
-2.65	0.35	7545980.13	7545980.13\\
-2.65	0.50	7445358.82	7445358.82\\
-2.65	0.65	7493658.01	7493658.01\\
-2.65	0.80	7693790.67	7693790.67\\
-2.65	0.95	8057960.30	8057960.30\\
-2.65	1.10	8608909.49	8608909.49\\
-2.65	1.25	9382302.27	9382302.27\\
-2.65	1.40	10430592.04	10430592.04\\
-2.65	1.55	11828965.35	11828965.35\\
-2.65	1.70	13684307.36	13684307.36\\
-2.65	1.85	16148684.19	16148684.19\\
-2.65	2.00	19439707.37	19439707.37\\
-2.50	-2.80	528877766.13	528877766.13\\
-2.50	-2.65	348292398.00	348292398.00\\
-2.50	-2.50	233975800.44	233975800.44\\
-2.50	-2.35	160337884.26	160337884.26\\
-2.50	-2.20	112082966.28	112082966.28\\
-2.50	-2.05	79924760.60	79924760.60\\
-2.50	-1.90	58138166.70	58138166.70\\
-2.50	-1.75	43139944.50	43139944.50\\
-2.50	-1.60	32653977.71	32653977.71\\
-2.50	-1.45	25213367.99	25213367.99\\
-2.50	-1.30	19859298.04	19859298.04\\
-2.50	-1.15	15956409.94	15956409.94\\
-2.50	-1.00	13078102.52	13078102.52\\
-2.50	-0.85	10934339.41	10934339.41\\
-2.50	-0.70	9325639.52	9325639.52\\
-2.50	-0.55	8113401.32	8113401.32\\
-2.50	-0.40	7200548.05	7200548.05\\
-2.50	-0.25	6518781.26	6518781.26\\
-2.50	-0.10	6020125.21	6020125.21\\
-2.50	0.05	5671303.65	5671303.65\\
-2.50	0.20	5450025.54	5450025.54\\
-2.50	0.35	5342597.24	5342597.24\\
-2.50	0.50	5342500.78	5342500.78\\
-2.50	0.65	5449730.34	5449730.34\\
-2.50	0.80	5670791.67	5670791.67\\
-2.50	0.95	6019364.36	6019364.36\\
-2.50	1.10	6517722.01	6517722.01\\
-2.50	1.25	7199118.05	7199118.05\\
-2.50	1.40	8111497.09	8111497.09\\
-2.50	1.55	9323114.10	9323114.10\\
-2.50	1.70	10930983.60	10930983.60\\
-2.50	1.85	13073616.66	13073616.66\\
-2.50	2.00	15950360.78	15950360.78\\
-2.35	-2.80	285480712.65	285480712.65\\
-2.35	-2.65	190540649.18	190540649.18\\
-2.35	-2.50	129728910.11	129728910.11\\
-2.35	-2.35	90099868.88	90099868.88\\
-2.35	-2.20	63833669.26	63833669.26\\
-2.35	-2.05	46133214.83	46133214.83\\
-2.35	-1.90	34010725.11	34010725.11\\
-2.35	-1.75	25577395.74	25577395.74\\
-2.35	-1.60	19621627.00	19621627.00\\
-2.35	-1.45	15355075.95	15355075.95\\
-2.35	-1.30	12257649.04	12257649.04\\
-2.35	-1.15	9981611.37	9981611.37\\
-2.35	-1.00	8291486.38	8291486.38\\
-2.35	-0.85	7025906.85	7025906.85\\
-2.35	-0.70	6073102.96	6073102.96\\
-2.35	-0.55	5354971.46	5354971.46\\
-2.35	-0.40	4816615.04	4816615.04\\
-2.35	-0.25	4419416.94	4419416.94\\
-2.35	-0.10	4136435.67	4136435.67\\
-2.35	0.05	3949351.91	3949351.91\\
-2.35	0.20	3846481.55	3846481.55\\
-2.35	0.35	3821551.65	3821551.65\\
-2.35	0.50	3873058.66	3873058.66\\
-2.35	0.65	4004116.24	4004116.24\\
-2.35	0.80	4222771.08	4222771.08\\
-2.35	0.95	4542831.85	4542831.85\\
-2.35	1.10	4985331.53	4985331.53\\
-2.35	1.25	5580841.51	5580841.51\\
-2.35	1.40	6372995.23	6372995.23\\
-2.35	1.55	7423791.38	7423791.38\\
-2.35	1.70	8821576.10	8821576.10\\
-2.35	1.85	10693130.60	10693130.60\\
-2.35	2.00	13222142.04	13222142.04\\
-2.20	-2.80	155685711.01	155685711.01\\
-2.20	-2.65	105312956.57	105312956.57\\
-2.20	-2.50	72669658.11	72669658.11\\
-2.20	-2.35	51152008.27	51152008.27\\
-2.20	-2.20	36729119.90	36729119.90\\
-2.20	-2.05	26902746.25	26902746.25\\
-2.20	-1.90	20101153.00	20101153.00\\
-2.20	-1.75	15320875.52	15320875.52\\
-2.20	-1.60	11911993.81	11911993.81\\
-2.20	-1.45	9447645.63	9447645.63\\
-2.20	-1.30	7643653.50	7643653.50\\
-2.20	-1.15	6308362.60	6308362.60\\
-2.20	-1.00	5310929.86	5310929.86\\
-2.20	-0.85	4561028.06	4561028.06\\
-2.20	-0.70	3995702.85	3995702.85\\
-2.20	-0.55	3570770.04	3570770.04\\
-2.20	-0.40	3255133.78	3255133.78\\
-2.20	-0.25	3027011.45	3027011.45\\
-2.20	-0.10	2871425.44	2871425.44\\
-2.20	0.05	2778556.85	2778556.85\\
-2.20	0.20	2742706.20	2742706.20\\
-2.20	0.35	2761706.69	2761706.69\\
-2.20	0.50	2836704.36	2836704.36\\
-2.20	0.65	2972274.12	2972274.12\\
-2.20	0.80	3176888.00	3176888.00\\
-2.20	0.95	3463803.21	3463803.21\\
-2.20	1.10	3852501.14	3852501.14\\
-2.20	1.25	4370897.31	4370897.31\\
-2.20	1.40	5058673.97	5058673.97\\
-2.20	1.55	5972292.07	5972292.07\\
-2.20	1.70	7192562.54	7192562.54\\
-2.20	1.85	8836179.22	8836179.22\\
-2.20	2.00	11073467.76	11073467.76\\
-2.05	-2.80	85777079.68	85777079.68\\
-2.05	-2.65	58806656.78	58806656.78\\
-2.05	-2.50	41126332.67	41126332.67\\
-2.05	-2.35	29339434.52	29339434.52\\
-2.05	-2.20	21351173.74	21351173.74\\
-2.05	-2.05	15850027.88	15850027.88\\
-2.05	-1.90	12002634.11	12002634.11\\
-2.05	-1.75	9271742.56	9271742.56\\
-2.05	-1.60	7306079.98	7306079.98\\
-2.05	-1.45	5872807.16	5872807.16\\
-2.05	-1.30	4815543.61	4815543.61\\
-2.05	-1.15	4027941.55	4027941.55\\
-2.05	-1.00	3436839.57	3436839.57\\
-2.05	-0.85	2991394.02	2991394.02\\
-2.05	-0.70	2655988.82	2655988.82\\
-2.05	-0.55	2405565.15	2405565.15\\
-2.05	-0.40	2222523.02	2222523.02\\
-2.05	-0.25	2094660.63	2094660.63\\
-2.05	-0.10	2013813.90	2013813.90\\
-2.05	0.05	1974982.53	1974982.53\\
-2.05	0.20	1975811.20	1975811.20\\
-2.05	0.35	2016349.86	2016349.86\\
-2.05	0.50	2099058.76	2099058.76\\
-2.05	0.65	2229058.99	2229058.99\\
-2.05	0.80	2414664.46	2414664.46\\
-2.05	0.95	2668273.10	2668273.10\\
-2.05	1.10	3007752.02	3007752.02\\
-2.05	1.25	3458533.90	3458533.90\\
-2.05	1.40	4056769.26	4056769.26\\
-2.05	1.55	4854078.97	4854078.97\\
-2.05	1.70	5924771.81	5924771.81\\
-2.05	1.85	7376913.31	7376913.31\\
-2.05	2.00	9369490.87	9369490.87\\
-1.90	-2.80	47746800.85	47746800.85\\
-1.90	-2.65	33175822.16	33175822.16\\
-1.90	-2.50	23514587.86	23514587.86\\
-1.90	-2.35	17001659.29	17001659.29\\
-1.90	-2.20	12539594.91	12539594.91\\
-1.90	-2.05	9434393.05	9434393.05\\
-1.90	-1.90	7240735.53	7240735.53\\
-1.90	-1.75	5668780.81	5668780.81\\
-1.90	-1.60	4527254.46	4527254.46\\
-1.90	-1.45	3688233.48	3688233.48\\
-1.90	-1.30	3065068.38	3065068.38\\
-1.90	-1.15	2598365.20	2598365.20\\
-1.90	-1.00	2246976.21	2246976.21\\
-1.90	-0.85	1982143.15	1982143.15\\
-1.90	-0.70	1783650.75	1783650.75\\
-1.90	-0.55	1637279.73	1637279.73\\
-1.90	-0.40	1533113.16	1533113.16\\
-1.90	-0.25	1464413.75	1464413.75\\
-1.90	-0.10	1426893.78	1426893.78\\
-1.90	0.05	1418266.19	1418266.19\\
-1.90	0.20	1438010.70	1438010.70\\
-1.90	0.35	1487321.13	1487321.13\\
-1.90	0.50	1569226.51	1569226.51\\
-1.90	0.65	1688903.31	1688903.31\\
-1.90	0.80	1854223.99	1854223.99\\
-1.90	0.95	2076623.99	2076623.99\\
-1.90	1.10	2372421.21	2372421.21\\
-1.90	1.25	2764801.71	2764801.71\\
-1.90	1.40	3286808.81	3286808.81\\
-1.90	1.55	3985870.24	3985870.24\\
-1.90	1.70	4930717.60	4930717.60\\
-1.90	1.85	6222076.75	6222076.75\\
-1.90	2.00	8009379.26	8009379.26\\
-1.75	-2.80	26851454.74	26851454.74\\
-1.75	-2.65	18908950.22	18908950.22\\
-1.75	-2.50	13583299.54	13583299.54\\
-1.75	-2.35	9953627.83	9953627.83\\
-1.75	-2.20	7440391.35	7440391.35\\
-1.75	-2.05	5673465.53	5673465.53\\
-1.75	-1.90	4413054.82	4413054.82\\
-1.75	-1.75	3501615.73	3501615.73\\
-1.75	-1.60	2834235.33	2834235.33\\
-1.75	-1.45	2340138.66	2340138.66\\
-1.75	-1.30	1970995.06	1970995.06\\
-1.75	-1.15	1693431.92	1693431.92\\
-1.75	-1.00	1484185.58	1484185.58\\
-1.75	-0.85	1326926.70	1326926.70\\
-1.75	-0.70	1210163.14	1210163.14\\
-1.75	-0.55	1125846.46	1125846.46\\
-1.75	-0.40	1068446.24	1068446.24\\
-1.75	-0.25	1034342.68	1034342.68\\
-1.75	-0.10	1021443.79	1021443.79\\
-1.75	0.05	1028970.12	1028970.12\\
-1.75	0.20	1057375.67	1057375.67\\
-1.75	0.35	1108393.89	1108393.89\\
-1.75	0.50	1185215.15	1185215.15\\
-1.75	0.65	1292821.39	1292821.39\\
-1.75	0.80	1438527.34	1438527.34\\
-1.75	0.95	1632811.21	1632811.21\\
-1.75	1.10	1890567.07	1890567.07\\
-1.75	1.25	2232988.38	2232988.38\\
-1.75	1.40	2690413.90	2690413.90\\
-1.75	1.55	3306663.45	3306663.45\\
-1.75	1.70	4145712.21	4145712.21\\
-1.75	1.85	5302083.56	5302083.56\\
-1.75	2.00	6917230.47	6917230.47\\
-1.60	-2.80	15256042.18	15256042.18\\
-1.60	-2.65	10888389.86	10888389.86\\
-1.60	-2.50	7927270.99	7927270.99\\
-1.60	-2.35	5887379.18	5887379.18\\
-1.60	-2.20	4460243.52	4460243.52\\
-1.60	-2.05	3446937.39	3446937.39\\
-1.60	-1.90	2717355.63	2717355.63\\
-1.60	-1.75	2185233.37	2185233.37\\
-1.60	-1.60	1792616.67	1792616.67\\
-1.60	-1.45	1500083.03	1500083.03\\
-1.60	-1.30	1280505.43	1280505.43\\
-1.60	-1.15	1115028.10	1115028.10\\
-1.60	-1.00	990440.66	990440.66\\
-1.60	-0.85	897448.15	897448.15\\
-1.60	-0.70	829523.20	829523.20\\
-1.60	-0.55	782142.65	782142.65\\
-1.60	-0.40	752283.71	752283.71\\
-1.60	-0.25	738100.69	738100.69\\
-1.60	-0.10	738733.55	738733.55\\
-1.60	0.05	754220.43	754220.43\\
-1.60	0.20	785501.52	785501.52\\
-1.60	0.35	834514.76	834514.76\\
-1.60	0.50	904397.34	904397.34\\
-1.60	0.65	999822.23	999822.23\\
-1.60	0.80	1127520.81	1127520.81\\
-1.60	0.95	1297073.55	1297073.55\\
-1.60	1.10	1522098.99	1522098.99\\
-1.60	1.25	1822046.50	1822046.50\\
-1.60	1.40	2224919.35	2224919.35\\
-1.60	1.55	2771451.94	2771451.94\\
-1.60	1.70	3521589.19	3521589.19\\
-1.60	1.85	4564658.50	4564658.50\\
-1.60	2.00	6035540.27	6035540.27\\
-1.45	-2.80	8757224.06	8757224.06\\
-1.45	-2.65	6334472.08	6334472.08\\
-1.45	-2.50	4674042.43	4674042.43\\
-1.45	-2.35	3518140.15	3518140.15\\
-1.45	-2.20	2701294.18	2701294.18\\
-1.45	-2.05	2115772.12	2115772.12\\
-1.45	-1.90	1690457.11	1690457.11\\
-1.45	-1.75	1377773.05	1377773.05\\
-1.45	-1.60	1145485.17	1145485.17\\
-1.45	-1.45	971492.66	971492.66\\
-1.45	-1.30	840480.87	840480.87\\
-1.45	-1.15	741744.61	741744.61\\
-1.45	-1.00	667758.20	667758.20\\
-1.45	-0.85	613228.49	613228.49\\
-1.45	-0.70	574465.14	574465.14\\
-1.45	-0.55	548963.27	548963.27\\
-1.45	-0.40	535132.30	535132.30\\
-1.45	-0.25	532129.45	532129.45\\
-1.45	-0.10	539773.66	539773.66\\
-1.45	0.05	558527.22	558527.22\\
-1.45	0.20	589542.69	589542.69\\
-1.45	0.35	634781.76	634781.76\\
-1.45	0.50	697223.27	697223.27\\
-1.45	0.65	781191.61	781191.61\\
-1.45	0.80	892856.22	892856.22\\
-1.45	0.95	1040983.27	1040983.27\\
-1.45	1.10	1238067.23	1238067.23\\
-1.45	1.25	1502045.09	1502045.09\\
-1.45	1.40	1858916.87	1858916.87\\
-1.45	1.55	2346795.41	2346795.41\\
-1.45	1.70	3022238.63	3022238.63\\
-1.45	1.85	3970274.50	3970274.50\\
-1.45	2.00	5320477.09	5320477.09\\
-1.30	-2.80	5078571.40	5078571.40\\
-1.30	-2.65	3723125.79	3723125.79\\
-1.30	-2.50	2784274.93	2784274.93\\
-1.30	-2.35	2124001.33	2124001.33\\
-1.30	-2.20	1652858.65	1652858.65\\
-1.30	-2.05	1312063.65	1312063.65\\
-1.30	-1.90	1062459.42	1062459.42\\
-1.30	-1.75	877623.14	877623.14\\
-1.30	-1.60	739506.57	739506.57\\
-1.30	-1.45	635644.46	635644.46\\
-1.30	-1.30	557345.83	557345.83\\
-1.30	-1.15	498509.58	498509.58\\
-1.30	-1.00	454841.96	454841.96\\
-1.30	-0.85	423336.58	423336.58\\
-1.30	-0.70	401928.99	401928.99\\
-1.30	-0.55	389270.17	389270.17\\
-1.30	-0.40	384583.98	384583.98\\
-1.30	-0.25	387587.27	387587.27\\
-1.30	-0.10	398461.24	398461.24\\
-1.30	0.05	417869.73	417869.73\\
-1.30	0.20	447027.27	447027.27\\
-1.30	0.35	487826.48	487826.48\\
-1.30	0.50	543043.97	543043.97\\
-1.30	0.65	616655.88	616655.88\\
-1.30	0.80	714313.76	714313.76\\
-1.30	0.95	844060.19	844060.19\\
-1.30	1.10	1017410.18	1017410.18\\
-1.30	1.25	1250999.11	1250999.11\\
-1.30	1.40	1569120.10	1569120.10\\
-1.30	1.55	2007676.00	2007676.00\\
-1.30	1.70	2620410.51	2620410.51\\
-1.30	1.85	3488858.03	3488858.03\\
-1.30	2.00	4738441.50	4738441.50\\
-1.15	-2.80	2975549.52	2975549.52\\
-1.15	-2.65	2210830.83	2210830.83\\
-1.15	-2.50	1675645.34	1675645.34\\
-1.15	-2.35	1295528.45	1295528.45\\
-1.15	-2.20	1021762.75	1021762.75\\
-1.15	-2.05	822037.13	822037.13\\
-1.15	-1.90	674638.42	674638.42\\
-1.15	-1.75	564792.55	564792.55\\
-1.15	-1.60	482330.93	482330.93\\
-1.15	-1.45	420184.02	420184.02\\
-1.15	-1.30	373398.19	373398.19\\
-1.15	-1.15	338487.89	338487.89\\
-1.15	-1.00	313005.77	313005.77\\
-1.15	-0.85	295256.72	295256.72\\
-1.15	-0.70	284109.34	284109.34\\
-1.15	-0.55	278874.94	278874.94\\
-1.15	-0.40	279236.21	279236.21\\
-1.15	-0.25	285214.91	285214.91\\
-1.15	-0.10	297174.12	297174.12\\
-1.15	0.05	315855.18	315855.18\\
-1.15	0.20	342454.82	342454.82\\
-1.15	0.35	378753.65	378753.65\\
-1.15	0.50	427315.49	427315.49\\
-1.15	0.65	491788.87	491788.87\\
-1.15	0.80	577360.43	577360.43\\
-1.15	0.95	691438.57	691438.57\\
-1.15	1.10	844692.13	844692.13\\
-1.15	1.25	1052644.08	1052644.08\\
-1.15	1.40	1338144.21	1338144.21\\
-1.15	1.55	1735251.97	1735251.97\\
-1.15	1.70	2295410.86	2295410.86\\
-1.15	1.85	3097394.91	3097394.91\\
-1.15	2.00	4263546.35	4263546.35\\
-1.00	-2.80	1761340.51	1761340.51\\
-1.00	-2.65	1326336.85	1326336.85\\
-1.00	-2.50	1018832.07	1018832.07\\
-1.00	-2.35	798343.22	798343.22\\
-1.00	-2.20	638138.48	638138.48\\
-1.00	-2.05	520329.55	520329.55\\
-1.00	-1.90	432793.08	432793.08\\
-1.00	-1.75	367214.97	367214.97\\
-1.00	-1.60	317832.80	317832.80\\
-1.00	-1.45	280617.86	280617.86\\
-1.00	-1.30	252737.78	252737.78\\
-1.00	-1.15	232200.57	232200.57\\
-1.00	-1.00	217617.93	217617.93\\
-1.00	-0.85	208048.37	208048.37\\
-1.00	-0.70	202895.40	202895.40\\
-1.00	-0.55	201845.17	201845.17\\
-1.00	-0.40	204834.35	204834.35\\
-1.00	-0.25	212043.75	212043.75\\
-1.00	-0.10	223916.66	223916.66\\
-1.00	0.05	241204.62	241204.62\\
-1.00	0.20	265047.12	265047.12\\
-1.00	0.35	297097.37	297097.37\\
-1.00	0.50	339713.49	339713.49\\
-1.00	0.65	396246.13	396246.13\\
-1.00	0.80	471471.61	471471.61\\
-1.00	0.95	572248.04	572248.04\\
-1.00	1.10	708518.73	708518.73\\
-1.00	1.25	894863.15	894863.15\\
-1.00	1.40	1152922.66	1152922.66\\
-1.00	1.55	1515241.94	1515241.94\\
-1.00	1.70	2031430.89	2031430.89\\
-1.00	1.85	2778180.02	2778180.02\\
-1.00	2.00	3875760.87	3875760.87\\
-0.85	-2.80	1053343.41	1053343.41\\
-0.85	-2.65	803901.20	803901.20\\
-0.85	-2.50	625854.87	625854.87\\
-0.85	-2.35	497030.28	497030.28\\
-0.85	-2.20	402652.43	402652.43\\
-0.85	-2.05	332748.47	332748.47\\
-0.85	-1.90	280504.65	280504.65\\
-0.85	-1.75	241213.89	241213.89\\
-0.85	-1.60	211593.75	211593.75\\
-0.85	-1.45	189339.67	189339.67\\
-0.85	-1.30	172829.81	172829.81\\
-0.85	-1.15	160928.87	160928.87\\
-0.85	-1.00	152857.77	152857.77\\
-0.85	-0.85	148108.28	148108.28\\
-0.85	-0.70	146389.33	146389.33\\
-0.85	-0.55	147597.08	147597.08\\
-0.85	-0.40	151804.40	151804.40\\
-0.85	-0.25	159268.26	159268.26\\
-0.85	-0.10	170456.03	170456.03\\
-0.85	0.05	186094.59	186094.59\\
-0.85	0.20	207249.45	207249.45\\
-0.85	0.35	235445.99	235445.99\\
-0.85	0.50	272852.21	272852.21\\
-0.85	0.65	322553.60	322553.60\\
-0.85	0.80	388968.64	388968.64\\
-0.85	0.95	478481.92	478481.92\\
-0.85	1.10	600419.44	600419.44\\
-0.85	1.25	768567.86	768567.86\\
-0.85	1.40	1003570.64	1003570.64\\
-0.85	1.55	1336755.38	1336755.38\\
-0.85	1.70	1816327.64	1816327.64\\
-0.85	1.85	2517530.30	2517530.30\\
-0.85	2.00	3559536.73	3559536.73\\
-0.70	-2.80	636424.90	636424.90\\
-0.70	-2.65	492268.49	492268.49\\
-0.70	-2.50	388414.27	388414.27\\
-0.70	-2.35	312627.07	312627.07\\
-0.70	-2.20	256682.49	256682.49\\
-0.70	-2.05	214983.02	214983.02\\
-0.70	-1.90	183675.12	183675.12\\
-0.70	-1.75	160079.15	160079.15\\
-0.70	-1.60	142317.23	142317.23\\
-0.70	-1.45	129067.96	129067.96\\
-0.70	-1.30	119403.67	119403.67\\
-0.70	-1.15	112682.16	112682.16\\
-0.70	-1.00	108475.31	108475.31\\
-0.70	-0.85	106523.37	106523.37\\
-0.70	-0.70	106708.05	106708.05\\
-0.70	-0.55	109040.46	109040.46\\
-0.70	-0.40	113662.31	113662.31\\
-0.70	-0.25	120860.26	120860.26\\
-0.70	-0.10	131095.82	131095.82\\
-0.70	0.05	145054.90	145054.90\\
-0.70	0.20	163724.71	163724.71\\
-0.70	0.35	188509.97	188509.97\\
-0.70	0.50	221407.67	221407.67\\
-0.70	0.65	265270.70	265270.70\\
-0.70	0.80	324208.32	324208.32\\
-0.70	0.95	404200.91	404200.91\\
-0.70	1.10	514053.92	514053.92\\
-0.70	1.25	666896.33	666896.33\\
-0.70	1.40	882564.07	882564.07\\
-0.70	1.55	1191440.66	1191440.66\\
-0.70	1.70	1640728.99	1640728.99\\
-0.70	1.85	2304833.45	2304833.45\\
-0.70	2.00	3302786.52	3302786.52\\
-0.55	-2.80	388485.54	388485.54\\
-0.55	-2.65	304545.31	304545.31\\
-0.55	-2.50	243538.29	243538.29\\
-0.55	-2.35	198664.75	198664.75\\
-0.55	-2.20	165315.15	165315.15\\
-0.55	-2.05	140327.49	140327.49\\
-0.55	-1.90	121509.75	121509.75\\
-0.55	-1.75	107329.17	107329.17\\
-0.55	-1.60	96708.05	96708.05\\
-0.55	-1.45	88888.54	88888.54\\
-0.55	-1.30	83342.63	83342.63\\
-0.55	-1.15	79712.58	79712.58\\
-0.55	-1.00	77772.27	77772.27\\
-0.55	-0.85	77403.57	77403.57\\
-0.55	-0.70	78584.23	78584.23\\
-0.55	-0.55	81385.70	81385.70\\
-0.55	-0.40	85980.33	85980.33\\
-0.55	-0.25	92659.15	92659.15\\
-0.55	-0.10	101862.85	101862.85\\
-0.55	0.05	114230.37	114230.37\\
-0.55	0.20	130672.92	130672.92\\
-0.55	0.35	152485.26	152485.26\\
-0.55	0.50	181513.27	181513.27\\
-0.55	0.65	220407.91	220407.91\\
-0.55	0.80	273013.57	273013.57\\
-0.55	0.95	344968.59	344968.59\\
-0.55	1.10	444644.70	444644.70\\
-0.55	1.25	584635.20	584635.20\\
-0.55	1.40	784142.61	784142.61\\
-0.55	1.55	1072860.89	1072860.89\\
-0.55	1.70	1497373.11	1497373.11\\
-0.55	1.85	2131841.49	2131841.49\\
-0.55	2.00	3096121.88	3096121.88\\
-0.40	-2.80	239581.39	239581.39\\
-0.40	-2.65	190349.76	190349.76\\
-0.40	-2.50	154272.97	154272.97\\
-0.40	-2.35	127545.63	127545.63\\
-0.40	-2.20	107567.14	107567.14\\
-0.40	-2.05	92540.51	92540.51\\
-0.40	-1.90	81212.42	81212.42\\
-0.40	-1.75	72702.82	72702.82\\
-0.40	-1.60	66392.39	66392.39\\
-0.40	-1.45	61847.71	61847.71\\
-0.40	-1.30	58771.56	58771.56\\
-0.40	-1.15	56970.38	56970.38\\
-0.40	-1.00	56333.82	56333.82\\
-0.40	-0.85	56823.45	56823.45\\
-0.40	-0.70	58468.80	58468.80\\
-0.40	-0.55	61370.42	61370.42\\
-0.40	-0.40	65710.12	65710.12\\
-0.40	-0.25	71770.12	71770.12\\
-0.40	-0.10	79963.78	79963.78\\
-0.40	0.05	90882.71	90882.71\\
-0.40	0.20	105367.69	105367.69\\
-0.40	0.35	124615.46	124615.46\\
-0.40	0.50	150340.03	150340.03\\
-0.40	0.65	185018.69	185018.69\\
-0.40	0.80	232270.92	232270.92\\
-0.40	0.95	297448.88	297448.88\\
-0.40	1.10	388568.95	388568.95\\
-0.40	1.25	517800.08	517800.08\\
-0.40	1.40	703873.14	703873.14\\
-0.40	1.55	976033.96	976033.96\\
-0.40	1.70	1380618.63	1380618.63\\
-0.40	1.85	1992144.29	1992144.29\\
-0.40	2.00	2932284.55	2932284.55\\
-0.25	-2.80	149273.19	149273.19\\
-0.25	-2.65	120199.66	120199.66\\
-0.25	-2.50	98733.14	98733.14\\
-0.25	-2.35	82729.59	82729.59\\
-0.25	-2.20	70712.65	70712.65\\
-0.25	-2.05	61655.46	61655.46\\
-0.25	-1.90	54838.34	54838.34\\
-0.25	-1.75	49754.83	49754.83\\
-0.25	-1.60	46049.46	46049.46\\
-0.25	-1.45	43476.25	43476.25\\
-0.25	-1.30	41871.43	41871.43\\
-0.25	-1.15	41135.98	41135.98\\
-0.25	-1.00	41225.33	41225.33\\
-0.25	-0.85	42144.87	42144.87\\
-0.25	-0.70	43950.47	43950.47\\
-0.25	-0.55	46754.19	46754.19\\
-0.25	-0.40	50735.96	50735.96\\
-0.25	-0.25	56162.90	56162.90\\
-0.25	-0.10	63419.29	63419.29\\
-0.25	0.05	73051.89	73051.89\\
-0.25	0.20	85838.05	85838.05\\
-0.25	0.35	102888.42	102888.42\\
-0.25	0.50	125803.11	125803.11\\
-0.25	0.65	156911.42	156911.42\\
-0.25	0.80	199643.86	199643.86\\
-0.25	0.95	259116.84	259116.84\\
-0.25	1.10	343062.76	343062.76\\
-0.25	1.25	463329.32	463329.32\\
-0.25	1.40	638328.51	638328.51\\
-0.25	1.55	897091.96	897091.96\\
-0.25	1.70	1286079.89	1286079.89\\
-0.25	1.85	1880776.51	1880776.51\\
-0.25	2.00	2805722.38	2805722.38\\
-0.10	-2.80	93963.91	93963.91\\
-0.10	-2.65	76683.99	76683.99\\
-0.10	-2.50	63839.07	63839.07\\
-0.10	-2.35	54213.41	54213.41\\
-0.10	-2.20	46964.00	46964.00\\
-0.10	-2.05	41501.30	41501.30\\
-0.10	-1.90	37410.77	37410.77\\
-0.10	-1.75	34400.90	34400.90\\
-0.10	-1.60	32268.68	32268.68\\
-0.10	-1.45	30876.70	30876.70\\
-0.10	-1.30	30138.31	30138.31\\
-0.10	-1.15	30008.56	30008.56\\
-0.10	-1.00	30479.62	30479.62\\
-0.10	-0.85	31580.01	31580.01\\
-0.10	-0.70	33377.46	33377.46\\
-0.10	-0.55	35985.92	35985.92\\
-0.10	-0.40	39577.66	39577.66\\
-0.10	-0.25	44402.34	44402.34\\
-0.10	-0.10	50815.93	50815.93\\
-0.10	0.05	59324.24	59324.24\\
-0.10	0.20	70648.47	70648.47\\
-0.10	0.35	85824.56	85824.56\\
-0.10	0.50	106355.18	106355.18\\
-0.10	0.65	134444.80	134444.80\\
-0.10	0.80	173367.47	173367.47\\
-0.10	0.95	228049.68	228049.68\\
-0.10	1.10	306005.75	306005.75\\
-0.10	1.25	418859.12	418859.12\\
-0.10	1.40	584850.16	584850.16\\
-0.10	1.55	833027.85	833027.85\\
-0.10	1.70	1210354.77	1210354.77\\
-0.10	1.85	1793924.30	1793924.30\\
-0.10	2.00	2712275.50	2712275.50\\
0.05	-2.80	59757.29	59757.29\\
0.05	-2.65	49426.13	49426.13\\
0.05	-2.50	41702.37	41702.37\\
0.05	-2.35	35892.44	35892.44\\
0.05	-2.20	31512.55	31512.55\\
0.05	-2.05	28222.95	28222.95\\
0.05	-1.90	25784.55	25784.55\\
0.05	-1.75	24030.06	24030.06\\
0.05	-1.60	22844.86	22844.86\\
0.05	-1.45	22154.42	22154.42\\
0.05	-1.30	21916.46	21916.46\\
0.05	-1.15	22116.63	22116.63\\
0.05	-1.00	22766.99	22766.99\\
0.05	-0.85	23907.30	23907.30\\
0.05	-0.70	25609.06	25609.06\\
0.05	-0.55	27983.05	27983.05\\
0.05	-0.40	31191.39	31191.39\\
0.05	-0.25	35466.04	35466.04\\
0.05	-0.10	41136.65	41136.65\\
0.05	0.05	48672.48	48672.48\\
0.05	0.20	58745.71	58745.71\\
0.05	0.35	72328.12	72328.12\\
0.05	0.50	90839.86	90839.86\\
0.05	0.65	116381.51	116381.51\\
0.05	0.80	152100.19	152100.19\\
0.05	0.95	202774.72	202774.72\\
0.05	1.10	275763.08	275763.08\\
0.05	1.25	382557.47	382557.47\\
0.05	1.40	541371.65	541371.65\\
0.05	1.55	781506.50	781506.50\\
0.05	1.70	1150821.44	1150821.44\\
0.05	1.85	1728707.65	1728707.65\\
0.05	2.00	2648947.94	2648947.94\\
0.20	-2.80	38394.69	38394.69\\
0.20	-2.65	32185.42	32185.42\\
0.20	-2.50	27522.34	27522.34\\
0.20	-2.35	24007.66	24007.66\\
0.20	-2.20	21362.53	21362.53\\
0.20	-2.05	19390.71	19390.71\\
0.20	-1.90	17954.48	17954.48\\
0.20	-1.75	16958.62	16958.62\\
0.20	-1.60	16339.78	16339.78\\
0.20	-1.45	16059.81	16059.81\\
0.20	-1.30	16101.73	16101.73\\
0.20	-1.15	16468.09	16468.09\\
0.20	-1.00	17181.14	17181.14\\
0.20	-0.85	18285.18	18285.18\\
0.20	-0.70	19851.10	19851.10\\
0.20	-0.55	21984.07	21984.07\\
0.20	-0.40	24835.33	24835.33\\
0.20	-0.25	28620.03	28620.03\\
0.20	-0.10	33644.07	33644.07\\
0.20	0.05	40344.58	40344.58\\
0.20	0.20	49351.48	49351.48\\
0.20	0.35	61581.93	61581.93\\
0.20	0.50	78387.13	78387.13\\
0.20	0.65	101782.82	101782.82\\
0.20	0.80	134816.31	134816.31\\
0.20	0.95	182158.18	182158.18\\
0.20	1.10	251069.04	251069.04\\
0.20	1.25	353000.98	353000.98\\
0.20	1.40	506287.17	506287.17\\
0.20	1.55	740723.62	740723.62\\
0.20	1.70	1105487.21	1105487.21\\
0.20	1.85	1683020.92	1683020.92\\
0.20	2.00	2613747.12	2613747.12\\
0.35	-2.80	24923.10	24923.10\\
0.35	-2.65	21174.45	21174.45\\
0.35	-2.50	18351.03	18351.03\\
0.35	-2.35	16223.60	16223.60\\
0.35	-2.20	14630.94	14630.94\\
0.35	-2.05	13459.70	13459.70\\
0.35	-1.90	12630.97	12630.97\\
0.35	-1.75	12091.40	12091.40\\
0.35	-1.60	11807.41	11807.41\\
0.35	-1.45	11761.72	11761.72\\
0.35	-1.30	11951.58	11951.58\\
0.35	-1.15	12388.48	12388.48\\
0.35	-1.00	13099.33	13099.33\\
0.35	-0.85	14129.22	14129.22\\
0.35	-0.70	15546.26	15546.26\\
0.35	-0.55	17449.05	17449.05\\
0.35	-0.40	19978.17	19978.17\\
0.35	-0.25	23333.40	23333.40\\
0.35	-0.10	27799.61	27799.61\\
0.35	0.05	33786.06	33786.06\\
0.35	0.20	41886.55	41886.55\\
0.35	0.35	52972.44	52972.44\\
0.35	0.50	68338.21	68338.21\\
0.35	0.65	89932.26	89932.26\\
0.35	0.80	120727.34	120727.34\\
0.35	0.95	165323.30	165323.30\\
0.35	1.10	230940.83	230940.83\\
0.35	1.25	329083.16	329083.16\\
0.35	1.40	478353.39	478353.39\\
0.35	1.55	709300.57	709300.57\\
0.35	1.70	1072877.21	1072877.21\\
0.35	1.85	1655419.23	1655419.23\\
0.35	2.00	2605578.91	2605578.91\\
0.50	-2.80	16344.94	16344.94\\
0.50	-2.65	14073.94	14073.94\\
0.50	-2.50	12361.93	12361.93\\
0.50	-2.35	11076.31	11076.31\\
0.50	-2.20	10123.77	10123.77\\
0.50	-2.05	9439.04	9439.04\\
0.50	-1.90	8977.41	8977.41\\
0.50	-1.75	8709.90	8709.90\\
0.50	-1.60	8620.12	8620.12\\
0.50	-1.45	8702.65	8702.65\\
0.50	-1.30	8962.48	8962.48\\
0.50	-1.15	9415.50	9415.50\\
0.50	-1.00	10090.13	10090.13\\
0.50	-0.85	11030.32	11030.32\\
0.50	-0.70	12300.36	12300.36\\
0.50	-0.55	13992.19	13992.19\\
0.50	-0.40	16236.48	16236.48\\
0.50	-0.25	19219.25	19219.25\\
0.50	-0.10	23207.01	23207.01\\
0.50	0.05	28585.14	28585.14\\
0.50	0.20	35916.96	35916.96\\
0.50	0.35	46035.95	46035.95\\
0.50	0.50	60191.19	60191.19\\
0.50	0.65	80279.94	80279.94\\
0.50	0.80	109224.33	109224.33\\
0.50	0.95	151589.79	151589.79\\
0.50	1.10	214614.38	214614.38\\
0.50	1.25	309945.93	309945.93\\
0.50	1.40	456616.19	456616.19\\
0.50	1.55	686206.68	686206.68\\
0.50	1.70	1051954.22	1051954.22\\
0.50	1.85	1645042.02	1645042.02\\
0.50	2.00	2624190.83	2624190.83\\
0.65	-2.80	10829.67	10829.67\\
0.65	-2.65	9450.82	9450.82\\
0.65	-2.50	8413.22	8413.22\\
0.65	-2.35	7640.00	7640.00\\
0.65	-2.20	7077.22	7077.22\\
0.65	-2.05	6687.60	6687.60\\
0.65	-1.90	6446.38	6446.38\\
0.65	-1.75	6338.70	6338.70\\
0.65	-1.60	6358.03	6358.03\\
0.65	-1.45	6505.54	6505.54\\
0.65	-1.30	6790.19	6790.19\\
0.65	-1.15	7229.68	7229.68\\
0.65	-1.00	7852.26	7852.26\\
0.65	-0.85	8699.78	8699.78\\
0.65	-0.70	9832.41	9832.41\\
0.65	-0.55	11335.75	11335.75\\
0.65	-0.40	13331.49	13331.49\\
0.65	-0.25	15993.57	15993.57\\
0.65	-0.10	19572.69	19572.69\\
0.65	0.05	24433.95	24433.95\\
0.65	0.20	31115.38	31115.38\\
0.65	0.35	40419.86	40419.86\\
0.65	0.50	53561.51	53561.51\\
0.65	0.65	72401.76	72401.76\\
0.65	0.80	99835.18	99835.18\\
0.65	0.95	140428.87	140428.87\\
0.65	1.10	201496.46	201496.46\\
0.65	1.25	294928.49	294928.49\\
0.65	1.40	440356.37	440356.37\\
0.65	1.55	670702.77	670702.77\\
0.65	1.70	1042063.48	1042063.48\\
0.65	1.85	1651568.21	1651568.21\\
0.65	2.00	2670158.95	2670158.95\\
0.80	-2.80	7249.32	7249.32\\
0.80	-2.65	6411.72	6411.72\\
0.80	-2.50	5784.81	5784.81\\
0.80	-2.35	5324.05	5324.05\\
0.80	-2.20	4998.43	4998.43\\
0.80	-2.05	4787.00	4787.00\\
0.80	-1.90	4676.62	4676.62\\
0.80	-1.75	4660.56	4660.56\\
0.80	-1.60	4737.86	4737.86\\
0.80	-1.45	4913.21	4913.21\\
0.80	-1.30	5197.40	5197.40\\
0.80	-1.15	5608.49	5608.49\\
0.80	-1.00	6173.67	6173.67\\
0.80	-0.85	6932.33	6932.33\\
0.80	-0.70	7940.60	7940.60\\
0.80	-0.55	9278.24	9278.24\\
0.80	-0.40	11059.01	11059.01\\
0.80	-0.25	13446.37	13446.37\\
0.80	-0.10	16677.54	16677.54\\
0.80	0.05	21100.73	21100.73\\
0.80	0.20	27233.35	27233.35\\
0.80	0.35	35854.45	35854.45\\
0.80	0.50	48152.99	48152.99\\
0.80	0.65	65969.27	65969.27\\
0.80	0.80	92193.09	92193.09\\
0.80	0.95	131429.64	131429.64\\
0.80	1.10	191128.98	191128.98\\
0.80	1.25	283529.35	283529.35\\
0.80	1.40	429049.87	429049.87\\
0.80	1.55	662301.55	662301.55\\
0.80	1.70	1042898.48	1042898.48\\
0.80	1.85	1675199.58	1675199.58\\
0.80	2.00	2744917.76	2744917.76\\
0.95	-2.80	4902.64	4902.64\\
0.95	-2.65	4394.70	4394.70\\
0.95	-2.50	4018.52	4018.52\\
0.95	-2.35	3748.37	3748.37\\
0.95	-2.20	3566.61	3566.61\\
0.95	-2.05	3461.84	3461.84\\
0.95	-1.90	3427.66	3427.66\\
0.95	-1.75	3461.99	3461.99\\
0.95	-1.60	3566.92	3566.92\\
0.95	-1.45	3748.85	3748.85\\
0.95	-1.30	4019.21	4019.21\\
0.95	-1.15	4395.64	4395.64\\
0.95	-1.00	4903.91	4903.91\\
0.95	-0.85	5580.85	5580.85\\
0.95	-0.70	6478.83	6478.83\\
0.95	-0.55	7672.40	7672.40\\
0.95	-0.40	9268.38	9268.38\\
0.95	-0.25	11421.28	11421.28\\
0.95	-0.10	14357.02	14357.02\\
0.95	0.05	18409.92	18409.92\\
0.95	0.20	24081.18	24081.18\\
0.95	0.35	32132.30	32132.30\\
0.95	0.50	43736.51	43736.51\\
0.95	0.65	60727.41	60727.41\\
0.95	0.80	86012.91	86012.91\\
0.95	0.95	124274.15	124274.15\\
0.95	1.10	183162.35	183162.35\\
0.95	1.25	275378.39	275378.39\\
0.95	1.40	422339.58	422339.58\\
0.95	1.55	660742.08	660742.08\\
0.95	1.70	1054485.01	1054485.01\\
0.95	1.85	1716671.17	1716671.17\\
0.95	2.00	2850834.99	2850834.99\\
1.10	-2.80	3349.76	3349.76\\
1.10	-2.65	3043.23	3043.23\\
1.10	-2.50	2820.29	2820.29\\
1.10	-2.35	2666.19	2666.19\\
1.10	-2.20	2571.15	2571.15\\
1.10	-2.05	2529.31	2529.31\\
1.10	-1.90	2538.13	2538.13\\
1.10	-1.75	2598.15	2598.15\\
1.10	-1.60	2713.02	2713.02\\
1.10	-1.45	2889.89	2889.89\\
1.10	-1.30	3140.12	3140.12\\
1.10	-1.15	3480.57	3480.57\\
1.10	-1.00	3935.43	3935.43\\
1.10	-0.85	4539.12	4539.12\\
1.10	-0.70	5340.61	5340.61\\
1.10	-0.55	6409.84	6409.84\\
1.10	-0.40	7847.70	7847.70\\
1.10	-0.25	9801.11	9801.11\\
1.10	-0.10	12486.68	12486.68\\
1.10	0.05	16227.69	16227.69\\
1.10	0.20	21513.19	21513.19\\
1.10	0.35	29093.17	29093.17\\
1.10	0.50	40134.29	40134.29\\
1.10	0.65	56477.88	56477.88\\
1.10	0.80	81073.60	81073.60\\
1.10	0.95	118718.60	118718.60\\
1.10	1.10	177335.78	177335.78\\
1.10	1.25	270216.71	270216.71\\
1.10	1.40	420016.46	420016.46\\
1.10	1.55	665976.15	665976.15\\
1.10	1.70	1077182.54	1077182.54\\
1.10	1.85	1777289.58	1777289.58\\
1.10	2.00	2991336.99	2991336.99\\
1.25	-2.80	2312.32	2312.32\\
1.25	-2.65	2129.08	2129.08\\
1.25	-2.50	1999.74	1999.74\\
1.25	-2.35	1915.99	1915.99\\
1.25	-2.20	1872.62	1872.62\\
1.25	-2.05	1867.01	1867.01\\
1.25	-1.90	1898.81	1898.81\\
1.25	-1.75	1969.94	1969.94\\
1.25	-1.60	2084.80	2084.80\\
1.25	-1.45	2250.68	2250.68\\
1.25	-1.30	2478.57	2478.57\\
1.25	-1.15	2784.38	2784.38\\
1.25	-1.00	3190.74	3190.74\\
1.25	-0.85	3729.88	3729.88\\
1.25	-0.70	4447.70	4447.70\\
1.25	-0.55	5410.21	5410.21\\
1.25	-0.40	6713.22	6713.22\\
1.25	-0.25	8497.41	8497.41\\
1.25	-0.10	10971.86	10971.86\\
1.25	0.05	14451.47	14451.47\\
1.25	0.20	19417.01	19417.01\\
1.25	0.35	26612.82	26612.82\\
1.25	0.50	37208.10	37208.10\\
1.25	0.65	53066.75	53066.75\\
1.25	0.80	77205.06	77205.06\\
1.25	0.95	114579.59	114579.59\\
1.25	1.10	173463.08	173463.08\\
1.25	1.25	267882.95	267882.95\\
1.25	1.40	422008.66	422008.66\\
1.25	1.55	678165.84	678165.84\\
1.25	1.70	1111702.86	1111702.86\\
1.25	1.85	1859001.72	1859001.72\\
1.25	2.00	3171094.05	3171094.05\\
1.40	-2.80	1612.62	1612.62\\
1.40	-2.65	1504.87	1504.87\\
1.40	-2.50	1432.52	1432.52\\
1.40	-2.35	1391.05	1391.05\\
1.40	-2.20	1377.92	1377.92\\
1.40	-2.05	1392.33	1392.33\\
1.40	-1.90	1435.15	1435.15\\
1.40	-1.75	1509.02	1509.02\\
1.40	-1.60	1618.55	1618.55\\
1.40	-1.45	1770.92	1770.92\\
1.40	-1.30	1976.55	1976.55\\
1.40	-1.15	2250.38	2250.38\\
1.40	-1.00	2613.62	2613.62\\
1.40	-0.85	3096.47	3096.47\\
1.40	-0.70	3742.23	3742.23\\
1.40	-0.55	4613.51	4613.51\\
1.40	-0.40	5801.90	5801.90\\
1.40	-0.25	7443.00	7443.00\\
1.40	-0.10	9740.11	9740.11\\
1.40	0.05	13002.24	13002.24\\
1.40	0.20	17705.59	17705.59\\
1.40	0.35	24594.68	24594.68\\
1.40	0.50	34850.58	34850.58\\
1.40	0.65	50375.25	50375.25\\
1.40	0.80	74278.41	74278.41\\
1.40	0.95	111723.96	111723.96\\
1.40	1.10	171422.67	171422.67\\
1.40	1.25	268304.81	268304.81\\
1.40	1.40	428377.78	428377.78\\
1.40	1.55	697691.88	697691.88\\
1.40	1.70	1159147.38	1159147.38\\
1.40	1.85	1964499.45	1964499.45\\
1.40	2.00	3396279.51	3396279.51\\
1.55	-2.80	1136.23	1136.23\\
1.55	-2.65	1074.62	1074.62\\
1.55	-2.50	1036.77	1036.77\\
1.55	-2.35	1020.34	1020.34\\
1.55	-2.20	1024.35	1024.35\\
1.55	-2.05	1049.03	1049.03\\
1.55	-1.90	1095.89	1095.89\\
1.55	-1.75	1167.84	1167.84\\
1.55	-1.60	1269.52	1269.52\\
1.55	-1.45	1407.77	1407.77\\
1.55	-1.30	1592.45	1592.45\\
1.55	-1.15	1837.53	1837.53\\
1.55	-1.00	2162.94	2162.94\\
1.55	-0.85	2597.11	2597.11\\
1.55	-0.70	3181.09	3181.09\\
1.55	-0.55	3974.65	3974.65\\
1.55	-0.40	5065.94	5065.94\\
1.55	-0.25	6586.58	6586.58\\
1.55	-0.10	8735.71	8735.71\\
1.55	0.05	11818.83	11818.83\\
1.55	0.20	16311.32	16311.32\\
1.55	0.35	22963.70	22963.70\\
1.55	0.50	32978.66	32978.66\\
1.55	0.65	48312.82	48312.82\\
1.55	0.80	72198.80	72198.80\\
1.55	0.95	110061.61	110061.61\\
1.55	1.10	171151.21	171151.21\\
1.55	1.25	271495.33	271495.33\\
1.55	1.40	439322.07	439322.07\\
1.55	1.55	725173.52	725173.52\\
1.55	1.70	1221065.93	1221065.93\\
1.55	1.85	2097367.57	2097367.57\\
1.55	2.00	3674923.02	3674923.02\\
1.70	-2.80	808.82	808.82\\
1.70	-2.65	775.29	775.29\\
1.70	-2.50	758.07	758.07\\
1.70	-2.35	756.13	756.13\\
1.70	-2.20	769.35	769.35\\
1.70	-2.05	798.52	798.52\\
1.70	-1.90	845.45	845.45\\
1.70	-1.75	913.11	913.11\\
1.70	-1.60	1006.01	1006.01\\
1.70	-1.45	1130.62	1130.62\\
1.70	-1.30	1296.20	1296.20\\
1.70	-1.15	1515.88	1515.88\\
1.70	-1.00	1808.40	1808.40\\
1.70	-0.85	2200.72	2200.72\\
1.70	-0.70	2731.94	2731.94\\
1.70	-0.55	3459.53	3459.53\\
1.70	-0.40	4468.90	4468.90\\
1.70	-0.25	5888.75	5888.75\\
1.70	-0.10	7915.58	7915.58\\
1.70	0.05	10853.79	10853.79\\
1.70	0.20	15181.62	15181.62\\
1.70	0.35	21661.74	21661.74\\
1.70	0.50	31528.73	31528.73\\
1.70	0.65	46812.09	46812.09\\
1.70	0.80	70900.26	70900.26\\
1.70	0.95	109540.80	109540.80\\
1.70	1.10	172640.32	172640.32\\
1.70	1.25	277553.56	277553.56\\
1.70	1.40	455186.77	455186.77\\
1.70	1.55	761501.45	761501.45\\
1.70	1.70	1299541.31	1299541.31\\
1.70	1.85	2262287.04	2262287.04\\
1.70	2.00	4017386.23	4017386.23\\
1.85	-2.80	581.69	581.69\\
1.85	-2.65	565.09	565.09\\
1.85	-2.50	560.00	560.00\\
1.85	-2.35	566.11	566.11\\
1.85	-2.20	583.78	583.78\\
1.85	-2.05	614.09	614.09\\
1.85	-1.90	658.95	658.95\\
1.85	-1.75	721.30	721.30\\
1.85	-1.60	805.41	805.41\\
1.85	-1.45	917.39	917.39\\
1.85	-1.30	1065.93	1065.93\\
1.85	-1.15	1263.41	1263.41\\
1.85	-1.00	1527.56	1527.56\\
1.85	-0.85	1884.03	1884.03\\
1.85	-0.70	2370.38	2370.38\\
1.85	-0.55	3042.19	3042.19\\
1.85	-0.40	3982.83	3982.83\\
1.85	-0.25	5319.07	5319.07\\
1.85	-0.10	7246.33	7246.33\\
1.85	0.05	10070.22	10070.22\\
1.85	0.20	14275.72	14275.72\\
1.85	0.35	20644.06	20644.06\\
1.85	0.50	30453.03	30453.03\\
1.85	0.65	45825.19	45825.19\\
1.85	0.80	70342.25	70342.25\\
1.85	0.95	110145.43	110145.43\\
1.85	1.10	175936.11	175936.11\\
1.85	1.25	286669.67	286669.67\\
1.85	1.40	476482.29	476482.29\\
1.85	1.55	807885.93	807885.93\\
1.85	1.70	1397306.20	1397306.20\\
1.85	1.85	2465309.14	2465309.14\\
1.85	2.00	4437000.19	4437000.19\\
2.00	-2.80	422.64	422.64\\
2.00	-2.65	416.13	416.13\\
2.00	-2.50	417.95	417.95\\
2.00	-2.35	428.21	428.21\\
2.00	-2.20	447.53	447.53\\
2.00	-2.05	477.12	477.12\\
2.00	-1.90	518.89	518.89\\
2.00	-1.75	575.65	575.65\\
2.00	-1.60	651.45	651.45\\
2.00	-1.45	752.04	752.04\\
2.00	-1.30	885.60	885.60\\
2.00	-1.15	1063.84	1063.84\\
2.00	-1.00	1303.62	1303.62\\
2.00	-0.85	1629.54	1629.54\\
2.00	-0.70	2077.86	2077.86\\
2.00	-0.55	2702.75	2702.75\\
2.00	-0.40	3586.19	3586.19\\
2.00	-0.25	4854.00	4854.00\\
2.00	-0.10	6702.00	6702.00\\
2.00	0.05	9439.46	9439.46\\
2.00	0.20	13562.14	13562.14\\
2.00	0.35	19876.85	19876.85\\
2.00	0.50	29717.01	29717.01\\
2.00	0.65	45321.16	45321.16\\
2.00	0.80	70507.48	70507.48\\
2.00	0.95	111894.20	111894.20\\
2.00	1.10	181141.63	181141.63\\
2.00	1.25	299134.99	299134.99\\
2.00	1.40	503911.67	503911.67\\
2.00	1.55	865924.20	865924.20\\
2.00	1.70	1517901.53	1517901.53\\
2.00	1.85	2714223.36	2714223.36\\
2.00	2.00	4950919.07	4950919.07\\
};

\addplot3[%
surf,
fill opacity=0.65, draw opacity=0.70, fill=white!65!black, faceted color=white!15!black, z buffer=sort, colormap={mymap}{[1pt] rgb(0pt)=(0.2422,0.1504,0.6603); rgb(1pt)=(0.2444,0.1534,0.6728); rgb(2pt)=(0.2464,0.1569,0.6847); rgb(3pt)=(0.2484,0.1607,0.6961); rgb(4pt)=(0.2503,0.1648,0.7071); rgb(5pt)=(0.2522,0.1689,0.7179); rgb(6pt)=(0.254,0.1732,0.7286); rgb(7pt)=(0.2558,0.1773,0.7393); rgb(8pt)=(0.2576,0.1814,0.7501); rgb(9pt)=(0.2594,0.1854,0.761); rgb(11pt)=(0.2628,0.1932,0.7828); rgb(12pt)=(0.2645,0.1972,0.7937); rgb(13pt)=(0.2661,0.2011,0.8043); rgb(14pt)=(0.2676,0.2052,0.8148); rgb(15pt)=(0.2691,0.2094,0.8249); rgb(16pt)=(0.2704,0.2138,0.8346); rgb(17pt)=(0.2717,0.2184,0.8439); rgb(18pt)=(0.2729,0.2231,0.8528); rgb(19pt)=(0.274,0.228,0.8612); rgb(20pt)=(0.2749,0.233,0.8692); rgb(21pt)=(0.2758,0.2382,0.8767); rgb(22pt)=(0.2766,0.2435,0.884); rgb(23pt)=(0.2774,0.2489,0.8908); rgb(24pt)=(0.2781,0.2543,0.8973); rgb(25pt)=(0.2788,0.2598,0.9035); rgb(26pt)=(0.2794,0.2653,0.9094); rgb(27pt)=(0.2798,0.2708,0.915); rgb(28pt)=(0.2802,0.2764,0.9204); rgb(29pt)=(0.2806,0.2819,0.9255); rgb(30pt)=(0.2809,0.2875,0.9305); rgb(31pt)=(0.2811,0.293,0.9352); rgb(32pt)=(0.2813,0.2985,0.9397); rgb(33pt)=(0.2814,0.304,0.9441); rgb(34pt)=(0.2814,0.3095,0.9483); rgb(35pt)=(0.2813,0.315,0.9524); rgb(36pt)=(0.2811,0.3204,0.9563); rgb(37pt)=(0.2809,0.3259,0.96); rgb(38pt)=(0.2807,0.3313,0.9636); rgb(39pt)=(0.2803,0.3367,0.967); rgb(40pt)=(0.2798,0.3421,0.9702); rgb(41pt)=(0.2791,0.3475,0.9733); rgb(42pt)=(0.2784,0.3529,0.9763); rgb(43pt)=(0.2776,0.3583,0.9791); rgb(44pt)=(0.2766,0.3638,0.9817); rgb(45pt)=(0.2754,0.3693,0.984); rgb(46pt)=(0.2741,0.3748,0.9862); rgb(47pt)=(0.2726,0.3804,0.9881); rgb(48pt)=(0.271,0.386,0.9898); rgb(49pt)=(0.2691,0.3916,0.9912); rgb(50pt)=(0.267,0.3973,0.9924); rgb(51pt)=(0.2647,0.403,0.9935); rgb(52pt)=(0.2621,0.4088,0.9946); rgb(53pt)=(0.2591,0.4145,0.9955); rgb(54pt)=(0.2556,0.4203,0.9965); rgb(55pt)=(0.2517,0.4261,0.9974); rgb(56pt)=(0.2473,0.4319,0.9983); rgb(57pt)=(0.2424,0.4378,0.9991); rgb(58pt)=(0.2369,0.4437,0.9996); rgb(59pt)=(0.2311,0.4497,0.9995); rgb(60pt)=(0.225,0.4559,0.9985); rgb(61pt)=(0.2189,0.462,0.9968); rgb(62pt)=(0.2128,0.4682,0.9948); rgb(63pt)=(0.2066,0.4743,0.9926); rgb(64pt)=(0.2006,0.4803,0.9906); rgb(65pt)=(0.195,0.4861,0.9887); rgb(66pt)=(0.1903,0.4919,0.9867); rgb(67pt)=(0.1869,0.4975,0.9844); rgb(68pt)=(0.1847,0.503,0.9819); rgb(69pt)=(0.1831,0.5084,0.9793); rgb(70pt)=(0.1818,0.5138,0.9766); rgb(71pt)=(0.1806,0.5191,0.9738); rgb(72pt)=(0.1795,0.5244,0.9709); rgb(73pt)=(0.1785,0.5296,0.9677); rgb(74pt)=(0.1778,0.5349,0.9641); rgb(75pt)=(0.1773,0.5401,0.9602); rgb(76pt)=(0.1768,0.5452,0.956); rgb(77pt)=(0.1764,0.5504,0.9516); rgb(78pt)=(0.1755,0.5554,0.9473); rgb(79pt)=(0.174,0.5605,0.9432); rgb(80pt)=(0.1716,0.5655,0.9393); rgb(81pt)=(0.1686,0.5705,0.9357); rgb(82pt)=(0.1649,0.5755,0.9323); rgb(83pt)=(0.161,0.5805,0.9289); rgb(84pt)=(0.1573,0.5854,0.9254); rgb(85pt)=(0.154,0.5902,0.9218); rgb(86pt)=(0.1513,0.595,0.9182); rgb(87pt)=(0.1492,0.5997,0.9147); rgb(88pt)=(0.1475,0.6043,0.9113); rgb(89pt)=(0.1461,0.6089,0.908); rgb(90pt)=(0.1446,0.6135,0.905); rgb(91pt)=(0.1429,0.618,0.9022); rgb(92pt)=(0.1408,0.6226,0.8998); rgb(93pt)=(0.1383,0.6272,0.8975); rgb(94pt)=(0.1354,0.6317,0.8953); rgb(95pt)=(0.1321,0.6363,0.8932); rgb(96pt)=(0.1288,0.6408,0.891); rgb(97pt)=(0.1253,0.6453,0.8887); rgb(98pt)=(0.1219,0.6497,0.8862); rgb(99pt)=(0.1185,0.6541,0.8834); rgb(100pt)=(0.1152,0.6584,0.8804); rgb(101pt)=(0.1119,0.6627,0.877); rgb(102pt)=(0.1085,0.6669,0.8734); rgb(103pt)=(0.1048,0.671,0.8695); rgb(104pt)=(0.1009,0.675,0.8653); rgb(105pt)=(0.0964,0.6789,0.8609); rgb(106pt)=(0.0914,0.6828,0.8562); rgb(107pt)=(0.0855,0.6865,0.8513); rgb(108pt)=(0.0789,0.6902,0.8462); rgb(109pt)=(0.0713,0.6938,0.8409); rgb(110pt)=(0.0628,0.6972,0.8355); rgb(111pt)=(0.0535,0.7006,0.8299); rgb(112pt)=(0.0433,0.7039,0.8242); rgb(113pt)=(0.0328,0.7071,0.8183); rgb(114pt)=(0.0234,0.7103,0.8124); rgb(115pt)=(0.0155,0.7133,0.8064); rgb(116pt)=(0.0091,0.7163,0.8003); rgb(117pt)=(0.0046,0.7192,0.7941); rgb(118pt)=(0.0019,0.722,0.7878); rgb(119pt)=(0.0009,0.7248,0.7815); rgb(120pt)=(0.0018,0.7275,0.7752); rgb(121pt)=(0.0046,0.7301,0.7688); rgb(122pt)=(0.0094,0.7327,0.7623); rgb(123pt)=(0.0162,0.7352,0.7558); rgb(124pt)=(0.0253,0.7376,0.7492); rgb(125pt)=(0.0369,0.74,0.7426); rgb(126pt)=(0.0504,0.7423,0.7359); rgb(127pt)=(0.0638,0.7446,0.7292); rgb(128pt)=(0.077,0.7468,0.7224); rgb(129pt)=(0.0899,0.7489,0.7156); rgb(130pt)=(0.1023,0.751,0.7088); rgb(131pt)=(0.1141,0.7531,0.7019); rgb(132pt)=(0.1252,0.7552,0.695); rgb(133pt)=(0.1354,0.7572,0.6881); rgb(134pt)=(0.1448,0.7593,0.6812); rgb(135pt)=(0.1532,0.7614,0.6741); rgb(136pt)=(0.1609,0.7635,0.6671); rgb(137pt)=(0.1678,0.7656,0.6599); rgb(138pt)=(0.1741,0.7678,0.6527); rgb(139pt)=(0.1799,0.7699,0.6454); rgb(140pt)=(0.1853,0.7721,0.6379); rgb(141pt)=(0.1905,0.7743,0.6303); rgb(142pt)=(0.1954,0.7765,0.6225); rgb(143pt)=(0.2003,0.7787,0.6146); rgb(144pt)=(0.2061,0.7808,0.6065); rgb(145pt)=(0.2118,0.7828,0.5983); rgb(146pt)=(0.2178,0.7849,0.5899); rgb(147pt)=(0.2244,0.7869,0.5813); rgb(148pt)=(0.2318,0.7887,0.5725); rgb(149pt)=(0.2401,0.7905,0.5636); rgb(150pt)=(0.2491,0.7922,0.5546); rgb(151pt)=(0.2589,0.7937,0.5454); rgb(152pt)=(0.2695,0.7951,0.536); rgb(153pt)=(0.2809,0.7964,0.5266); rgb(154pt)=(0.2929,0.7975,0.517); rgb(155pt)=(0.3052,0.7985,0.5074); rgb(156pt)=(0.3176,0.7994,0.4975); rgb(157pt)=(0.3301,0.8002,0.4876); rgb(158pt)=(0.3424,0.8009,0.4774); rgb(159pt)=(0.3548,0.8016,0.4669); rgb(160pt)=(0.3671,0.8021,0.4563); rgb(161pt)=(0.3795,0.8026,0.4454); rgb(162pt)=(0.3921,0.8029,0.4344); rgb(163pt)=(0.405,0.8031,0.4233); rgb(164pt)=(0.4184,0.803,0.4122); rgb(165pt)=(0.4322,0.8028,0.4013); rgb(166pt)=(0.4463,0.8024,0.3904); rgb(167pt)=(0.4608,0.8018,0.3797); rgb(168pt)=(0.4753,0.8011,0.3691); rgb(169pt)=(0.4899,0.8002,0.3586); rgb(170pt)=(0.5044,0.7993,0.348); rgb(171pt)=(0.5187,0.7982,0.3374); rgb(172pt)=(0.5329,0.797,0.3267); rgb(173pt)=(0.547,0.7957,0.3159); rgb(175pt)=(0.5748,0.7929,0.2941); rgb(176pt)=(0.5886,0.7913,0.2833); rgb(177pt)=(0.6024,0.7896,0.2726); rgb(178pt)=(0.6161,0.7878,0.2622); rgb(179pt)=(0.6297,0.7859,0.2521); rgb(180pt)=(0.6433,0.7839,0.2423); rgb(181pt)=(0.6567,0.7818,0.2329); rgb(182pt)=(0.6701,0.7796,0.2239); rgb(183pt)=(0.6833,0.7773,0.2155); rgb(184pt)=(0.6963,0.775,0.2075); rgb(185pt)=(0.7091,0.7727,0.1998); rgb(186pt)=(0.7218,0.7703,0.1924); rgb(187pt)=(0.7344,0.7679,0.1852); rgb(188pt)=(0.7468,0.7654,0.1782); rgb(189pt)=(0.759,0.7629,0.1717); rgb(190pt)=(0.771,0.7604,0.1658); rgb(191pt)=(0.7829,0.7579,0.1608); rgb(192pt)=(0.7945,0.7554,0.157); rgb(193pt)=(0.806,0.7529,0.1546); rgb(194pt)=(0.8172,0.7505,0.1535); rgb(195pt)=(0.8281,0.7481,0.1536); rgb(196pt)=(0.8389,0.7457,0.1546); rgb(197pt)=(0.8495,0.7435,0.1564); rgb(198pt)=(0.86,0.7413,0.1587); rgb(199pt)=(0.8703,0.7392,0.1615); rgb(200pt)=(0.8804,0.7372,0.165); rgb(201pt)=(0.8903,0.7353,0.1695); rgb(202pt)=(0.9,0.7336,0.1749); rgb(203pt)=(0.9093,0.7321,0.1815); rgb(204pt)=(0.9184,0.7308,0.189); rgb(205pt)=(0.9272,0.7298,0.1973); rgb(206pt)=(0.9357,0.729,0.2061); rgb(207pt)=(0.944,0.7285,0.2151); rgb(208pt)=(0.9523,0.7284,0.2237); rgb(209pt)=(0.9606,0.7285,0.2312); rgb(210pt)=(0.9689,0.7292,0.2373); rgb(211pt)=(0.977,0.7304,0.2418); rgb(212pt)=(0.9842,0.733,0.2446); rgb(213pt)=(0.99,0.7365,0.2429); rgb(214pt)=(0.9946,0.7407,0.2394); rgb(215pt)=(0.9966,0.7458,0.2351); rgb(216pt)=(0.9971,0.7513,0.2309); rgb(217pt)=(0.9972,0.7569,0.2267); rgb(218pt)=(0.9971,0.7626,0.2224); rgb(219pt)=(0.9969,0.7683,0.2181); rgb(220pt)=(0.9966,0.774,0.2138); rgb(221pt)=(0.9962,0.7798,0.2095); rgb(222pt)=(0.9957,0.7856,0.2053); rgb(223pt)=(0.9949,0.7915,0.2012); rgb(224pt)=(0.9938,0.7974,0.1974); rgb(225pt)=(0.9923,0.8034,0.1939); rgb(226pt)=(0.9906,0.8095,0.1906); rgb(227pt)=(0.9885,0.8156,0.1875); rgb(228pt)=(0.9861,0.8218,0.1846); rgb(229pt)=(0.9835,0.828,0.1817); rgb(230pt)=(0.9807,0.8342,0.1787); rgb(231pt)=(0.9778,0.8404,0.1757); rgb(232pt)=(0.9748,0.8467,0.1726); rgb(233pt)=(0.972,0.8529,0.1695); rgb(234pt)=(0.9694,0.8591,0.1665); rgb(235pt)=(0.9671,0.8654,0.1636); rgb(236pt)=(0.9651,0.8716,0.1608); rgb(237pt)=(0.9634,0.8778,0.1582); rgb(238pt)=(0.9619,0.884,0.1557); rgb(239pt)=(0.9608,0.8902,0.1532); rgb(240pt)=(0.9601,0.8963,0.1507); rgb(241pt)=(0.9596,0.9023,0.148); rgb(242pt)=(0.9595,0.9084,0.145); rgb(243pt)=(0.9597,0.9143,0.1418); rgb(244pt)=(0.9601,0.9203,0.1382); rgb(245pt)=(0.9608,0.9262,0.1344); rgb(246pt)=(0.9618,0.932,0.1304); rgb(247pt)=(0.9629,0.9379,0.1261); rgb(248pt)=(0.9642,0.9437,0.1216); rgb(249pt)=(0.9657,0.9494,0.1168); rgb(250pt)=(0.9674,0.9552,0.1116); rgb(251pt)=(0.9692,0.9609,0.1061); rgb(252pt)=(0.9711,0.9667,0.1001); rgb(253pt)=(0.973,0.9724,0.0938); rgb(254pt)=(0.9749,0.9782,0.0872); rgb(255pt)=(0.9769,0.9839,0.0805)}, mesh/rows=37]
table[row sep=crcr, point meta=\thisrow{c}] {%
%
x	y	z	c\\
-3.40	-2.80	20639824.66	20639824.66\\
-3.40	-2.65	19958958.82	19958958.82\\
-3.40	-2.50	19714509.18	19714509.18\\
-3.40	-2.35	19890709.04	19890709.04\\
-3.40	-2.20	20498909.99	20498909.99\\
-3.40	-2.05	21578809.49	21578809.49\\
-3.40	-1.90	23202800.28	23202800.28\\
-3.40	-1.75	25484113.53	25484113.53\\
-3.40	-1.60	28590046.98	28590046.98\\
-3.40	-1.45	32762453.33	32762453.33\\
-3.40	-1.30	38349011.16	38349011.16\\
-3.40	-1.15	45850928.35	45850928.35\\
-3.40	-1.00	55996167.64	55996167.64\\
-3.40	-0.85	69852939.30	69852939.30\\
-3.40	-0.70	89007637.52	89007637.52\\
-3.40	-0.55	115847341.29	115847341.29\\
-3.40	-0.40	154014310.80	154014310.80\\
-3.40	-0.25	209147313.72	209147313.72\\
-3.40	-0.10	290108005.63	290108005.63\\
-3.40	0.05	411039297.37	411039297.37\\
-3.40	0.20	594871529.43	594871529.43\\
-3.40	0.35	879385361.77	879385361.77\\
-3.40	0.50	1327857554.05	1327857554.05\\
-3.40	0.65	2048047163.34	2048047163.34\\
-3.40	0.80	3226596266.78	3226596266.78\\
-3.40	0.95	5192368527.92	5192368527.92\\
-3.40	1.10	8534981798.63	8534981798.63\\
-3.40	1.25	14330320321.29	14330320321.29\\
-3.40	1.40	24576803971.98	24576803971.98\\
-3.40	1.55	43053763546.30	43053763546.30\\
-3.40	1.70	77039427293.49	77039427293.49\\
-3.40	1.85	140809247758.83	140809247758.83\\
-3.40	2.00	262884826449.57	262884826449.57\\
-3.25	-2.80	12456421.31	12456421.31\\
-3.25	-2.65	11944427.82	11944427.82\\
-3.25	-2.50	11699131.45	11699131.45\\
-3.25	-2.35	11704641.05	11704641.05\\
-3.25	-2.20	11961311.12	11961311.12\\
-3.25	-2.05	12485780.12	12485780.12\\
-3.25	-1.90	13312780.97	13312780.97\\
-3.25	-1.75	14499001.65	14499001.65\\
-3.25	-1.60	16129600.86	16129600.86\\
-3.25	-1.45	18328433.58	18328433.58\\
-3.25	-1.30	21273712.62	21273712.62\\
-3.25	-1.15	25221878.82	25221878.82\\
-3.25	-1.00	30544132.14	30544132.14\\
-3.25	-0.85	37782818.61	37782818.61\\
-3.25	-0.70	47739418.57	47739418.57\\
-3.25	-0.55	61613532.05	61613532.05\\
-3.25	-0.40	81225294.52	81225294.52\\
-3.25	-0.25	109376165.08	109376165.08\\
-3.25	-0.10	150442415.58	150442415.58\\
-3.25	0.05	211365505.41	211365505.41\\
-3.25	0.20	303329141.03	303329141.03\\
-3.25	0.35	444641880.92	444641880.92\\
-3.25	0.50	665767830.35	665767830.35\\
-3.25	0.65	1018243050.36	1018243050.36\\
-3.25	0.80	1590729413.11	1590729413.11\\
-3.25	0.95	2538384321.28	2538384321.28\\
-3.25	1.10	4137468023.18	4137468023.18\\
-3.25	1.25	6888555034.96	6888555034.96\\
-3.25	1.40	11714879428.53	11714879428.53\\
-3.25	1.55	20349967606.68	20349967606.68\\
-3.25	1.70	36108199212.75	36108199212.75\\
-3.25	1.85	65443142569.07	65443142569.07\\
-3.25	2.00	121154254741.51	121154254741.51\\
-3.10	-2.80	7713253.22	7713253.22\\
-3.10	-2.65	7334150.78	7334150.78\\
-3.10	-2.50	7123251.66	7123251.66\\
-3.10	-2.35	7066802.43	7066802.43\\
-3.10	-2.20	7161167.31	7161167.31\\
-3.10	-2.05	7412435.02	7412435.02\\
-3.10	-1.90	7837078.32	7837078.32\\
-3.10	-1.75	8463766.68	8463766.68\\
-3.10	-1.60	9336613.64	9336613.64\\
-3.10	-1.45	10520376.93	10520376.93\\
-3.10	-1.30	12108474.23	12108474.23\\
-3.10	-1.15	14235205.84	14235205.84\\
-3.10	-1.00	17094416.84	17094416.84\\
-3.10	-0.85	20968194.51	20968194.51\\
-3.10	-0.70	26271446.89	26271446.89\\
-3.10	-0.55	33621969.96	33621969.96\\
-3.10	-0.40	43951988.46	43951988.46\\
-3.10	-0.25	58688107.98	58688107.98\\
-3.10	-0.10	80045679.92	80045679.92\\
-3.10	0.05	111517210.04	111517210.04\\
-3.10	0.20	158694582.74	158694582.74\\
-3.10	0.35	230673930.85	230673930.85\\
-3.10	0.50	342492558.24	342492558.24\\
-3.10	0.65	519421558.58	519421558.58\\
-3.10	0.80	804646286.07	804646286.07\\
-3.10	0.95	1273228231.15	1273228231.15\\
-3.10	1.10	2057897391.22	2057897391.22\\
-3.10	1.25	3397483702.22	3397483702.22\\
-3.10	1.40	5729375096.93	5729375096.93\\
-3.10	1.55	9869004306.09	9869004306.09\\
-3.10	1.70	17364233984.76	17364233984.76\\
-3.10	1.85	31207151459.25	31207151459.25\\
-3.10	2.00	57288682756.31	57288682756.31\\
-2.95	-2.80	4900483.20	4900483.20\\
-2.95	-2.65	4620525.09	4620525.09\\
-2.95	-2.50	4449999.58	4449999.58\\
-2.95	-2.35	4377688.10	4377688.10\\
-2.95	-2.20	4398918.03	4398918.03\\
-2.95	-2.05	4515055.90	4515055.90\\
-2.95	-1.90	4733655.00	4733655.00\\
-2.95	-1.75	5069280.02	5069280.02\\
-2.95	-1.60	5545135.58	5545135.58\\
-2.95	-1.45	6195755.42	6195755.42\\
-2.95	-1.30	7071191.04	7071191.04\\
-2.95	-1.15	8243413.47	8243413.47\\
-2.95	-1.00	9816073.58	9816073.58\\
-2.95	-0.85	11939461.39	11939461.39\\
-2.95	-0.70	14833644.75	14833644.75\\
-2.95	-0.55	18824663.47	18824663.47\\
-2.95	-0.40	24401851.32	24401851.32\\
-2.95	-0.25	32309820.15	32309820.15\\
-2.95	-0.10	43698095.87	43698095.87\\
-2.95	0.05	60367986.96	60367986.96\\
-2.95	0.20	85185775.91	85185775.91\\
-2.95	0.35	122784538.28	122784538.28\\
-2.95	0.50	180774216.47	180774216.47\\
-2.95	0.65	271860119.04	271860119.04\\
-2.95	0.80	417609842.73	417609842.73\\
-2.95	0.95	655257737.46	655257737.46\\
-2.95	1.10	1050194641.98	1050194641.98\\
-2.95	1.25	1719268265.55	1719268265.55\\
-2.95	1.40	2874972561.99	2874972561.99\\
-2.95	1.55	4910661177.17	4910661177.17\\
-2.95	1.70	8567664352.10	8567664352.10\\
-2.95	1.85	15268667117.73	15268667117.73\\
-2.95	2.00	27794317741.67	27794317741.67\\
-2.80	-2.80	3194458.39	3194458.39\\
-2.80	-2.65	2986687.93	2986687.93\\
-2.80	-2.50	2852322.76	2852322.76\\
-2.80	-2.35	2782426.47	2782426.47\\
-2.80	-2.20	2772457.73	2772457.73\\
-2.80	-2.05	2821774.99	2821774.99\\
-2.80	-1.90	2933567.15	2933567.15\\
-2.80	-1.75	3115199.73	3115199.73\\
-2.80	-1.60	3379029.37	3379029.37\\
-2.80	-1.45	3743813.85	3743813.85\\
-2.80	-1.30	4236944.10	4236944.10\\
-2.80	-1.15	4897872.00	4897872.00\\
-2.80	-1.00	5783334.82	5783334.82\\
-2.80	-0.85	6975341.09	6975341.09\\
-2.80	-0.70	8593474.08	8593474.08\\
-2.80	-0.55	10814048.21	10814048.21\\
-2.80	-0.40	13900295.74	13900295.74\\
-2.80	-0.25	18250551.30	18250551.30\\
-2.80	-0.10	24476208.47	24476208.47\\
-2.80	0.05	33529610.45	33529610.45\\
-2.80	0.20	46916875.36	46916875.36\\
-2.80	0.35	67057268.94	67057268.94\\
-2.80	0.50	97899134.32	97899134.32\\
-2.80	0.65	145991661.00	145991661.00\\
-2.80	0.80	222378837.89	222378837.89\\
-2.80	0.95	345999181.56	345999181.56\\
-2.80	1.10	549886294.18	549886294.18\\
-2.80	1.25	892661752.60	892661752.60\\
-2.80	1.40	1480188984.02	1480188984.02\\
-2.80	1.55	2507053586.40	2507053586.40\\
-2.80	1.70	4337368032.54	4337368032.54\\
-2.80	1.85	7664876121.87	7664876121.87\\
-2.80	2.00	13835670630.52	13835670630.52\\
-2.65	-2.80	2136547.65	2136547.65\\
-2.65	-2.65	1980821.66	1980821.66\\
-2.65	-2.50	1875833.90	1875833.90\\
-2.65	-2.35	1814510.95	1814510.95\\
-2.65	-2.20	1792837.85	1792837.85\\
-2.65	-2.05	1809416.90	1809416.90\\
-2.65	-1.90	1865316.27	1865316.27\\
-2.65	-1.75	1964185.61	1964185.61\\
-2.65	-1.60	2112655.96	2112655.96\\
-2.65	-1.45	2321086.06	2321086.06\\
-2.65	-1.30	2604773.19	2604773.19\\
-2.65	-1.15	2985828.05	2985828.05\\
-2.65	-1.00	3496035.88	3496035.88\\
-2.65	-0.85	4181221.53	4181221.53\\
-2.65	-0.70	5107950.54	5107950.54\\
-2.65	-0.55	6373916.91	6373916.91\\
-2.65	-0.40	8124232.05	8124232.05\\
-2.65	-0.25	10577290.94	10577290.94\\
-2.65	-0.10	14066394.37	14066394.37\\
-2.65	0.05	19107652.61	19107652.61\\
-2.65	0.20	26512341.97	26512341.97\\
-2.65	0.35	37575525.02	37575525.02\\
-2.65	0.50	54397411.02	54397411.02\\
-2.65	0.65	80439175.79	80439175.79\\
-2.65	0.80	121499139.45	121499139.45\\
-2.65	0.95	187454124.44	187454124.44\\
-2.65	1.10	295415315.62	295415315.62\\
-2.65	1.25	475540185.30	475540185.30\\
-2.65	1.40	781911582.73	781911582.73\\
-2.65	1.55	1313240533.49	1313240533.49\\
-2.65	1.70	2252926999.31	2252926999.31\\
-2.65	1.85	3947900231.61	3947900231.61\\
-2.65	2.00	7066452218.73	7066452218.73\\
-2.50	-2.80	1466172.08	1466172.08\\
-2.50	-2.65	1347900.77	1347900.77\\
-2.50	-2.50	1265747.59	1265747.59\\
-2.50	-2.35	1214094.53	1214094.53\\
-2.50	-2.20	1189526.46	1189526.46\\
-2.50	-2.05	1190452.09	1190452.09\\
-2.50	-1.90	1216930.96	1216930.96\\
-2.50	-1.75	1270679.93	1270679.93\\
-2.50	-1.60	1355259.96	1355259.96\\
-2.50	-1.45	1476472.13	1476472.13\\
-2.50	-1.30	1643024.76	1643024.76\\
-2.50	-1.15	1867579.82	1867579.82\\
-2.50	-1.00	2168355.29	2168355.29\\
-2.50	-0.85	2571567.47	2571567.47\\
-2.50	-0.70	3115168.99	3115168.99\\
-2.50	-0.55	3854619.45	3854619.45\\
-2.50	-0.40	4871891.57	4871891.57\\
-2.50	-0.25	6289699.95	6289699.95\\
-2.50	-0.10	8294275.38	8294275.38\\
-2.50	0.05	11172315.24	11172315.24\\
-2.50	0.20	15371778.60	15371778.60\\
-2.50	0.35	21603358.29	21603358.29\\
-2.50	0.50	31012346.29	31012346.29\\
-2.50	0.65	45474108.07	45474108.07\\
-2.50	0.80	68109856.11	68109856.11\\
-2.50	0.95	104201013.84	104201013.84\\
-2.50	1.10	162835890.96	162835890.96\\
-2.50	1.25	259922893.77	259922893.77\\
-2.50	1.40	423794328.71	423794328.71\\
-2.50	1.55	705800491.17	705800491.17\\
-2.50	1.70	1200673651.83	1200673651.83\\
-2.50	1.85	2086335812.78	2086335812.78\\
-2.50	2.00	3703050664.45	3703050664.45\\
-2.35	-2.80	1032319.95	1032319.95\\
-2.35	-2.65	941082.03	941082.03\\
-2.35	-2.50	876308.15	876308.15\\
-2.35	-2.35	833493.91	833493.91\\
-2.35	-2.20	809774.73	809774.73\\
-2.35	-2.05	803604.23	803604.23\\
-2.35	-1.90	814585.02	814585.02\\
-2.35	-1.75	843425.70	843425.70\\
-2.35	-1.60	892017.66	892017.66\\
-2.35	-1.45	963643.24	963643.24\\
-2.35	-1.30	1063347.75	1063347.75\\
-2.35	-1.15	1198534.57	1198534.57\\
-2.35	-1.00	1379882.22	1379882.22\\
-2.35	-0.85	1622742.77	1622742.77\\
-2.35	-0.70	1949276.98	1949276.98\\
-2.35	-0.55	2391738.17	2391738.17\\
-2.35	-0.40	2997574.07	2997574.07\\
-2.35	-0.25	3837447.27	3837447.27\\
-2.35	-0.10	5018005.43	5018005.43\\
-2.35	0.05	6702487.93	6702487.93\\
-2.35	0.20	9144441.03	9144441.03\\
-2.35	0.35	12743669.35	12743669.35\\
-2.35	0.50	18140449.09	18140449.09\\
-2.35	0.65	26376537.39	26376537.39\\
-2.35	0.80	39174524.97	39174524.97\\
-2.35	0.95	59430023.77	59430023.77\\
-2.35	1.10	92092503.13	92092503.13\\
-2.35	1.25	145766880.43	145766880.43\\
-2.35	1.40	235672902.20	235672902.20\\
-2.35	1.55	389203420.97	389203420.97\\
-2.35	1.70	656537982.62	656537982.62\\
-2.35	1.85	1131251771.33	1131251771.33\\
-2.35	2.00	1991016671.79	1991016671.79\\
-2.20	-2.80	745762.79	745762.79\\
-2.20	-2.65	674146.15	674146.15\\
-2.20	-2.50	622477.43	622477.43\\
-2.20	-2.35	587096.34	587096.34\\
-2.20	-2.20	565602.53	565602.53\\
-2.20	-2.05	556582.47	556582.47\\
-2.20	-1.90	559453.39	559453.39\\
-2.20	-1.75	574400.10	574400.10\\
-2.20	-1.60	602394.93	602394.93\\
-2.20	-1.45	645303.95	645303.95\\
-2.20	-1.30	706095.66	706095.66\\
-2.20	-1.15	789185.27	789185.27\\
-2.20	-1.00	900970.58	900970.58\\
-2.20	-0.85	1050650.95	1050650.95\\
-2.20	-0.70	1251475.96	1251475.96\\
-2.20	-0.55	1522659.43	1522659.43\\
-2.20	-0.40	1892340.35	1892340.35\\
-2.20	-0.25	2402215.30	2402215.30\\
-2.20	-0.10	3114876.40	3114876.40\\
-2.20	0.05	4125588.58	4125588.58\\
-2.20	0.20	5581452.35	5581452.35\\
-2.20	0.35	7713024.95	7713024.95\\
-2.20	0.50	10887256.61	10887256.61\\
-2.20	0.65	15697424.05	15697424.05\\
-2.20	0.80	23118227.35	23118227.35\\
-2.20	0.95	34777381.44	34777381.44\\
-2.20	1.10	53438644.75	53438644.75\\
-2.20	1.25	83874549.56	83874549.56\\
-2.20	1.40	134468690.20	134468690.20\\
-2.20	1.55	220205598.82	220205598.82\\
-2.20	1.70	368342428.93	368342428.93\\
-2.20	1.85	629348710.68	629348710.68\\
-2.20	2.00	1098366186.31	1098366186.31\\
-2.05	-2.80	552769.56	552769.56\\
-2.05	-2.65	495493.14	495493.14\\
-2.05	-2.50	453677.65	453677.65\\
-2.05	-2.35	424300.29	424300.29\\
-2.05	-2.20	405336.29	405336.29\\
-2.05	-2.05	395524.92	395524.92\\
-2.05	-1.90	394228.86	394228.86\\
-2.05	-1.75	401364.72	401364.72\\
-2.05	-1.60	417393.98	417393.98\\
-2.05	-1.45	443373.14	443373.14\\
-2.05	-1.30	481070.56	481070.56\\
-2.05	-1.15	533168.39	533168.39\\
-2.05	-1.00	603581.91	603581.91\\
-2.05	-0.85	697949.90	697949.90\\
-2.05	-0.70	824381.97	824381.97\\
-2.05	-0.55	994601.10	994601.10\\
-2.05	-0.40	1225703.96	1225703.96\\
-2.05	-0.25	1542902.39	1542902.39\\
-2.05	-0.10	1983844.03	1983844.03\\
-2.05	0.05	2605510.37	2605510.37\\
-2.05	0.20	3495379.18	3495379.18\\
-2.05	0.35	4789740.59	4789740.59\\
-2.05	0.50	6704183.60	6704183.60\\
-2.05	0.65	9585086.00	9585086.00\\
-2.05	0.80	13997882.29	13997882.29\\
-2.05	0.95	20880691.11	20880691.11\\
-2.05	1.10	31815855.75	31815855.75\\
-2.05	1.25	49517482.58	49517482.58\\
-2.05	1.40	78720839.08	78720839.08\\
-2.05	1.55	127831263.08	127831263.08\\
-2.05	1.70	212031631.05	212031631.05\\
-2.05	1.85	359236483.63	359236483.63\\
-2.05	2.00	621693644.50	621693644.50\\
-1.90	-2.80	420382.42	420382.42\\
-1.90	-2.65	373661.42	373661.42\\
-1.90	-2.50	339256.50	339256.50\\
-1.90	-2.35	314625.78	314625.78\\
-1.90	-2.20	298041.43	298041.43\\
-1.90	-2.05	288386.67	288386.67\\
-1.90	-1.90	285029.58	285029.58\\
-1.90	-1.75	287753.69	287753.69\\
-1.90	-1.60	296734.52	296734.52\\
-1.90	-1.45	312558.60	312558.60\\
-1.90	-1.30	336287.74	336287.74\\
-1.90	-1.15	369578.62	369578.62\\
-1.90	-1.00	414876.50	414876.50\\
-1.90	-0.85	475715.23	475715.23\\
-1.90	-0.70	557174.81	557174.81\\
-1.90	-0.55	666579.73	666579.73\\
-1.90	-0.40	814571.00	814571.00\\
-1.90	-0.25	1016768.29	1016768.29\\
-1.90	-0.10	1296376.79	1296376.79\\
-1.90	0.05	1688327.57	1688327.57\\
-1.90	0.20	2245941.21	2245941.21\\
-1.90	0.35	3051801.38	3051801.38\\
-1.90	0.50	4235750.34	4235750.34\\
-1.90	0.65	6005105.65	6005105.65\\
-1.90	0.80	8696152.74	8696152.74\\
-1.90	0.95	12863225.23	12863225.23\\
-1.90	1.10	19435190.31	19435190.31\\
-1.90	1.25	29994659.03	29994659.03\\
-1.90	1.40	47284115.53	47284115.53\\
-1.90	1.55	76138237.40	76138237.40\\
-1.90	1.70	125229479.65	125229479.65\\
-1.90	1.85	210390701.57	210390701.57\\
-1.90	2.00	361046148.20	361046148.20\\
-1.75	-2.80	328021.25	328021.25\\
-1.75	-2.65	289118.50	289118.50\\
-1.75	-2.50	260295.09	260295.09\\
-1.75	-2.35	239371.42	239371.42\\
-1.75	-2.20	224850.99	224850.99\\
-1.75	-2.05	215741.42	215741.42\\
-1.75	-1.90	211440.65	211440.65\\
-1.75	-1.75	211670.15	211670.15\\
-1.75	-1.60	216444.71	216444.71\\
-1.75	-1.45	226073.96	226073.96\\
-1.75	-1.30	241196.12	241196.12\\
-1.75	-1.15	262848.98	262848.98\\
-1.75	-1.00	292589.33	292589.33\\
-1.75	-0.85	332680.14	332680.14\\
-1.75	-0.70	386377.19	386377.19\\
-1.75	-0.55	458365.89	458365.89\\
-1.75	-0.40	555429.98	555429.98\\
-1.75	-0.25	687483.92	687483.92\\
-1.75	-0.10	869184.52	869184.52\\
-1.75	0.05	1122477.47	1122477.47\\
-1.75	0.20	1480674.13	1480674.13\\
-1.75	0.35	1995067.31	1995067.31\\
-1.75	0.50	2745818.49	2745818.49\\
-1.75	0.65	3860133.34	3860133.34\\
-1.75	0.80	5543052.56	5543052.56\\
-1.75	0.95	8130399.67	8130399.67\\
-1.75	1.10	12181226.02	12181226.02\\
-1.75	1.25	18641734.43	18641734.43\\
-1.75	1.40	29140556.44	29140556.44\\
-1.75	1.55	46529198.60	46529198.60\\
-1.75	1.70	75887371.97	75887371.97\\
-1.75	1.85	126424039.35	126424039.35\\
-1.75	2.00	215132497.91	215132497.91\\
-1.60	-2.80	262613.13	262613.13\\
-1.60	-2.65	229525.27	229525.27\\
-1.60	-2.50	204908.89	204908.89\\
-1.60	-2.35	186856.11	186856.11\\
-1.60	-2.20	174048.38	174048.38\\
-1.60	-2.05	165595.64	165595.64\\
-1.60	-1.90	160932.60	160932.60\\
-1.60	-1.75	159755.33	159755.33\\
-1.60	-1.60	161988.02	161988.02\\
-1.60	-1.45	167774.77	167774.77\\
-1.60	-1.30	177495.19	177495.19\\
-1.60	-1.15	191806.25	191806.25\\
-1.60	-1.00	211716.71	211716.71\\
-1.60	-0.85	238706.21	238706.21\\
-1.60	-0.70	274908.71	274908.71\\
-1.60	-0.55	323392.17	323392.17\\
-1.60	-0.40	388585.60	388585.60\\
-1.60	-0.25	476936.01	476936.01\\
-1.60	-0.10	597929.17	597929.17\\
-1.60	0.05	765694.57	765694.57\\
-1.60	0.20	1001561.47	1001561.47\\
-1.60	0.35	1338184.07	1338184.07\\
-1.60	0.50	1826292.40	1826292.40\\
-1.60	0.65	2545897.86	2545897.86\\
-1.60	0.80	3625165.59	3625165.59\\
-1.60	0.95	5272674.12	5272674.12\\
-1.60	1.10	7833398.50	7833398.50\\
-1.60	1.25	11887368.32	11887368.32\\
-1.60	1.40	18426269.81	18426269.81\\
-1.60	1.55	29174628.97	29174628.97\\
-1.60	1.70	47183425.31	47183425.31\\
-1.60	1.85	77945275.78	77945275.78\\
-1.60	2.00	131524402.91	131524402.91\\
-1.45	-2.80	215718.75	215718.75\\
-1.45	-2.65	186957.19	186957.19\\
-1.45	-2.50	165505.58	165505.58\\
-1.45	-2.35	149657.78	149657.78\\
-1.45	-2.20	138229.97	138229.97\\
-1.45	-2.05	130413.12	130413.12\\
-1.45	-1.90	125677.23	125677.23\\
-1.45	-1.75	123710.94	123710.94\\
-1.45	-1.60	124387.24	124387.24\\
-1.45	-1.45	127749.66	127749.66\\
-1.45	-1.30	134017.00	134017.00\\
-1.45	-1.15	143607.21	143607.21\\
-1.45	-1.00	157184.17	157184.17\\
-1.45	-0.85	175734.72	175734.72\\
-1.45	-0.70	200688.52	200688.52\\
-1.45	-0.55	234101.24	234101.24\\
-1.45	-0.40	278933.77	278933.77\\
-1.45	-0.25	339480.40	339480.40\\
-1.45	-0.10	422031.15	422031.15\\
-1.45	0.05	535908.34	535908.34\\
-1.45	0.20	695108.73	695108.73\\
-1.45	0.35	920939.65	920939.65\\
-1.45	0.50	1246309.23	1246309.23\\
-1.45	0.65	1722807.13	1722807.13\\
-1.45	0.80	2432560.87	2432560.87\\
-1.45	0.95	3508382.87	3508382.87\\
-1.45	1.10	5168523.02	5168523.02\\
-1.45	1.25	7777537.92	7777537.92\\
-1.45	1.40	11954571.87	11954571.87\\
-1.45	1.55	18769042.86	18769042.86\\
-1.45	1.70	30099995.68	30099995.68\\
-1.45	1.85	49306814.89	49306814.89\\
-1.45	2.00	82501847.35	82501847.35\\
-1.30	-2.80	181809.42	181809.42\\
-1.30	-2.65	156246.69	156246.69\\
-1.30	-2.50	137158.11	137158.11\\
-1.30	-2.35	122983.92	122983.92\\
-1.30	-2.20	112639.68	112639.68\\
-1.30	-2.05	105378.18	105378.18\\
-1.30	-1.90	100699.24	100699.24\\
-1.30	-1.75	98291.93	98291.93\\
-1.30	-1.60	97999.93	97999.93\\
-1.30	-1.45	99804.45	99804.45\\
-1.30	-1.30	103822.20	103822.20\\
-1.30	-1.15	110318.09	110318.09\\
-1.30	-1.00	119734.55	119734.55\\
-1.30	-0.85	132742.02	132742.02\\
-1.30	-0.70	150318.90	150318.90\\
-1.30	-0.55	173874.12	173874.12\\
-1.30	-0.40	205434.09	205434.09\\
-1.30	-0.25	247928.41	247928.41\\
-1.30	-0.10	305630.22	305630.22\\
-1.30	0.05	384842.02	384842.02\\
-1.30	0.20	494976.86	494976.86\\
-1.30	0.35	650284.66	650284.66\\
-1.30	0.50	872646.46	872646.46\\
-1.30	0.65	1196160.23	1196160.23\\
-1.30	0.80	1674775.50	1674775.50\\
-1.30	0.95	2395190.43	2395190.43\\
-1.30	1.10	3498965.86	3498965.86\\
-1.30	1.25	5221022.56	5221022.56\\
-1.30	1.40	7957701.49	7957701.49\\
-1.30	1.55	12388990.67	12388990.67\\
-1.30	1.70	19701550.96	19701550.96\\
-1.30	1.85	32002294.34	32002294.34\\
-1.30	2.00	53097982.84	53097982.84\\
-1.15	-2.80	157217.86	157217.86\\
-1.15	-2.65	133978.93	133978.93\\
-1.15	-2.50	116623.84	116623.84\\
-1.15	-2.35	103694.18	103694.18\\
-1.15	-2.20	94175.44	94175.44\\
-1.15	-2.05	87364.93	87364.93\\
-1.15	-1.90	82785.22	82785.22\\
-1.15	-1.75	80128.07	80128.07\\
-1.15	-1.60	79219.62	79219.62\\
-1.15	-1.45	80001.30	80001.30\\
-1.15	-1.30	82523.48	82523.48\\
-1.15	-1.15	86950.94	86950.94\\
-1.15	-1.00	93580.89	93580.89\\
-1.15	-0.85	102876.53	102876.53\\
-1.15	-0.70	115521.19	115521.19\\
-1.15	-0.55	132502.23	132502.23\\
-1.15	-0.40	155239.04	155239.04\\
-1.15	-0.25	185778.27	185778.27\\
-1.15	-0.10	227093.70	227093.70\\
-1.15	0.05	283551.19	283551.19\\
-1.15	0.20	361638.00	361638.00\\
-1.15	0.35	471121.43	471121.43\\
-1.15	0.50	626913.85	626913.85\\
-1.15	0.65	852116.70	852116.70\\
-1.15	0.80	1183059.26	1183059.26\\
-1.15	0.95	1677761.14	1677761.14\\
-1.15	1.10	2430356.43	2430356.43\\
-1.15	1.25	3596052.46	3596052.46\\
-1.15	1.40	5434984.07	5434984.07\\
-1.15	1.55	8390478.73	8390478.73\\
-1.15	1.70	13230961.77	13230961.77\\
-1.15	1.85	21311415.53	21311415.53\\
-1.15	2.00	35063025.52	35063025.52\\
-1.00	-2.80	139490.43	139490.43\\
-1.00	-2.65	117874.32	117874.32\\
-1.00	-2.50	101744.33	101744.33\\
-1.00	-2.35	89705.17	89705.17\\
-1.00	-2.20	80786.89	80786.89\\
-1.00	-2.05	74315.69	74315.69\\
-1.00	-1.90	69829.09	69829.09\\
-1.00	-1.75	67020.62	67020.62\\
-1.00	-1.60	65704.74	65704.74\\
-1.00	-1.45	65796.26	65796.26\\
-1.00	-1.30	67301.06	67301.06\\
-1.00	-1.15	70316.75	70316.75\\
-1.00	-1.00	75043.30	75043.30\\
-1.00	-0.85	81805.26	81805.26\\
-1.00	-0.70	91089.17	91089.17\\
-1.00	-0.55	103602.08	103602.08\\
-1.00	-0.40	120361.17	120361.17\\
-1.00	-0.25	142830.37	142830.37\\
-1.00	-0.10	173129.43	173129.43\\
-1.00	0.05	214356.90	214356.90\\
-1.00	0.20	271094.21	271094.21\\
-1.00	0.35	350202.48	350202.48\\
-1.00	0.50	462098.34	462098.34\\
-1.00	0.65	622824.67	622824.67\\
-1.00	0.80	857459.10	857459.10\\
-1.00	0.95	1205805.41	1205805.41\\
-1.00	1.10	1732037.34	1732037.34\\
-1.00	1.25	2541285.64	2541285.64\\
-1.00	1.40	3808604.90	3808604.90\\
-1.00	1.55	5830349.25	5830349.25\\
-1.00	1.70	9116737.04	9116737.04\\
-1.00	1.85	14561311.54	14561311.54\\
-1.00	2.00	23756243.49	23756243.49\\
-0.85	-2.80	126982.52	126982.52\\
-0.85	-2.65	106404.24	106404.24\\
-0.85	-2.50	91073.10	91073.10\\
-0.85	-2.35	79622.82	79622.82\\
-0.85	-2.20	71105.17	71105.17\\
-0.85	-2.05	64860.60	64860.60\\
-0.85	-1.90	60433.40	60433.40\\
-0.85	-1.75	57516.08	57516.08\\
-0.85	-1.60	55913.63	55913.63\\
-0.85	-1.45	55521.65	55521.65\\
-0.85	-1.30	56314.89	56314.89\\
-0.85	-1.15	58344.56	58344.56\\
-0.85	-1.00	61743.84	61743.84\\
-0.85	-0.85	66742.60	66742.60\\
-0.85	-0.70	73693.44	73693.44\\
-0.85	-0.55	83113.34	83113.34\\
-0.85	-0.40	95747.81	95747.81\\
-0.85	-0.25	112668.68	112668.68\\
-0.85	-0.10	135423.41	135423.41\\
-0.85	0.05	166264.86	166264.86\\
-0.85	0.20	208508.34	208508.34\\
-0.85	0.35	267093.06	267093.06\\
-0.85	0.50	349476.49	349476.49\\
-0.85	0.65	467078.15	467078.15\\
-0.85	0.80	637642.63	637642.63\\
-0.85	0.95	889162.91	889162.91\\
-0.85	1.10	1266489.36	1266489.36\\
-0.85	1.25	1842629.25	1842629.25\\
-0.85	1.40	2738360.27	2738360.27\\
-0.85	1.55	4156802.75	4156802.75\\
-0.85	1.70	6445319.60	6445319.60\\
-0.85	1.85	10208118.57	10208118.57\\
-0.85	2.00	16514413.84	16514413.84\\
-0.70	-2.80	118604.31	118604.31\\
-0.70	-2.65	98549.78	98549.78\\
-0.70	-2.50	83642.51	83642.51\\
-0.70	-2.35	72512.79	72512.79\\
-0.70	-2.20	64212.33	64212.33\\
-0.70	-2.05	58081.58	58081.58\\
-0.70	-1.90	53662.97	53662.97\\
-0.70	-1.75	50643.89	50643.89\\
-0.70	-1.60	48819.77	48819.77\\
-0.70	-1.45	48070.72	48070.72\\
-0.70	-1.30	48348.35	48348.35\\
-0.70	-1.15	49670.54	49670.54\\
-0.70	-1.00	52123.36	52123.36\\
-0.70	-0.85	55870.43	55870.43\\
-0.70	-0.70	61171.33	61171.33\\
-0.70	-0.55	68411.65	68411.65\\
-0.70	-0.40	78149.88	78149.88\\
-0.70	-0.25	91189.08	91189.08\\
-0.70	-0.10	108685.99	108685.99\\
-0.70	0.05	132318.48	132318.48\\
-0.70	0.20	164544.60	164544.60\\
-0.70	0.35	209008.04	209008.04\\
-0.70	0.50	271180.56	271180.56\\
-0.70	0.65	359393.58	359393.58\\
-0.70	0.80	486517.30	486517.30\\
-0.70	0.95	672732.59	672732.59\\
-0.70	1.10	950173.38	950173.38\\
-0.70	1.25	1370816.92	1370816.92\\
-0.70	1.40	2020097.18	2020097.18\\
-0.70	1.55	3040753.94	3040753.94\\
-0.70	1.70	4675267.90	4675267.90\\
-0.70	1.85	7342567.26	7342567.26\\
-0.70	2.00	11778922.77	11778922.77\\
-0.55	-2.80	113661.67	113661.67\\
-0.55	-2.65	93650.35	93650.35\\
-0.55	-2.50	78817.19	78817.19\\
-0.55	-2.35	67756.15	67756.15\\
-0.55	-2.20	59496.68	59496.68\\
-0.55	-2.05	53364.56	53364.56\\
-0.55	-1.90	48891.05	48891.05\\
-0.55	-1.75	45753.25	45753.25\\
-0.55	-1.60	43735.16	43735.16\\
-0.55	-1.45	42702.74	42702.74\\
-0.55	-1.30	42588.96	42588.96\\
-0.55	-1.15	43386.49	43386.49\\
-0.55	-1.00	45146.92	45146.92\\
-0.55	-0.85	47986.38	47986.38\\
-0.55	-0.70	52098.36	52098.36\\
-0.55	-0.55	57775.85	57775.85\\
-0.55	-0.40	65446.26	65446.26\\
-0.55	-0.25	75725.04	75725.04\\
-0.55	-0.10	89497.40	89497.40\\
-0.55	0.05	108043.23	108043.23\\
-0.55	0.20	133229.65	133229.65\\
-0.55	0.35	167810.99	167810.99\\
-0.55	0.50	215901.74	215901.74\\
-0.55	0.65	283731.85	283731.85\\
-0.55	0.80	380869.54	380869.54\\
-0.55	0.95	522228.55	522228.55\\
-0.55	1.10	731410.49	731410.49\\
-0.55	1.25	1046352.38	1046352.38\\
-0.55	1.40	1529012.22	1529012.22\\
-0.55	1.55	2282233.91	2282233.91\\
-0.55	1.70	3479569.86	3479569.86\\
-0.55	1.85	5418850.58	5418850.58\\
-0.55	2.00	8619955.03	8619955.03\\
-0.40	-2.80	111759.54	111759.54\\
-0.40	-2.65	91310.38	91310.38\\
-0.40	-2.50	76202.97	76202.97\\
-0.40	-2.35	64959.08	64959.08\\
-0.40	-2.20	56561.91	56561.91\\
-0.40	-2.05	50306.54	50306.54\\
-0.40	-1.90	45702.61	45702.61\\
-0.40	-1.75	42410.54	42410.54\\
-0.40	-1.60	40199.70	40199.70\\
-0.40	-1.45	38921.36	38921.36\\
-0.40	-1.30	38491.91	38491.91\\
-0.40	-1.15	38883.65	38883.65\\
-0.40	-1.00	40121.85	40121.85\\
-0.40	-0.85	42287.40	42287.40\\
-0.40	-0.70	45525.76	45525.76\\
-0.40	-0.55	50063.32	50063.32\\
-0.40	-0.40	56233.92	56233.92\\
-0.40	-0.25	64519.83	64519.83\\
-0.40	-0.10	75614.37	75614.37\\
-0.40	0.05	90517.30	90517.30\\
-0.40	0.20	110681.52	110681.52\\
-0.40	0.35	138240.35	138240.35\\
-0.40	0.50	176364.33	176364.33\\
-0.40	0.65	229827.97	229827.97\\
-0.40	0.80	305922.36	305922.36\\
-0.40	0.95	415944.87	415944.87\\
-0.40	1.10	577665.62	577665.62\\
-0.40	1.25	819470.81	819470.81\\
-0.40	1.40	1187426.30	1187426.30\\
-0.40	1.55	1757502.88	1757502.88\\
-0.40	1.70	2657061.68	2657061.68\\
-0.40	1.85	4103207.07	4103207.07\\
-0.40	2.00	6472341.76	6472341.76\\
-0.25	-2.80	112748.87	112748.87\\
-0.25	-2.65	91345.66	91345.66\\
-0.25	-2.50	75592.70	75592.70\\
-0.25	-2.35	63898.11	63898.11\\
-0.25	-2.20	55171.19	55171.19\\
-0.25	-2.05	48657.85	48657.85\\
-0.25	-1.90	43833.86	43833.86\\
-0.25	-1.75	40335.06	40335.06\\
-0.25	-1.60	37911.58	37911.58\\
-0.25	-1.45	36397.98	36397.98\\
-0.25	-1.30	35694.30	35694.30\\
-0.25	-1.15	35754.99	35754.99\\
-0.25	-1.00	36583.96	36583.96\\
-0.25	-0.85	38234.99	38234.99\\
-0.25	-0.70	40817.59	40817.59\\
-0.25	-0.55	44509.22	44509.22\\
-0.25	-0.40	49575.70	49575.70\\
-0.25	-0.25	56403.22	56403.22\\
-0.25	-0.10	65547.36	65547.36\\
-0.25	0.05	77807.72	77807.72\\
-0.25	0.20	94342.28	94342.28\\
-0.25	0.35	116843.96	116843.96\\
-0.25	0.50	147816.32	147816.32\\
-0.25	0.65	191009.38	191009.38\\
-0.25	0.80	252117.63	252117.63\\
-0.25	0.95	339913.13	339913.13\\
-0.25	1.10	468111.04	468111.04\\
-0.25	1.25	658485.22	658485.22\\
-0.25	1.40	946148.68	946148.68\\
-0.25	1.55	1388637.70	1388637.70\\
-0.25	1.70	2081779.49	2081779.49\\
-0.25	1.85	3187841.41	3187841.41\\
-0.25	2.00	4986259.75	4986259.75\\
-0.10	-2.80	116706.97	116706.97\\
-0.10	-2.65	93758.95	93758.95\\
-0.10	-2.50	76938.70	76938.70\\
-0.10	-2.35	64490.12	64490.12\\
-0.10	-2.20	55215.08	55215.08\\
-0.10	-2.05	48287.91	48287.91\\
-0.10	-1.90	43135.56	43135.56\\
-0.10	-1.75	39359.41	39359.41\\
-0.10	-1.60	36684.11	36684.11\\
-0.10	-1.45	34923.96	34923.96\\
-0.10	-1.30	33961.38	33961.38\\
-0.10	-1.15	33733.65	33733.65\\
-0.10	-1.00	34226.11	34226.11\\
-0.10	-0.85	35470.55	35470.55\\
-0.10	-0.70	37548.67	37548.67\\
-0.10	-0.55	40601.06	40601.06\\
-0.10	-0.40	44843.18	44843.18\\
-0.10	-0.25	50590.81	50590.81\\
-0.10	-0.10	58299.27	58299.27\\
-0.10	0.05	68623.17	68623.17\\
-0.10	0.20	82507.73	82507.73\\
-0.10	0.35	101329.23	101329.23\\
-0.10	0.50	127113.31	127113.31\\
-0.10	0.65	162878.41	162878.41\\
-0.10	0.80	213182.83	213182.83\\
-0.10	0.95	285008.05	285008.05\\
-0.10	1.10	389204.86	389204.86\\
-0.10	1.25	542894.68	542894.68\\
-0.10	1.40	773515.69	773515.69\\
-0.10	1.55	1125742.01	1125742.01\\
-0.10	1.70	1673496.59	1673496.59\\
-0.10	1.85	2541130.66	2541130.66\\
-0.10	2.00	3941353.33	3941353.33\\
0.05	-2.80	123947.69	123947.69\\
0.05	-2.65	98740.32	98740.32\\
0.05	-2.50	80346.47	80346.47\\
0.05	-2.35	66781.37	66781.37\\
0.05	-2.20	56697.00	56697.00\\
0.05	-2.05	49167.82	49167.82\\
0.05	-1.90	43553.00	43553.00\\
0.05	-1.75	39406.83	39406.83\\
0.05	-1.60	36420.09	36420.09\\
0.05	-1.45	34381.66	34381.66\\
0.05	-1.30	33153.45	33153.45\\
0.05	-1.15	32654.79	32654.79\\
0.05	-1.00	32853.48	32853.48\\
0.05	-0.85	33762.29	33762.29\\
0.05	-0.70	35440.41	35440.41\\
0.05	-0.55	37999.84	37999.84\\
0.05	-0.40	41617.98	41617.98\\
0.05	-0.25	46558.22	46558.22\\
0.05	-0.10	53202.01	53202.01\\
0.05	0.05	62097.75	62097.75\\
0.05	0.20	74035.49	74035.49\\
0.05	0.35	90161.31	90161.31\\
0.05	0.50	112154.49	112154.49\\
0.05	0.65	142504.75	142504.75\\
0.05	0.80	184951.66	184951.66\\
0.05	0.95	245190.33	245190.33\\
0.05	1.10	332020.30	332020.30\\
0.05	1.25	459242.59	459242.59\\
0.05	1.40	648837.41	648837.41\\
0.05	1.55	936366.38	936366.38\\
0.05	1.70	1380295.12	1380295.12\\
0.05	1.85	2078329.01	2078329.01\\
0.05	2.00	3196486.37	3196486.37\\
0.20	-2.80	135063.21	135063.21\\
0.20	-2.65	106692.37	106692.37\\
0.20	-2.50	86088.64	86088.64\\
0.20	-2.35	70953.61	70953.61\\
0.20	-2.20	59733.70	59733.70\\
0.20	-2.05	51366.57	51366.57\\
0.20	-1.90	45118.83	45118.83\\
0.20	-1.75	40481.01	40481.01\\
0.20	-1.60	37098.91	37098.91\\
0.20	-1.45	34728.58	34728.58\\
0.20	-1.30	33206.97	33206.97\\
0.20	-1.15	32433.03	32433.03\\
0.20	-1.00	32356.55	32356.55\\
0.20	-0.85	32972.58	32972.58\\
0.20	-0.70	34321.00	34321.00\\
0.20	-0.55	36490.78	36490.78\\
0.20	-0.40	39629.86	39629.86\\
0.20	-0.25	43962.07	43962.07\\
0.20	-0.10	49813.83	49813.83\\
0.20	0.05	57655.14	57655.14\\
0.20	0.20	68161.99	68161.99\\
0.20	0.35	82311.91	82311.91\\
0.20	0.50	101531.16	101531.16\\
0.20	0.65	127924.05	127924.05\\
0.20	0.80	164634.66	164634.66\\
0.20	0.95	216424.57	216424.57\\
0.20	1.10	290608.32	290608.32\\
0.20	1.25	398589.40	398589.40\\
0.20	1.40	558418.31	558418.31\\
0.20	1.55	799115.89	799115.89\\
0.20	1.70	1168089.41	1168089.41\\
0.20	1.85	1744048.75	1744048.75\\
0.20	2.00	2659851.31	2659851.31\\
0.35	-2.80	151005.48	151005.48\\
0.35	-2.65	118284.87	118284.87\\
0.35	-2.50	94641.55	94641.55\\
0.35	-2.35	77348.29	77348.29\\
0.35	-2.20	64570.75	64570.75\\
0.35	-2.05	55060.11	55060.11\\
0.35	-1.90	47957.29	47957.29\\
0.35	-1.75	42666.63	42666.63\\
0.35	-1.60	38773.79	38773.79\\
0.35	-1.45	35991.87	35991.87\\
0.35	-1.30	34126.10	34126.10\\
0.35	-1.15	33051.05	33051.05\\
0.35	-1.00	32696.41	32696.41\\
0.35	-0.85	33039.31	33039.31\\
0.35	-0.70	34101.87	34101.87\\
0.35	-0.55	35953.53	35953.53\\
0.35	-0.40	38718.73	38718.73\\
0.35	-0.25	42590.91	42590.91\\
0.35	-0.10	47855.18	47855.18\\
0.35	0.05	54923.36	54923.36\\
0.35	0.20	64387.50	64387.50\\
0.35	0.35	77101.39	77101.39\\
0.35	0.50	94305.93	94305.93\\
0.35	0.65	117823.54	117823.54\\
0.35	0.80	150363.12	150363.12\\
0.35	0.95	196004.84	196004.84\\
0.35	1.10	260980.74	260980.74\\
0.35	1.25	354949.32	354949.32\\
0.35	1.40	493106.17	493106.17\\
0.35	1.55	699730.38	699730.38\\
0.35	1.70	1014231.84	1014231.84\\
0.35	1.85	1501619.69	1501619.69\\
0.35	2.00	2270904.54	2270904.54\\
0.50	-2.80	173222.93	173222.93\\
0.50	-2.65	134549.48	134549.48\\
0.50	-2.50	106751.72	106751.72\\
0.50	-2.35	86513.51	86513.51\\
0.50	-2.20	71615.86	71615.86\\
0.50	-2.05	60555.10	60555.10\\
0.50	-1.90	52300.81	52300.81\\
0.50	-1.75	46140.50	46140.50\\
0.50	-1.60	41578.84	41578.84\\
0.50	-1.45	38271.78	38271.78\\
0.50	-1.30	35983.32	35983.32\\
0.50	-1.15	34557.31	34557.31\\
0.50	-1.00	33899.62	33899.62\\
0.50	-0.85	33967.69	33967.69\\
0.50	-0.70	34765.89	34765.89\\
0.50	-0.55	36346.03	36346.03\\
0.50	-0.40	38812.96	38812.96\\
0.50	-0.25	42336.28	42336.28\\
0.50	-0.10	47169.89	47169.89\\
0.50	0.05	53682.57	53682.57\\
0.50	0.20	62404.78	62404.78\\
0.50	0.35	74100.08	74100.08\\
0.50	0.50	89874.34	89874.34\\
0.50	0.65	111344.55	111344.55\\
0.50	0.80	140902.40	140902.40\\
0.50	0.95	182131.07	182131.07\\
0.50	1.10	240472.76	240472.76\\
0.50	1.25	324312.70	324312.70\\
0.50	1.40	446764.06	446764.06\\
0.50	1.55	628649.70	628649.70\\
0.50	1.70	903556.69	903556.69\\
0.50	1.85	1326533.75	1326533.75\\
0.50	2.00	1989286.85	1989286.85\\
0.65	-2.80	203880.21	203880.21\\
0.65	-2.65	157033.35	157033.35\\
0.65	-2.50	123544.93	123544.93\\
0.65	-2.35	99282.84	99282.84\\
0.65	-2.20	81496.63	81496.63\\
0.65	-2.05	68331.55	68331.55\\
0.65	-1.90	58522.00	58522.00\\
0.65	-1.75	51195.67	51195.67\\
0.65	-1.60	45747.10	45747.10\\
0.65	-1.45	41755.15	41755.15\\
0.65	-1.30	38928.95	38928.95\\
0.65	-1.15	37072.48	37072.48\\
0.65	-1.00	36061.75	36061.75\\
0.65	-0.85	35830.93	35830.93\\
0.65	-0.70	36365.17	36365.17\\
0.65	-0.55	37698.96	37698.96\\
0.65	-0.40	39919.89	39919.89\\
0.65	-0.25	43178.30	43178.30\\
0.65	-0.10	47704.34	47704.34\\
0.65	0.05	53835.21	53835.21\\
0.65	0.20	62057.07	62057.07\\
0.65	0.35	73068.84	73068.84\\
0.65	0.50	87879.88	87879.88\\
0.65	0.65	107959.99	107959.99\\
0.65	0.80	135472.91	135472.91\\
0.65	0.95	173643.41	173643.41\\
0.65	1.10	227342.35	227342.35\\
0.65	1.25	304031.51	304031.51\\
0.65	1.40	415310.63	415310.63\\
0.65	1.55	579487.02	579487.02\\
0.65	1.70	825905.93	825905.93\\
0.65	1.85	1202357.64	1202357.64\\
0.65	2.00	1787940.08	1787940.08\\
0.80	-2.80	246207.78	246207.78\\
0.80	-2.65	188043.70	188043.70\\
0.80	-2.50	146700.64	146700.64\\
0.80	-2.35	116901.86	116901.86\\
0.80	-2.20	95154.00	95154.00\\
0.80	-2.05	79113.19	79113.19\\
0.80	-1.90	67187.26	67187.26\\
0.80	-1.75	58282.91	58282.91\\
0.80	-1.60	51643.03	51643.03\\
0.80	-1.45	46741.04	46741.04\\
0.80	-1.30	43211.69	43211.69\\
0.80	-1.15	40805.66	40805.66\\
0.80	-1.00	39360.05	39360.05\\
0.80	-0.85	38779.94	38779.94\\
0.80	-0.70	39027.87	39027.87\\
0.80	-0.55	40119.81	40119.81\\
0.80	-0.40	42126.85	42126.85\\
0.80	-0.25	45183.03	45183.03\\
0.80	-0.10	49500.30	49500.30\\
0.80	0.05	55393.22	55393.22\\
0.80	0.20	63317.18	63317.18\\
0.80	0.35	73926.95	73926.95\\
0.80	0.50	88165.80	88165.80\\
0.80	0.65	107402.35	107402.35\\
0.80	0.80	133642.19	133642.19\\
0.80	0.95	169859.40	169859.40\\
0.80	1.10	220521.94	220521.94\\
0.80	1.25	292435.60	292435.60\\
0.80	1.40	396118.27	396118.27\\
0.80	1.55	548069.62	548069.62\\
0.80	1.70	774573.77	774573.77\\
0.80	1.85	1118165.40	1118165.40\\
0.80	2.00	1648790.68	1648790.68\\
0.95	-2.80	305060.15	305060.15\\
0.95	-2.65	231037.61	231037.61\\
0.95	-2.50	178729.43	178729.43\\
0.95	-2.35	141229.59	141229.59\\
0.95	-2.20	113991.25	113991.25\\
0.95	-2.05	93979.59	93979.59\\
0.95	-1.90	79142.87	79142.87\\
0.95	-1.75	68077.91	68077.91\\
0.95	-1.60	59815.93	59815.93\\
0.95	-1.45	53683.86	53683.86\\
0.95	-1.30	49213.79	49213.79\\
0.95	-1.15	46083.57	46083.57\\
0.95	-1.00	44077.97	44077.97\\
0.95	-0.85	43063.89	43063.89\\
0.95	-0.70	42975.52	42975.52\\
0.95	-0.55	43807.18	43807.18\\
0.95	-0.40	45612.68	45612.68\\
0.95	-0.25	48511.21	48511.21\\
0.95	-0.10	52700.52	52700.52\\
0.95	0.05	58479.52	58479.52\\
0.95	0.20	66284.04	66284.04\\
0.95	0.35	76741.51	76741.51\\
0.95	0.50	90754.45	90754.45\\
0.95	0.65	109628.05	109628.05\\
0.95	0.80	135266.95	135266.95\\
0.95	0.95	170481.74	170481.74\\
0.95	1.10	219472.59	219472.59\\
0.95	1.25	288601.71	288601.71\\
0.95	1.40	387644.59	387644.59\\
0.95	1.55	531844.60	531844.60\\
0.95	1.70	745335.82	745335.82\\
0.95	1.85	1066928.82	1066928.82\\
0.95	2.00	1560037.73	1560037.73\\
1.10	-2.80	387816.43	387816.43\\
1.10	-2.65	291248.43	291248.43\\
1.10	-2.50	223417.49	223417.49\\
1.10	-2.35	175060.02	175060.02\\
1.10	-2.20	140111.24	140111.24\\
1.10	-2.05	114544.75	114544.75\\
1.10	-1.90	95651.90	95651.90\\
1.10	-1.75	81588.36	81588.36\\
1.10	-1.60	71085.18	71085.18\\
1.10	-1.45	63262.47	63262.47\\
1.10	-1.30	57508.15	57508.15\\
1.10	-1.15	53398.47	53398.47\\
1.10	-1.00	50645.92	50645.92\\
1.10	-0.85	49065.52	49065.52\\
1.10	-0.70	48553.94	48553.94\\
1.10	-0.55	49078.22	49078.22\\
1.10	-0.40	50672.14	50672.14\\
1.10	-0.25	53439.94	53439.94\\
1.10	-0.10	57567.70	57567.70\\
1.10	0.05	63344.37	63344.37\\
1.10	0.20	71195.64	71195.64\\
1.10	0.35	81736.29	81736.29\\
1.10	0.50	95850.12	95850.12\\
1.10	0.65	114811.83	114811.83\\
1.10	0.80	140474.28	140474.28\\
1.10	0.95	175559.04	175559.04\\
1.10	1.10	224112.35	224112.35\\
1.10	1.25	292229.85	292229.85\\
1.10	1.40	389223.98	389223.98\\
1.10	1.55	529530.28	529530.28\\
1.10	1.70	735865.13	735865.13\\
1.10	1.85	1044532.26	1044532.26\\
1.10	2.00	1514473.56	1514473.56\\
1.25	-2.80	505852.56	505852.56\\
1.25	-2.65	376705.11	376705.11\\
1.25	-2.50	286546.62	286546.62\\
1.25	-2.35	222641.06	222641.06\\
1.25	-2.20	176697.92	176697.92\\
1.25	-2.05	143243.14	143243.14\\
1.25	-1.90	118613.04	118613.04\\
1.25	-1.75	100324.56	100324.56\\
1.25	-1.60	86675.88	86675.88\\
1.25	-1.45	76490.15	76490.15\\
1.25	-1.30	68949.15	68949.15\\
1.25	-1.15	63484.63	63484.63\\
1.25	-1.00	59706.89	59706.89\\
1.25	-0.85	57358.33	57358.33\\
1.25	-0.70	56283.98	56283.98\\
1.25	-0.55	56414.31	56414.31\\
1.25	-0.40	57757.71	57757.71\\
1.25	-0.25	60401.38	60401.38\\
1.25	-0.10	64520.83	64520.83\\
1.25	0.05	70399.45	70399.45\\
1.25	0.20	78461.17	78461.17\\
1.25	0.35	89321.60	89321.60\\
1.25	0.50	103866.25	103866.25\\
1.25	0.65	123369.73	123369.73\\
1.25	0.80	149678.34	149678.34\\
1.25	0.95	185492.14	185492.14\\
1.25	1.10	234805.51	234805.51\\
1.25	1.25	303603.85	303603.85\\
1.25	1.40	400979.77	400979.77\\
1.25	1.55	540945.93	540945.93\\
1.25	1.70	745420.73	745420.73\\
1.25	1.85	1049216.91	1049216.91\\
1.25	2.00	1508499.97	1508499.97\\
1.40	-2.80	676984.48	676984.48\\
1.40	-2.65	499915.33	499915.33\\
1.40	-2.50	377077.32	377077.32\\
1.40	-2.35	290523.03	290523.03\\
1.40	-2.20	228637.22	228637.22\\
1.40	-2.05	183793.23	183793.23\\
1.40	-1.90	150913.57	150913.57\\
1.40	-1.75	126573.65	126573.65\\
1.40	-1.60	108436.25	108436.25\\
1.40	-1.45	94890.32	94890.32\\
1.40	-1.30	84817.51	84817.51\\
1.40	-1.15	77440.00	77440.00\\
1.40	-1.00	72220.65	72220.65\\
1.40	-0.85	68797.66	68797.66\\
1.40	-0.70	66942.53	66942.53\\
1.40	-0.55	66534.48	66534.48\\
1.40	-0.40	67547.24	67547.24\\
1.40	-0.25	70046.22	70046.22\\
1.40	-0.10	74195.57	74195.57\\
1.40	0.05	80276.33	80276.33\\
1.40	0.20	88718.30	88718.30\\
1.40	0.35	100150.96	100150.96\\
1.40	0.50	115481.72	115481.72\\
1.40	0.65	136015.24	136015.24\\
1.40	0.80	163635.72	163635.72\\
1.40	0.95	201087.40	201087.40\\
1.40	1.10	252410.73	252410.73\\
1.40	1.25	323628.66	323628.66\\
1.40	1.40	423840.40	423840.40\\
1.40	1.55	566988.09	566988.09\\
1.40	1.70	774750.25	774750.25\\
1.40	1.85	1081348.60	1081348.60\\
1.40	2.00	1541650.51	1541650.51\\
1.55	-2.80	929587.95	929587.95\\
1.55	-2.65	680688.50	680688.50\\
1.55	-2.50	509122.80	509122.80\\
1.55	-2.35	388967.12	388967.12\\
1.55	-2.20	303542.46	303542.46\\
1.55	-2.05	241959.24	241959.24\\
1.55	-1.90	197006.78	197006.78\\
1.55	-1.75	163846.19	163846.19\\
1.55	-1.60	139189.91	139189.91\\
1.55	-1.45	120780.09	120780.09\\
1.55	-1.30	107053.07	107053.07\\
1.55	-1.15	96921.28	96921.28\\
1.55	-1.00	89630.41	89630.41\\
1.55	-0.85	84665.76	84665.76\\
1.55	-0.70	81691.42	81691.42\\
1.55	-0.55	80512.12	80512.12\\
1.55	-0.40	81051.74	81051.74\\
1.55	-0.25	83345.01	83345.01\\
1.55	-0.10	87541.32	87541.32\\
1.55	0.05	93921.01	93921.01\\
1.55	0.20	102926.84	102926.84\\
1.55	0.35	115215.46	115215.46\\
1.55	0.50	131737.39	131737.39\\
1.55	0.65	153859.23	153859.23\\
1.55	0.80	183549.96	183549.96\\
1.55	0.95	223666.64	223666.64\\
1.55	1.10	278396.88	278396.88\\
1.55	1.25	353951.46	353951.46\\
1.55	1.40	459662.70	459662.70\\
1.55	1.55	609748.90	609748.90\\
1.55	1.70	826188.22	826188.22\\
1.55	1.85	1143465.80	1143465.80\\
1.55	2.00	1616529.26	1616529.26\\
1.70	-2.80	1309662.15	1309662.15\\
1.70	-2.65	950949.35	950949.35\\
1.70	-2.50	705296.46	705296.46\\
1.70	-2.35	534320.99	534320.99\\
1.70	-2.20	413474.73	413474.73\\
1.70	-2.05	326822.47	326822.47\\
1.70	-1.90	263870.64	263870.64\\
1.70	-1.75	217613.80	217613.80\\
1.70	-1.60	183315.00	183315.00\\
1.70	-1.45	157734.17	157734.17\\
1.70	-1.30	138634.00	138634.00\\
1.70	-1.15	124460.05	124460.05\\
1.70	-1.00	114131.72	114131.72\\
1.70	-0.85	106905.24	106905.24\\
1.70	-0.70	102284.02	102284.02\\
1.70	-0.55	99961.51	99961.51\\
1.70	-0.40	99787.02	99787.02\\
1.70	-0.25	101749.32	101749.32\\
1.70	-0.10	105975.43	105975.43\\
1.70	0.05	112744.41	112744.41\\
1.70	0.20	122518.34	122518.34\\
1.70	0.35	135995.14	135995.14\\
1.70	0.50	154192.01	154192.01\\
1.70	0.65	178573.32	178573.32\\
1.70	0.80	211245.50	211245.50\\
1.70	0.95	255255.19	255255.19\\
1.70	1.10	315048.86	315048.86\\
1.70	1.25	397189.21	397189.21\\
1.70	1.40	511485.33	511485.33\\
1.70	1.55	672798.68	672798.68\\
1.70	1.70	903968.48	903968.48\\
1.70	1.85	1240616.78	1240616.78\\
1.70	2.00	1739154.75	1739154.75\\
1.85	-2.80	1893150.24	1893150.24\\
1.85	-2.65	1363086.37	1363086.37\\
1.85	-2.50	1002484.99	1002484.99\\
1.85	-2.35	753092.96	753092.96\\
1.85	-2.20	577877.13	577877.13\\
1.85	-2.05	452937.87	452937.87\\
1.85	-1.90	362625.19	362625.19\\
1.85	-1.75	296546.99	296546.99\\
1.85	-1.60	247710.99	247710.99\\
1.85	-1.45	211355.34	211355.34\\
1.85	-1.30	184203.30	184203.30\\
1.85	-1.15	163982.61	163982.61\\
1.85	-1.00	149112.61	149112.61\\
1.85	-0.85	138499.17	138499.17\\
1.85	-0.70	131400.24	131400.24\\
1.85	-0.55	127338.98	127338.98\\
1.85	-0.40	126049.99	126049.99\\
1.85	-0.25	127450.17	127450.17\\
1.85	-0.10	131629.82	131629.82\\
1.85	0.05	138862.29	138862.29\\
1.85	0.20	149634.10	149634.10\\
1.85	0.35	164699.79	164699.79\\
1.85	0.50	185170.46	185170.46\\
1.85	0.65	212650.59	212650.59\\
1.85	0.80	249446.65	249446.65\\
1.85	0.95	298885.59	298885.59\\
1.85	1.10	365804.03	365804.03\\
1.85	1.25	457307.37	457307.37\\
1.85	1.40	583961.37	583961.37\\
1.85	1.55	761686.49	761686.49\\
1.85	1.70	1014809.66	1014809.66\\
1.85	1.85	1381049.17	1381049.17\\
1.85	2.00	1919773.08	1919773.08\\
2.00	-2.80	2807811.19	2807811.19\\
2.00	-2.65	2004686.10	2004686.10\\
2.00	-2.50	1461978.76	1461978.76\\
2.00	-2.35	1089060.38	1089060.38\\
2.00	-2.20	828665.12	828665.12\\
2.00	-2.05	644054.13	644054.13\\
2.00	-1.90	511307.18	511307.18\\
2.00	-1.75	414627.04	414627.04\\
2.00	-1.60	343438.96	343438.96\\
2.00	-1.45	290574.64	290574.64\\
2.00	-1.30	251120.44	251120.44\\
2.00	-1.15	221678.01	221678.01\\
2.00	-1.00	199884.62	199884.62\\
2.00	-0.85	184099.40	184099.40\\
2.00	-0.70	173197.47	173197.47\\
2.00	-0.55	166435.88	166435.88\\
2.00	-0.40	163368.59	163368.59\\
2.00	-0.25	163797.16	163797.16\\
2.00	-0.10	167749.17	167749.17\\
2.00	0.05	175481.21	175481.21\\
2.00	0.20	187506.82	187506.82\\
2.00	0.35	204653.75	204653.75\\
2.00	0.50	228159.50	228159.50\\
2.00	0.65	259820.63	259820.63\\
2.00	0.80	302221.18	302221.18\\
2.00	0.95	359080.96	359080.96\\
2.00	1.10	435788.80	435788.80\\
2.00	1.25	540226.60	540226.60\\
2.00	1.40	684056.66	684056.66\\
2.00	1.55	884757.77	884757.77\\
2.00	1.70	1168888.03	1168888.03\\
2.00	1.85	1577384.77	1577384.77\\
2.00	2.00	2174295.58	2174295.58\\
};
\addplot3 [color=white!60!black]
 table[row sep=crcr] {%
1.00	1.00	0.00\\
1.00	1.00	314144.00\\
};
 \addplot3 [color=white!60!black]
 table[row sep=crcr] {%
-1.00	1.00	0.00\\
-1.00	1.00	625776.00\\
};
 \addplot3 [color=white!60!black]
 table[row sep=crcr] {%
-1.00	-1.00	0.00\\
-1.00	-1.00	4760.00\\
};
 \addplot3 [color=white!60!black]
 table[row sep=crcr] {%
-1.41	0.00	0.00\\
-1.41	0.00	800000.00\\
};
 \addplot3 [color=white!60!black]
 table[row sep=crcr] {%
0.00	-1.41	0.00\\
0.00	-1.41	32072.00\\
};
 \addplot3 [color=white!60!black]
 table[row sep=crcr] {%
0.00	0.00	0.00\\
0.00	0.00	43472.00\\
};
 \addplot3 [color=white!60!black]
 table[row sep=crcr] {%
0.00	0.00	0.00\\
0.00	0.00	40144.00\\
};
 \addplot3 [color=white!60!black]
 table[row sep=crcr] {%
0.00	0.00	0.00\\
0.00	0.00	70440.00\\
};
 \addplot3 [color=white!60!black]
 table[row sep=crcr] {%
0.00	0.00	0.00\\
0.00	0.00	45296.00\\
};
 \addplot3 [color=white!60!black]
 table[row sep=crcr] {%
0.00	0.00	0.00\\
0.00	0.00	44560.00\\
};
 \addplot3 [color=white!60!black]
 table[row sep=crcr] {%
1.00	1.00	0.00\\
1.00	1.00	210584.00\\
};
 \addplot3 [color=white!60!black]
 table[row sep=crcr] {%
-1.00	1.00	0.00\\
-1.00	1.00	752936.00\\
};
 \addplot3 [color=white!60!black]
 table[row sep=crcr] {%
-1.00	-1.00	0.00\\
-1.00	-1.00	5264.00\\
};
 \addplot3 [color=white!60!black]
 table[row sep=crcr] {%
-1.41	0.00	0.00\\
-1.41	0.00	709336.00\\
};
 \addplot3 [color=white!60!black]
 table[row sep=crcr] {%
0.00	-1.41	0.00\\
0.00	-1.41	42728.00\\
};
 \addplot3 [color=white!60!black]
 table[row sep=crcr] {%
0.00	0.00	0.00\\
0.00	0.00	20936.00\\
};
 \addplot3 [color=white!60!black]
 table[row sep=crcr] {%
0.00	0.00	0.00\\
0.00	0.00	71744.00\\
};
 \addplot3 [color=white!60!black]
 table[row sep=crcr] {%
0.00	0.00	0.00\\
0.00	0.00	57960.00\\
};
 \addplot3 [color=white!60!black]
 table[row sep=crcr] {%
0.00	0.00	0.00\\
0.00	0.00	22432.00\\
};
 \addplot3 [color=white!60!black]
 table[row sep=crcr] {%
0.00	0.00	0.00\\
0.00	0.00	65816.00\\
};
 \addplot3 [color=white!30!black]
 table[row sep=crcr] {%
-2.42	-2.17	0.00\\
-2.42	-2.17	1850000.00\\
};
 \addplot3 [color=white!30!black]
 table[row sep=crcr] {%
-3.15	-2.56	0.00\\
-3.15	-2.56	4200000.00\\
};
 \addplot3 [color=white!30!black]
 table[row sep=crcr] {%
-3.15	-2.56	0.00\\
-3.15	-2.56	3950000.00\\
};
 \addplot3 [color=white!30!black]
 table[row sep=crcr] {%
-2.42	-2.17	0.00\\
-2.42	-2.17	1910000.00\\
};
 \addplot3 [color=white!30!black]
 table[row sep=crcr] {%
-3.15	-2.56	0.00\\
-3.15	-2.56	4120000.00\\
};
 \addplot3 [color=white!30!black]
 table[row sep=crcr] {%
-3.15	-2.56	0.00\\
-3.15	-2.56	4300000.00\\
};
 \addplot3 [color=mycolor1, only marks, mark size=2.5pt, mark=*, mark options={solid, fill=white, draw=white!40!black}]
 table[row sep=crcr] {%
1.00	1.00	0.00\\
-1.00	1.00	0.00\\
-1.00	-1.00	0.00\\
-1.41	0.00	0.00\\
0.00	-1.41	0.00\\
0.00	0.00	0.00\\
0.00	0.00	0.00\\
0.00	0.00	0.00\\
0.00	0.00	0.00\\
0.00	0.00	0.00\\
1.00	1.00	0.00\\
-1.00	1.00	0.00\\
-1.00	-1.00	0.00\\
-1.41	0.00	0.00\\
0.00	-1.41	0.00\\
0.00	0.00	0.00\\
0.00	0.00	0.00\\
0.00	0.00	0.00\\
0.00	0.00	0.00\\
0.00	0.00	0.00\\
};
 \addlegendentry{Basis-CCD}

\addplot3 [color=mycolor2, only marks, mark size=4.3pt, mark=diamond*, mark options={solid, fill=gray, draw=white!10!black}]
 table[row sep=crcr] {%
-2.42	-2.17	0.00\\
-3.15	-2.56	0.00\\
-3.15	-2.56	0.00\\
-2.42	-2.17	0.00\\
-3.15	-2.56	0.00\\
-3.15	-2.56	0.00\\
};
 \addlegendentry{Ankerpunkte}

\addplot3 [color=black, line width=1.2pt, only marks, mark=x, mark options={solid, black}]
 table[row sep=crcr] {%
1.00	1.00	314144.00\\
-1.00	1.00	625776.00\\
-1.00	-1.00	4760.00\\
-1.41	0.00	800000.00\\
0.00	-1.41	32072.00\\
0.00	0.00	43472.00\\
0.00	0.00	40144.00\\
0.00	0.00	70440.00\\
0.00	0.00	45296.00\\
0.00	0.00	44560.00\\
1.00	1.00	210584.00\\
-1.00	1.00	752936.00\\
-1.00	-1.00	5264.00\\
-1.41	0.00	709336.00\\
0.00	-1.41	42728.00\\
0.00	0.00	20936.00\\
0.00	0.00	71744.00\\
0.00	0.00	57960.00\\
0.00	0.00	22432.00\\
0.00	0.00	65816.00\\
-2.42	-2.17	1850000.00\\
-3.15	-2.56	4200000.00\\
-3.15	-2.56	3950000.00\\
-2.42	-2.17	1910000.00\\
-3.15	-2.56	4120000.00\\
-3.15	-2.56	4300000.00\\
};
 \addlegendentry{Versuchsergebnisse}

\addplot3 [color=white!20!black, dotted, line width=1.0pt]
 table[row sep=crcr] {%
-1.50	-1.80	0.00\\
-1.50	-1.80	134597.30\\
};
 \addplot3 [color=mycolor3, line width=1.2pt, only marks, mark size=3.5pt, mark=*, mark options={solid, fill=white, draw=white!20!black}]
 table[row sep=crcr] {%
-1.50	-1.80	0.00\\
};
 \addplot3 [color=mycolor4, line width=1.5pt, only marks, mark size=3.5pt, mark=*, mark options={solid, fill=white, draw=black}]
 table[row sep=crcr] {%
-1.50	-1.80	134597.30\\
};
 \addlegendentry{Prognosepunkt}

\addplot3 [color=white!40!black, dotted, line width=1.5pt]
 table[row sep=crcr] {%
nan	nan\\
};
 \addlegendentry{Antwort (Basis-CCD)}

\addplot3 [color=white!15!black, line width=1.5pt]
 table[row sep=crcr] {%
nan	nan\\
};
 \addlegendentry{Antwort (EmALT)}

\end{axis}
\end{tikzpicture}%
    \caption{Modellanpassung im Parameterraum: Das Basis-\ac{CCD} (links) extrapoliert starr aus dem Hochlastzentrum. Das \ac{EmALT}-Design (rechts) wird durch die \acp{AP} im südwestlichen Quadranten gestützt und passt die Response Surface der wahren physikalischen Degradation im Niedrigstressbereich an.}
    \label{fig:ma_abb6.1_response_surface}
\end{figure}

Die 3D-Darstellung der Regressionsflächen liefert eine klare ingenieurswissenschaftliche Interpretation: Das Basis-\ac{CCD} approximiert die Versuchspunkte im beschleunigten Zentrum exzellent, extrapoliert diesen steilen Gradienten jedoch ungebremst in den fernen \textit{Use-Space}.
Das \ac{EmALT}-Design hingegen registriert durch seine \acp{AP}, dass die Lebensdauer bei moderateren Lasten physikalisch nicht in dem Maße ansteigt, wie es die rein lineare Hochlast-Extrapolation vermuten ließe.
Die Response Surface erfährt durch die Ankerpunkte eine merkliche Dämpfung (Krümmung) in Richtung geringerer Lebensdauern.

\section{Konsequenzen für die Zuverlässigkeit und Falsifikation} \label{sec:case_falsifikation}

Diese Korrektur der Antwortfläche wirkt sich unmittelbar auf die statistische Feldabsicherung am anvisierten Prognosepunkt $(-1{,}80 ; -1{,}60)$ aus.
Abbildung~\ref{fig:ma_abb6.1_reliability_ci} transferiert das jeweilige Modell auf den resultierenden Verlauf der Überlebenswahrscheinlichkeit $\sym{R}(\sym{t})$.
Dargestellt sind die Zuverlässigkeitsfunktionen auf Grundlage der Gleichungen~\ref{eq:reldef}, \ref{eq:glm_link} und \ref{eq:var_T_general} sowie die korrespondierenden Bänder der $90\,\%$-\acp{CI}, resultierend aus der Fisher-Informationsmatrix (\ac{FIM}).

\begin{figure}[htbp]
    \centering
    \setlength\figurewidth{0.9\textwidth}
    \setlength\figureheight{6.5cm}
    % This file was created by matlab2tikz.
%
\definecolor{mycolor1}{rgb}{0.12941,0.12941,0.12941}%
%
\begin{tikzpicture}

\begin{axis}[%
width=\figurewidth,
height=0.94\figureheight,
at={(0\figurewidth,0\figureheight)},
scale only axis,
xmin=0.00000,
xmax=1000000.00000,
xlabel style={font=\color{mycolor1}},
xlabel={Zeit $\sym{t}$ [h]},
ymin=0.00000,
ymax=1.00000,
ylabel style={font=\color{mycolor1}},
ylabel={Zuverlässigkeit $\sym{R}(\sym{t})$},
axis background/.style={fill=white},
legend style={legend cell align=left, align=left},
grid=major,
grid style={dotted, gray!50},
scaled x ticks=base 10:-5,
x tick label style={/pgf/number format/fixed, /pgf/number format/precision=0}
]

\addplot[area legend, draw=none, fill=white!80!black, fill opacity=0.50000, forget plot]
table[row sep=crcr] {%
x	y\\
0.00000	1.00000\\
2004.00802	0.99951\\
4008.01603	0.99845\\
6012.02405	0.99694\\
8016.03206	0.99504\\
10020.04008	0.99278\\
12024.04810	0.99018\\
14028.05611	0.98725\\
16032.06413	0.98401\\
18036.07214	0.98047\\
20040.08016	0.97663\\
22044.08818	0.97252\\
24048.09619	0.96812\\
26052.10421	0.96345\\
28056.11222	0.95852\\
30060.12024	0.95333\\
32064.12826	0.94788\\
34068.13627	0.94219\\
36072.14429	0.93625\\
38076.15230	0.93007\\
40080.16032	0.92366\\
42084.16834	0.91702\\
44088.17635	0.91015\\
46092.18437	0.90305\\
48096.19238	0.89575\\
50100.20040	0.88822\\
52104.20842	0.88049\\
54108.21643	0.87256\\
56112.22445	0.86443\\
58116.23246	0.85610\\
60120.24048	0.84759\\
62124.24850	0.83889\\
64128.25651	0.83001\\
66132.26453	0.82095\\
68136.27255	0.81173\\
70140.28056	0.80234\\
72144.28858	0.79279\\
74148.29659	0.78309\\
76152.30461	0.77325\\
78156.31263	0.76326\\
80160.32064	0.75313\\
82164.32866	0.74287\\
84168.33667	0.73249\\
86172.34469	0.72199\\
88176.35271	0.71137\\
90180.36072	0.70065\\
92184.36874	0.68984\\
94188.37675	0.67892\\
96192.38477	0.66793\\
98196.39279	0.65685\\
100200.40080	0.64570\\
102204.40882	0.63448\\
104208.41683	0.62320\\
106212.42485	0.61188\\
108216.43287	0.60050\\
110220.44088	0.58909\\
112224.44890	0.57765\\
114228.45691	0.56619\\
116232.46493	0.55471\\
118236.47295	0.54322\\
120240.48096	0.53173\\
122244.48898	0.52025\\
124248.49699	0.50878\\
126252.50501	0.49733\\
128256.51303	0.48591\\
130260.52104	0.47452\\
132264.52906	0.46318\\
134268.53707	0.45188\\
136272.54509	0.44064\\
138276.55311	0.42947\\
140280.56112	0.41837\\
142284.56914	0.40734\\
144288.57715	0.39640\\
146292.58517	0.38555\\
148296.59319	0.37480\\
150300.60120	0.36415\\
152304.60922	0.35361\\
154308.61723	0.34319\\
156312.62525	0.33289\\
158316.63327	0.32272\\
160320.64128	0.31267\\
162324.64930	0.30277\\
164328.65731	0.29301\\
166332.66533	0.28340\\
168336.67335	0.27393\\
170340.68136	0.26463\\
172344.68938	0.25548\\
174348.69739	0.24650\\
176352.70541	0.23769\\
178356.71343	0.22905\\
180360.72144	0.22058\\
182364.72946	0.21229\\
184368.73747	0.20418\\
186372.74549	0.19624\\
188376.75351	0.18850\\
190380.76152	0.18093\\
192384.76954	0.17356\\
194388.77756	0.16637\\
196392.78557	0.15936\\
198396.79359	0.15255\\
200400.80160	0.14592\\
202404.80962	0.13949\\
204408.81764	0.13324\\
206412.82565	0.12718\\
208416.83367	0.12131\\
210420.84168	0.11562\\
212424.84970	0.11012\\
214428.85772	0.10480\\
216432.86573	0.09966\\
218436.87375	0.09470\\
220440.88176	0.08992\\
222444.88978	0.08531\\
224448.89780	0.08088\\
226452.90581	0.07661\\
228456.91383	0.07251\\
230460.92184	0.06858\\
232464.92986	0.06481\\
234468.93788	0.06119\\
236472.94589	0.05773\\
238476.95391	0.05442\\
240480.96192	0.05125\\
242484.96994	0.04823\\
244488.97796	0.04535\\
246492.98597	0.04260\\
248496.99399	0.03999\\
250501.00200	0.03750\\
252505.01002	0.03513\\
254509.01804	0.03289\\
256513.02605	0.03076\\
258517.03407	0.02874\\
260521.04208	0.02684\\
262525.05010	0.02503\\
264529.05812	0.02332\\
266533.06613	0.02172\\
268537.07415	0.02020\\
270541.08216	0.01877\\
272545.09018	0.01742\\
274549.09820	0.01616\\
276553.10621	0.01497\\
278557.11423	0.01386\\
280561.12224	0.01282\\
282565.13026	0.01184\\
284569.13828	0.01093\\
286573.14629	0.01007\\
288577.15431	0.00928\\
290581.16232	0.00854\\
292585.17034	0.00785\\
294589.17836	0.00721\\
296593.18637	0.00661\\
298597.19439	0.00606\\
300601.20240	0.00554\\
302605.21042	0.00507\\
304609.21844	0.00463\\
306613.22645	0.00422\\
308617.23447	0.00385\\
310621.24248	0.00350\\
312625.25050	0.00319\\
314629.25852	0.00290\\
316633.26653	0.00263\\
318637.27455	0.00238\\
320641.28257	0.00216\\
322645.29058	0.00195\\
324649.29860	0.00176\\
326653.30661	0.00159\\
328657.31463	0.00143\\
330661.32265	0.00129\\
332665.33066	0.00116\\
334669.33868	0.00104\\
336673.34669	0.00093\\
338677.35471	0.00084\\
340681.36273	0.00075\\
342685.37074	0.00067\\
344689.37876	0.00060\\
346693.38677	0.00053\\
348697.39479	0.00047\\
350701.40281	0.00042\\
352705.41082	0.00037\\
354709.41884	0.00033\\
356713.42685	0.00029\\
358717.43487	0.00026\\
360721.44289	0.00023\\
362725.45090	0.00020\\
364729.45892	0.00018\\
366733.46693	0.00016\\
368737.47495	0.00014\\
370741.48297	0.00012\\
372745.49098	0.00011\\
374749.49900	0.00009\\
376753.50701	0.00008\\
378757.51503	0.00007\\
380761.52305	0.00006\\
382765.53106	0.00005\\
384769.53908	0.00005\\
386773.54709	0.00004\\
388777.55511	0.00004\\
390781.56313	0.00003\\
392785.57114	0.00003\\
394789.57916	0.00002\\
396793.58717	0.00002\\
398797.59519	0.00002\\
400801.60321	0.00001\\
402805.61122	0.00001\\
404809.61924	0.00001\\
406813.62725	0.00001\\
408817.63527	0.00001\\
410821.64329	0.00001\\
412825.65130	0.00001\\
414829.65932	0.00000\\
416833.66733	0.00000\\
418837.67535	0.00000\\
420841.68337	0.00000\\
422845.69138	0.00000\\
424849.69940	0.00000\\
426853.70741	0.00000\\
428857.71543	0.00000\\
430861.72345	0.00000\\
432865.73146	0.00000\\
434869.73948	0.00000\\
436873.74749	0.00000\\
438877.75551	0.00000\\
440881.76353	0.00000\\
442885.77154	0.00000\\
444889.77956	0.00000\\
446893.78758	0.00000\\
448897.79559	0.00000\\
450901.80361	0.00000\\
452905.81162	0.00000\\
454909.81964	0.00000\\
456913.82766	0.00000\\
458917.83567	0.00000\\
460921.84369	0.00000\\
462925.85170	0.00000\\
464929.85972	0.00000\\
466933.86774	0.00000\\
468937.87575	0.00000\\
470941.88377	0.00000\\
472945.89178	0.00000\\
474949.89980	0.00000\\
476953.90782	0.00000\\
478957.91583	0.00000\\
480961.92385	0.00000\\
482965.93186	0.00000\\
484969.93988	0.00000\\
486973.94790	0.00000\\
488977.95591	0.00000\\
490981.96393	0.00000\\
492985.97194	0.00000\\
494989.97996	0.00000\\
496993.98798	0.00000\\
498997.99599	0.00000\\
501002.00401	0.00000\\
503006.01202	0.00000\\
505010.02004	0.00000\\
507014.02806	0.00000\\
509018.03607	0.00000\\
511022.04409	0.00000\\
513026.05210	0.00000\\
515030.06012	0.00000\\
517034.06814	0.00000\\
519038.07615	0.00000\\
521042.08417	0.00000\\
523046.09218	0.00000\\
525050.10020	0.00000\\
527054.10822	0.00000\\
529058.11623	0.00000\\
531062.12425	0.00000\\
533066.13226	0.00000\\
535070.14028	0.00000\\
537074.14830	0.00000\\
539078.15631	0.00000\\
541082.16433	0.00000\\
543086.17234	0.00000\\
545090.18036	0.00000\\
547094.18838	0.00000\\
549098.19639	0.00000\\
551102.20441	0.00000\\
553106.21242	0.00000\\
555110.22044	0.00000\\
557114.22846	0.00000\\
559118.23647	0.00000\\
561122.24449	0.00000\\
563126.25251	0.00000\\
565130.26052	0.00000\\
567134.26854	0.00000\\
569138.27655	0.00000\\
571142.28457	0.00000\\
573146.29259	0.00000\\
575150.30060	0.00000\\
577154.30862	0.00000\\
579158.31663	0.00000\\
581162.32465	0.00000\\
583166.33267	0.00000\\
585170.34068	0.00000\\
587174.34870	0.00000\\
589178.35671	0.00000\\
591182.36473	0.00000\\
593186.37275	0.00000\\
595190.38076	0.00000\\
597194.38878	0.00000\\
599198.39679	0.00000\\
601202.40481	0.00000\\
603206.41283	0.00000\\
605210.42084	0.00000\\
607214.42886	0.00000\\
609218.43687	0.00000\\
611222.44489	0.00000\\
613226.45291	0.00000\\
615230.46092	0.00000\\
617234.46894	0.00000\\
619238.47695	0.00000\\
621242.48497	0.00000\\
623246.49299	0.00000\\
625250.50100	0.00000\\
627254.50902	0.00000\\
629258.51703	0.00000\\
631262.52505	0.00000\\
633266.53307	0.00000\\
635270.54108	0.00000\\
637274.54910	0.00000\\
639278.55711	0.00000\\
641282.56513	0.00000\\
643286.57315	0.00000\\
645290.58116	0.00000\\
647294.58918	0.00000\\
649298.59719	0.00000\\
651302.60521	0.00000\\
653306.61323	0.00000\\
655310.62124	0.00000\\
657314.62926	0.00000\\
659318.63727	0.00000\\
661322.64529	0.00000\\
663326.65331	0.00000\\
665330.66132	0.00000\\
667334.66934	0.00000\\
669338.67735	0.00000\\
671342.68537	0.00000\\
673346.69339	0.00000\\
675350.70140	0.00000\\
677354.70942	0.00000\\
679358.71743	0.00000\\
681362.72545	0.00000\\
683366.73347	0.00000\\
685370.74148	0.00000\\
687374.74950	0.00000\\
689378.75752	0.00000\\
691382.76553	0.00000\\
693386.77355	0.00000\\
695390.78156	0.00000\\
697394.78958	0.00000\\
699398.79760	0.00000\\
701402.80561	0.00000\\
703406.81363	0.00000\\
705410.82164	0.00000\\
707414.82966	0.00000\\
709418.83768	0.00000\\
711422.84569	0.00000\\
713426.85371	0.00000\\
715430.86172	0.00000\\
717434.86974	0.00000\\
719438.87776	0.00000\\
721442.88577	0.00000\\
723446.89379	0.00000\\
725450.90180	0.00000\\
727454.90982	0.00000\\
729458.91784	0.00000\\
731462.92585	0.00000\\
733466.93387	0.00000\\
735470.94188	0.00000\\
737474.94990	0.00000\\
739478.95792	0.00000\\
741482.96593	0.00000\\
743486.97395	0.00000\\
745490.98196	0.00000\\
747494.98998	0.00000\\
749498.99800	0.00000\\
751503.00601	0.00000\\
753507.01403	0.00000\\
755511.02204	0.00000\\
757515.03006	0.00000\\
759519.03808	0.00000\\
761523.04609	0.00000\\
763527.05411	0.00000\\
765531.06212	0.00000\\
767535.07014	0.00000\\
769539.07816	0.00000\\
771543.08617	0.00000\\
773547.09419	0.00000\\
775551.10220	0.00000\\
777555.11022	0.00000\\
779559.11824	0.00000\\
781563.12625	0.00000\\
783567.13427	0.00000\\
785571.14228	0.00000\\
787575.15030	0.00000\\
789579.15832	0.00000\\
791583.16633	0.00000\\
793587.17435	0.00000\\
795591.18236	0.00000\\
797595.19038	0.00000\\
799599.19840	0.00000\\
801603.20641	0.00000\\
803607.21443	0.00000\\
805611.22244	0.00000\\
807615.23046	0.00000\\
809619.23848	0.00000\\
811623.24649	0.00000\\
813627.25451	0.00000\\
815631.26253	0.00000\\
817635.27054	0.00000\\
819639.27856	0.00000\\
821643.28657	0.00000\\
823647.29459	0.00000\\
825651.30261	0.00000\\
827655.31062	0.00000\\
829659.31864	0.00000\\
831663.32665	0.00000\\
833667.33467	0.00000\\
835671.34269	0.00000\\
837675.35070	0.00000\\
839679.35872	0.00000\\
841683.36673	0.00000\\
843687.37475	0.00000\\
845691.38277	0.00000\\
847695.39078	0.00000\\
849699.39880	0.00000\\
851703.40681	0.00000\\
853707.41483	0.00000\\
855711.42285	0.00000\\
857715.43086	0.00000\\
859719.43888	0.00000\\
861723.44689	0.00000\\
863727.45491	0.00000\\
865731.46293	0.00000\\
867735.47094	0.00000\\
869739.47896	0.00000\\
871743.48697	0.00000\\
873747.49499	0.00000\\
875751.50301	0.00000\\
877755.51102	0.00000\\
879759.51904	0.00000\\
881763.52705	0.00000\\
883767.53507	0.00000\\
885771.54309	0.00000\\
887775.55110	0.00000\\
889779.55912	0.00000\\
891783.56713	0.00000\\
893787.57515	0.00000\\
895791.58317	0.00000\\
897795.59118	0.00000\\
899799.59920	0.00000\\
901803.60721	0.00000\\
903807.61523	0.00000\\
905811.62325	0.00000\\
907815.63126	0.00000\\
909819.63928	0.00000\\
911823.64729	0.00000\\
913827.65531	0.00000\\
915831.66333	0.00000\\
917835.67134	0.00000\\
919839.67936	0.00000\\
921843.68737	0.00000\\
923847.69539	0.00000\\
925851.70341	0.00000\\
927855.71142	0.00000\\
929859.71944	0.00000\\
931863.72745	0.00000\\
933867.73547	0.00000\\
935871.74349	0.00000\\
937875.75150	0.00000\\
939879.75952	0.00000\\
941883.76754	0.00000\\
943887.77555	0.00000\\
945891.78357	0.00000\\
947895.79158	0.00000\\
949899.79960	0.00000\\
951903.80762	0.00000\\
953907.81563	0.00000\\
955911.82365	0.00000\\
957915.83166	0.00000\\
959919.83968	0.00000\\
961923.84770	0.00000\\
963927.85571	0.00000\\
965931.86373	0.00000\\
967935.87174	0.00000\\
969939.87976	0.00000\\
971943.88778	0.00000\\
973947.89579	0.00000\\
975951.90381	0.00000\\
977955.91182	0.00000\\
979959.91984	0.00000\\
981963.92786	0.00000\\
983967.93587	0.00000\\
985971.94389	0.00000\\
987975.95190	0.00000\\
989979.95992	0.00000\\
991983.96794	0.00000\\
993987.97595	0.00000\\
995991.98397	0.00000\\
997995.99198	0.00000\\
1000000.00000	0.00000\\
1000000.00000	0.00023\\
997995.99198	0.00024\\
995991.98397	0.00025\\
993987.97595	0.00025\\
991983.96794	0.00026\\
989979.95992	0.00027\\
987975.95190	0.00028\\
985971.94389	0.00029\\
983967.93587	0.00030\\
981963.92786	0.00031\\
979959.91984	0.00031\\
977955.91182	0.00032\\
975951.90381	0.00033\\
973947.89579	0.00034\\
971943.88778	0.00035\\
969939.87976	0.00037\\
967935.87174	0.00038\\
965931.86373	0.00039\\
963927.85571	0.00040\\
961923.84770	0.00041\\
959919.83968	0.00042\\
957915.83166	0.00044\\
955911.82365	0.00045\\
953907.81563	0.00046\\
951903.80762	0.00048\\
949899.79960	0.00049\\
947895.79158	0.00051\\
945891.78357	0.00052\\
943887.77555	0.00054\\
941883.76754	0.00055\\
939879.75952	0.00057\\
937875.75150	0.00058\\
935871.74349	0.00060\\
933867.73547	0.00062\\
931863.72745	0.00064\\
929859.71944	0.00066\\
927855.71142	0.00067\\
925851.70341	0.00069\\
923847.69539	0.00071\\
921843.68737	0.00074\\
919839.67936	0.00076\\
917835.67134	0.00078\\
915831.66333	0.00080\\
913827.65531	0.00082\\
911823.64729	0.00085\\
909819.63928	0.00087\\
907815.63126	0.00090\\
905811.62325	0.00092\\
903807.61523	0.00095\\
901803.60721	0.00098\\
899799.59920	0.00100\\
897795.59118	0.00103\\
895791.58317	0.00106\\
893787.57515	0.00109\\
891783.56713	0.00112\\
889779.55912	0.00116\\
887775.55110	0.00119\\
885771.54309	0.00122\\
883767.53507	0.00126\\
881763.52705	0.00129\\
879759.51904	0.00133\\
877755.51102	0.00137\\
875751.50301	0.00140\\
873747.49499	0.00144\\
871743.48697	0.00148\\
869739.47896	0.00152\\
867735.47094	0.00157\\
865731.46293	0.00161\\
863727.45491	0.00165\\
861723.44689	0.00170\\
859719.43888	0.00175\\
857715.43086	0.00180\\
855711.42285	0.00184\\
853707.41483	0.00190\\
851703.40681	0.00195\\
849699.39880	0.00200\\
847695.39078	0.00206\\
845691.38277	0.00211\\
843687.37475	0.00217\\
841683.36673	0.00223\\
839679.35872	0.00229\\
837675.35070	0.00235\\
835671.34269	0.00241\\
833667.33467	0.00248\\
831663.32665	0.00255\\
829659.31864	0.00261\\
827655.31062	0.00268\\
825651.30261	0.00276\\
823647.29459	0.00283\\
821643.28657	0.00291\\
819639.27856	0.00298\\
817635.27054	0.00306\\
815631.26253	0.00314\\
813627.25451	0.00323\\
811623.24649	0.00331\\
809619.23848	0.00340\\
807615.23046	0.00349\\
805611.22244	0.00358\\
803607.21443	0.00367\\
801603.20641	0.00377\\
799599.19840	0.00387\\
797595.19038	0.00397\\
795591.18236	0.00407\\
793587.17435	0.00418\\
791583.16633	0.00429\\
789579.15832	0.00440\\
787575.15030	0.00451\\
785571.14228	0.00463\\
783567.13427	0.00475\\
781563.12625	0.00487\\
779559.11824	0.00499\\
777555.11022	0.00512\\
775551.10220	0.00525\\
773547.09419	0.00539\\
771543.08617	0.00552\\
769539.07816	0.00566\\
767535.07014	0.00581\\
765531.06212	0.00596\\
763527.05411	0.00611\\
761523.04609	0.00626\\
759519.03808	0.00642\\
757515.03006	0.00658\\
755511.02204	0.00674\\
753507.01403	0.00691\\
751503.00601	0.00708\\
749498.99800	0.00726\\
747494.98998	0.00744\\
745490.98196	0.00762\\
743486.97395	0.00781\\
741482.96593	0.00801\\
739478.95792	0.00820\\
737474.94990	0.00841\\
735470.94188	0.00861\\
733466.93387	0.00882\\
731462.92585	0.00904\\
729458.91784	0.00926\\
727454.90982	0.00948\\
725450.90180	0.00971\\
723446.89379	0.00995\\
721442.88577	0.01019\\
719438.87776	0.01044\\
717434.86974	0.01069\\
715430.86172	0.01095\\
713426.85371	0.01121\\
711422.84569	0.01148\\
709418.83768	0.01175\\
707414.82966	0.01203\\
705410.82164	0.01232\\
703406.81363	0.01261\\
701402.80561	0.01291\\
699398.79760	0.01321\\
697394.78958	0.01352\\
695390.78156	0.01384\\
693386.77355	0.01417\\
691382.76553	0.01450\\
689378.75752	0.01484\\
687374.74950	0.01518\\
685370.74148	0.01554\\
683366.73347	0.01590\\
681362.72545	0.01627\\
679358.71743	0.01664\\
677354.70942	0.01702\\
675350.70140	0.01742\\
673346.69339	0.01782\\
671342.68537	0.01822\\
669338.67735	0.01864\\
667334.66934	0.01907\\
665330.66132	0.01950\\
663326.65331	0.01994\\
661322.64529	0.02039\\
659318.63727	0.02085\\
657314.62926	0.02132\\
655310.62124	0.02180\\
653306.61323	0.02229\\
651302.60521	0.02279\\
649298.59719	0.02330\\
647294.58918	0.02382\\
645290.58116	0.02435\\
643286.57315	0.02489\\
641282.56513	0.02544\\
639278.55711	0.02600\\
637274.54910	0.02657\\
635270.54108	0.02715\\
633266.53307	0.02775\\
631262.52505	0.02835\\
629258.51703	0.02897\\
627254.50902	0.02960\\
625250.50100	0.03024\\
623246.49299	0.03090\\
621242.48497	0.03157\\
619238.47695	0.03224\\
617234.46894	0.03294\\
615230.46092	0.03364\\
613226.45291	0.03436\\
611222.44489	0.03509\\
609218.43687	0.03584\\
607214.42886	0.03660\\
605210.42084	0.03737\\
603206.41283	0.03816\\
601202.40481	0.03896\\
599198.39679	0.03978\\
597194.38878	0.04061\\
595190.38076	0.04146\\
593186.37275	0.04232\\
591182.36473	0.04320\\
589178.35671	0.04410\\
587174.34870	0.04501\\
585170.34068	0.04593\\
583166.33267	0.04688\\
581162.32465	0.04784\\
579158.31663	0.04882\\
577154.30862	0.04981\\
575150.30060	0.05082\\
573146.29259	0.05185\\
571142.28457	0.05290\\
569138.27655	0.05397\\
567134.26854	0.05505\\
565130.26052	0.05615\\
563126.25251	0.05728\\
561122.24449	0.05842\\
559118.23647	0.05958\\
557114.22846	0.06076\\
555110.22044	0.06196\\
553106.21242	0.06318\\
551102.20441	0.06442\\
549098.19639	0.06568\\
547094.18838	0.06697\\
545090.18036	0.06827\\
543086.17234	0.06960\\
541082.16433	0.07095\\
539078.15631	0.07232\\
537074.14830	0.07371\\
535070.14028	0.07512\\
533066.13226	0.07656\\
531062.12425	0.07802\\
529058.11623	0.07951\\
527054.10822	0.08101\\
525050.10020	0.08255\\
523046.09218	0.08410\\
521042.08417	0.08568\\
519038.07615	0.08729\\
517034.06814	0.08892\\
515030.06012	0.09057\\
513026.05210	0.09225\\
511022.04409	0.09396\\
509018.03607	0.09569\\
507014.02806	0.09745\\
505010.02004	0.09924\\
503006.01202	0.10105\\
501002.00401	0.10289\\
498997.99599	0.10476\\
496993.98798	0.10666\\
494989.97996	0.10858\\
492985.97194	0.11053\\
490981.96393	0.11251\\
488977.95591	0.11452\\
486973.94790	0.11656\\
484969.93988	0.11863\\
482965.93186	0.12073\\
480961.92385	0.12285\\
478957.91583	0.12501\\
476953.90782	0.12720\\
474949.89980	0.12942\\
472945.89178	0.13167\\
470941.88377	0.13395\\
468937.87575	0.13626\\
466933.86774	0.13860\\
464929.85972	0.14098\\
462925.85170	0.14339\\
460921.84369	0.14583\\
458917.83567	0.14830\\
456913.82766	0.15081\\
454909.81964	0.15335\\
452905.81162	0.15592\\
450901.80361	0.15852\\
448897.79559	0.16116\\
446893.78758	0.16384\\
444889.77956	0.16655\\
442885.77154	0.16929\\
440881.76353	0.17207\\
438877.75551	0.17488\\
436873.74749	0.17772\\
434869.73948	0.18061\\
432865.73146	0.18352\\
430861.72345	0.18648\\
428857.71543	0.18947\\
426853.70741	0.19249\\
424849.69940	0.19556\\
422845.69138	0.19865\\
420841.68337	0.20179\\
418837.67535	0.20496\\
416833.66733	0.20817\\
414829.65932	0.21141\\
412825.65130	0.21469\\
410821.64329	0.21801\\
408817.63527	0.22137\\
406813.62725	0.22476\\
404809.61924	0.22820\\
402805.61122	0.23166\\
400801.60321	0.23517\\
398797.59519	0.23872\\
396793.58717	0.24230\\
394789.57916	0.24592\\
392785.57114	0.24958\\
390781.56313	0.25328\\
388777.55511	0.25701\\
386773.54709	0.26078\\
384769.53908	0.26460\\
382765.53106	0.26845\\
380761.52305	0.27233\\
378757.51503	0.27626\\
376753.50701	0.28022\\
374749.49900	0.28423\\
372745.49098	0.28827\\
370741.48297	0.29234\\
368737.47495	0.29646\\
366733.46693	0.30061\\
364729.45892	0.30481\\
362725.45090	0.30904\\
360721.44289	0.31330\\
358717.43487	0.31761\\
356713.42685	0.32195\\
354709.41884	0.32633\\
352705.41082	0.33075\\
350701.40281	0.33520\\
348697.39479	0.33969\\
346693.38677	0.34422\\
344689.37876	0.34878\\
342685.37074	0.35338\\
340681.36273	0.35802\\
338677.35471	0.36269\\
336673.34669	0.36740\\
334669.33868	0.37214\\
332665.33066	0.37691\\
330661.32265	0.38172\\
328657.31463	0.38657\\
326653.30661	0.39145\\
324649.29860	0.39636\\
322645.29058	0.40130\\
320641.28257	0.40628\\
318637.27455	0.41129\\
316633.26653	0.41633\\
314629.25852	0.42141\\
312625.25050	0.42651\\
310621.24248	0.43164\\
308617.23447	0.43681\\
306613.22645	0.44200\\
304609.21844	0.44722\\
302605.21042	0.45248\\
300601.20240	0.45775\\
298597.19439	0.46306\\
296593.18637	0.46839\\
294589.17836	0.47375\\
292585.17034	0.47913\\
290581.16232	0.48454\\
288577.15431	0.48997\\
286573.14629	0.49543\\
284569.13828	0.50091\\
282565.13026	0.50641\\
280561.12224	0.51193\\
278557.11423	0.51747\\
276553.10621	0.52303\\
274549.09820	0.52861\\
272545.09018	0.53421\\
270541.08216	0.53982\\
268537.07415	0.54545\\
266533.06613	0.55110\\
264529.05812	0.55676\\
262525.05010	0.56243\\
260521.04208	0.56811\\
258517.03407	0.57381\\
256513.02605	0.57952\\
254509.01804	0.58523\\
252505.01002	0.59096\\
250501.00200	0.59669\\
248496.99399	0.60243\\
246492.98597	0.60817\\
244488.97796	0.61392\\
242484.96994	0.61967\\
240480.96192	0.62542\\
238476.95391	0.63118\\
236472.94589	0.63693\\
234468.93788	0.64268\\
232464.92986	0.64843\\
230460.92184	0.65417\\
228456.91383	0.65991\\
226452.90581	0.66564\\
224448.89780	0.67136\\
222444.88978	0.67708\\
220440.88176	0.68278\\
218436.87375	0.68847\\
216432.86573	0.69415\\
214428.85772	0.69981\\
212424.84970	0.70546\\
210420.84168	0.71109\\
208416.83367	0.71670\\
206412.82565	0.72229\\
204408.81764	0.72786\\
202404.80962	0.73340\\
200400.80160	0.73892\\
198396.79359	0.74442\\
196392.78557	0.74988\\
194388.77756	0.75532\\
192384.76954	0.76073\\
190380.76152	0.76611\\
188376.75351	0.77145\\
186372.74549	0.77676\\
184368.73747	0.78203\\
182364.72946	0.78727\\
180360.72144	0.79246\\
178356.71343	0.79762\\
176352.70541	0.80273\\
174348.69739	0.80780\\
172344.68938	0.81282\\
170340.68136	0.81780\\
168336.67335	0.82273\\
166332.66533	0.82761\\
164328.65731	0.83244\\
162324.64930	0.83722\\
160320.64128	0.84194\\
158316.63327	0.84661\\
156312.62525	0.85123\\
154308.61723	0.85578\\
152304.60922	0.86028\\
150300.60120	0.86471\\
148296.59319	0.86909\\
146292.58517	0.87340\\
144288.57715	0.87765\\
142284.56914	0.88183\\
140280.56112	0.88595\\
138276.55311	0.89000\\
136272.54509	0.89398\\
134268.53707	0.89789\\
132264.52906	0.90174\\
130260.52104	0.90551\\
128256.51303	0.90920\\
126252.50501	0.91283\\
124248.49699	0.91638\\
122244.48898	0.91986\\
120240.48096	0.92326\\
118236.47295	0.92658\\
116232.46493	0.92983\\
114228.45691	0.93300\\
112224.44890	0.93609\\
110220.44088	0.93910\\
108216.43287	0.94204\\
106212.42485	0.94490\\
104208.41683	0.94767\\
102204.40882	0.95037\\
100200.40080	0.95299\\
98196.39279	0.95553\\
96192.38477	0.95799\\
94188.37675	0.96036\\
92184.36874	0.96266\\
90180.36072	0.96488\\
88176.35271	0.96702\\
86172.34469	0.96909\\
84168.33667	0.97107\\
82164.32866	0.97298\\
80160.32064	0.97481\\
78156.31263	0.97656\\
76152.30461	0.97823\\
74148.29659	0.97984\\
72144.28858	0.98136\\
70140.28056	0.98282\\
68136.27255	0.98420\\
66132.26453	0.98551\\
64128.25651	0.98675\\
62124.24850	0.98792\\
60120.24048	0.98903\\
58116.23246	0.99007\\
56112.22445	0.99104\\
54108.21643	0.99195\\
52104.20842	0.99280\\
50100.20040	0.99359\\
48096.19238	0.99433\\
46092.18437	0.99500\\
44088.17635	0.99562\\
42084.16834	0.99619\\
40080.16032	0.99671\\
38076.15230	0.99718\\
36072.14429	0.99760\\
34068.13627	0.99798\\
32064.12826	0.99832\\
30060.12024	0.99862\\
28056.11222	0.99888\\
26052.10421	0.99911\\
24048.09619	0.99930\\
22044.08818	0.99946\\
20040.08016	0.99960\\
18036.07214	0.99971\\
16032.06413	0.99980\\
14028.05611	0.99987\\
12024.04810	0.99992\\
10020.04008	0.99995\\
8016.03206	0.99998\\
6012.02405	0.99999\\
4008.01603	1.00000\\
2004.00802	1.00000\\
0.00000	1.00000\\
}--cycle;
\addplot [color=gray, line width=1.5pt]
  table[row sep=crcr]{%
0.00000	1.00000\\
2004.00802	0.99999\\
4008.01603	0.99993\\
6012.02405	0.99983\\
8016.03206	0.99966\\
10020.04008	0.99942\\
12024.04810	0.99910\\
14028.05611	0.99869\\
16032.06413	0.99820\\
18036.07214	0.99761\\
20040.08016	0.99693\\
22044.08818	0.99614\\
24048.09619	0.99525\\
26052.10421	0.99425\\
28056.11222	0.99314\\
30060.12024	0.99191\\
32064.12826	0.99056\\
34068.13627	0.98910\\
36072.14429	0.98751\\
38076.15230	0.98580\\
40080.16032	0.98395\\
42084.16834	0.98198\\
44088.17635	0.97988\\
46092.18437	0.97765\\
48096.19238	0.97528\\
50100.20040	0.97278\\
52104.20842	0.97014\\
54108.21643	0.96736\\
56112.22445	0.96444\\
58116.23246	0.96138\\
60120.24048	0.95819\\
62124.24850	0.95485\\
64128.25651	0.95137\\
66132.26453	0.94775\\
68136.27255	0.94399\\
70140.28056	0.94009\\
72144.28858	0.93605\\
74148.29659	0.93186\\
76152.30461	0.92753\\
78156.31263	0.92307\\
80160.32064	0.91846\\
82164.32866	0.91371\\
84168.33667	0.90883\\
86172.34469	0.90381\\
88176.35271	0.89865\\
90180.36072	0.89336\\
92184.36874	0.88793\\
94188.37675	0.88237\\
96192.38477	0.87668\\
98196.39279	0.87086\\
100200.40080	0.86491\\
102204.40882	0.85883\\
104208.41683	0.85264\\
106212.42485	0.84632\\
108216.43287	0.83988\\
110220.44088	0.83332\\
112224.44890	0.82665\\
114228.45691	0.81986\\
116232.46493	0.81297\\
118236.47295	0.80596\\
120240.48096	0.79886\\
122244.48898	0.79164\\
124248.49699	0.78434\\
126252.50501	0.77693\\
128256.51303	0.76943\\
130260.52104	0.76184\\
132264.52906	0.75416\\
134268.53707	0.74640\\
136272.54509	0.73856\\
138276.55311	0.73064\\
140280.56112	0.72264\\
142284.56914	0.71458\\
144288.57715	0.70645\\
146292.58517	0.69825\\
148296.59319	0.68999\\
150300.60120	0.68168\\
152304.60922	0.67332\\
154308.61723	0.66490\\
156312.62525	0.65644\\
158316.63327	0.64794\\
160320.64128	0.63940\\
162324.64930	0.63082\\
164328.65731	0.62221\\
166332.66533	0.61358\\
168336.67335	0.60492\\
170340.68136	0.59625\\
172344.68938	0.58755\\
174348.69739	0.57885\\
176352.70541	0.57013\\
178356.71343	0.56141\\
180360.72144	0.55269\\
182364.72946	0.54398\\
184368.73747	0.53527\\
186372.74549	0.52656\\
188376.75351	0.51788\\
190380.76152	0.50920\\
192384.76954	0.50055\\
194388.77756	0.49192\\
196392.78557	0.48332\\
198396.79359	0.47475\\
200400.80160	0.46622\\
202404.80962	0.45772\\
204408.81764	0.44926\\
206412.82565	0.44084\\
208416.83367	0.43247\\
210420.84168	0.42415\\
212424.84970	0.41588\\
214428.85772	0.40766\\
216432.86573	0.39950\\
218436.87375	0.39141\\
220440.88176	0.38337\\
222444.88978	0.37541\\
224448.89780	0.36751\\
226452.90581	0.35968\\
228456.91383	0.35192\\
230460.92184	0.34424\\
232464.92986	0.33663\\
234468.93788	0.32911\\
236472.94589	0.32166\\
238476.95391	0.31430\\
240480.96192	0.30702\\
242484.96994	0.29983\\
244488.97796	0.29273\\
246492.98597	0.28572\\
248496.99399	0.27879\\
250501.00200	0.27196\\
252505.01002	0.26523\\
254509.01804	0.25859\\
256513.02605	0.25204\\
258517.03407	0.24559\\
260521.04208	0.23924\\
262525.05010	0.23298\\
264529.05812	0.22683\\
266533.06613	0.22078\\
268537.07415	0.21482\\
270541.08216	0.20897\\
272545.09018	0.20322\\
274549.09820	0.19756\\
276553.10621	0.19202\\
278557.11423	0.18657\\
280561.12224	0.18122\\
282565.13026	0.17598\\
284569.13828	0.17084\\
286573.14629	0.16580\\
288577.15431	0.16086\\
290581.16232	0.15602\\
292585.17034	0.15129\\
294589.17836	0.14665\\
296593.18637	0.14212\\
298597.19439	0.13768\\
300601.20240	0.13334\\
302605.21042	0.12910\\
304609.21844	0.12496\\
306613.22645	0.12091\\
308617.23447	0.11696\\
310621.24248	0.11311\\
312625.25050	0.10935\\
314629.25852	0.10568\\
316633.26653	0.10210\\
318637.27455	0.09862\\
320641.28257	0.09522\\
322645.29058	0.09191\\
324649.29860	0.08869\\
326653.30661	0.08556\\
328657.31463	0.08251\\
330661.32265	0.07955\\
332665.33066	0.07667\\
334669.33868	0.07387\\
336673.34669	0.07115\\
338677.35471	0.06851\\
340681.36273	0.06594\\
342685.37074	0.06346\\
344689.37876	0.06104\\
346693.38677	0.05870\\
348697.39479	0.05643\\
350701.40281	0.05424\\
352705.41082	0.05211\\
354709.41884	0.05005\\
356713.42685	0.04805\\
358717.43487	0.04612\\
360721.44289	0.04426\\
362725.45090	0.04245\\
364729.45892	0.04071\\
366733.46693	0.03902\\
368737.47495	0.03739\\
370741.48297	0.03582\\
372745.49098	0.03431\\
374749.49900	0.03284\\
376753.50701	0.03143\\
378757.51503	0.03007\\
380761.52305	0.02876\\
382765.53106	0.02750\\
384769.53908	0.02628\\
386773.54709	0.02511\\
388777.55511	0.02399\\
390781.56313	0.02290\\
392785.57114	0.02186\\
394789.57916	0.02086\\
396793.58717	0.01990\\
398797.59519	0.01898\\
400801.60321	0.01809\\
402805.61122	0.01724\\
404809.61924	0.01642\\
406813.62725	0.01564\\
408817.63527	0.01489\\
410821.64329	0.01417\\
412825.65130	0.01348\\
414829.65932	0.01282\\
416833.66733	0.01218\\
418837.67535	0.01158\\
420841.68337	0.01100\\
422845.69138	0.01045\\
424849.69940	0.00992\\
426853.70741	0.00941\\
428857.71543	0.00893\\
430861.72345	0.00847\\
432865.73146	0.00803\\
434869.73948	0.00761\\
436873.74749	0.00721\\
438877.75551	0.00683\\
440881.76353	0.00647\\
442885.77154	0.00612\\
444889.77956	0.00579\\
446893.78758	0.00548\\
448897.79559	0.00518\\
450901.80361	0.00489\\
452905.81162	0.00462\\
454909.81964	0.00437\\
456913.82766	0.00412\\
458917.83567	0.00389\\
460921.84369	0.00367\\
462925.85170	0.00346\\
464929.85972	0.00326\\
466933.86774	0.00308\\
468937.87575	0.00290\\
470941.88377	0.00273\\
472945.89178	0.00257\\
474949.89980	0.00242\\
476953.90782	0.00228\\
478957.91583	0.00214\\
480961.92385	0.00201\\
482965.93186	0.00189\\
484969.93988	0.00178\\
486973.94790	0.00167\\
488977.95591	0.00157\\
490981.96393	0.00147\\
492985.97194	0.00138\\
494989.97996	0.00129\\
496993.98798	0.00121\\
498997.99599	0.00114\\
501002.00401	0.00106\\
503006.01202	0.00100\\
505010.02004	0.00093\\
507014.02806	0.00087\\
509018.03607	0.00082\\
511022.04409	0.00076\\
513026.05210	0.00071\\
515030.06012	0.00067\\
517034.06814	0.00062\\
519038.07615	0.00058\\
521042.08417	0.00054\\
523046.09218	0.00051\\
525050.10020	0.00047\\
527054.10822	0.00044\\
529058.11623	0.00041\\
531062.12425	0.00038\\
533066.13226	0.00036\\
535070.14028	0.00033\\
537074.14830	0.00031\\
539078.15631	0.00029\\
541082.16433	0.00027\\
543086.17234	0.00025\\
545090.18036	0.00023\\
547094.18838	0.00021\\
549098.19639	0.00020\\
551102.20441	0.00018\\
553106.21242	0.00017\\
555110.22044	0.00016\\
557114.22846	0.00015\\
559118.23647	0.00014\\
561122.24449	0.00013\\
563126.25251	0.00012\\
565130.26052	0.00011\\
567134.26854	0.00010\\
569138.27655	0.00009\\
571142.28457	0.00009\\
573146.29259	0.00008\\
575150.30060	0.00007\\
577154.30862	0.00007\\
579158.31663	0.00006\\
581162.32465	0.00006\\
583166.33267	0.00005\\
585170.34068	0.00005\\
587174.34870	0.00004\\
589178.35671	0.00004\\
591182.36473	0.00004\\
593186.37275	0.00004\\
595190.38076	0.00003\\
597194.38878	0.00003\\
599198.39679	0.00003\\
601202.40481	0.00003\\
603206.41283	0.00002\\
605210.42084	0.00002\\
607214.42886	0.00002\\
609218.43687	0.00002\\
611222.44489	0.00002\\
613226.45291	0.00001\\
615230.46092	0.00001\\
617234.46894	0.00001\\
619238.47695	0.00001\\
621242.48497	0.00001\\
623246.49299	0.00001\\
625250.50100	0.00001\\
627254.50902	0.00001\\
629258.51703	0.00001\\
631262.52505	0.00001\\
633266.53307	0.00001\\
635270.54108	0.00001\\
637274.54910	0.00001\\
639278.55711	0.00000\\
641282.56513	0.00000\\
643286.57315	0.00000\\
645290.58116	0.00000\\
647294.58918	0.00000\\
649298.59719	0.00000\\
651302.60521	0.00000\\
653306.61323	0.00000\\
655310.62124	0.00000\\
657314.62926	0.00000\\
659318.63727	0.00000\\
661322.64529	0.00000\\
663326.65331	0.00000\\
665330.66132	0.00000\\
667334.66934	0.00000\\
669338.67735	0.00000\\
671342.68537	0.00000\\
673346.69339	0.00000\\
675350.70140	0.00000\\
677354.70942	0.00000\\
679358.71743	0.00000\\
681362.72545	0.00000\\
683366.73347	0.00000\\
685370.74148	0.00000\\
687374.74950	0.00000\\
689378.75752	0.00000\\
691382.76553	0.00000\\
693386.77355	0.00000\\
695390.78156	0.00000\\
697394.78958	0.00000\\
699398.79760	0.00000\\
701402.80561	0.00000\\
703406.81363	0.00000\\
705410.82164	0.00000\\
707414.82966	0.00000\\
709418.83768	0.00000\\
711422.84569	0.00000\\
713426.85371	0.00000\\
715430.86172	0.00000\\
717434.86974	0.00000\\
719438.87776	0.00000\\
721442.88577	0.00000\\
723446.89379	0.00000\\
725450.90180	0.00000\\
727454.90982	0.00000\\
729458.91784	0.00000\\
731462.92585	0.00000\\
733466.93387	0.00000\\
735470.94188	0.00000\\
737474.94990	0.00000\\
739478.95792	0.00000\\
741482.96593	0.00000\\
743486.97395	0.00000\\
745490.98196	0.00000\\
747494.98998	0.00000\\
749498.99800	0.00000\\
751503.00601	0.00000\\
753507.01403	0.00000\\
755511.02204	0.00000\\
757515.03006	0.00000\\
759519.03808	0.00000\\
761523.04609	0.00000\\
763527.05411	0.00000\\
765531.06212	0.00000\\
767535.07014	0.00000\\
769539.07816	0.00000\\
771543.08617	0.00000\\
773547.09419	0.00000\\
775551.10220	0.00000\\
777555.11022	0.00000\\
779559.11824	0.00000\\
781563.12625	0.00000\\
783567.13427	0.00000\\
785571.14228	0.00000\\
787575.15030	0.00000\\
789579.15832	0.00000\\
791583.16633	0.00000\\
793587.17435	0.00000\\
795591.18236	0.00000\\
797595.19038	0.00000\\
799599.19840	0.00000\\
801603.20641	0.00000\\
803607.21443	0.00000\\
805611.22244	0.00000\\
807615.23046	0.00000\\
809619.23848	0.00000\\
811623.24649	0.00000\\
813627.25451	0.00000\\
815631.26253	0.00000\\
817635.27054	0.00000\\
819639.27856	0.00000\\
821643.28657	0.00000\\
823647.29459	0.00000\\
825651.30261	0.00000\\
827655.31062	0.00000\\
829659.31864	0.00000\\
831663.32665	0.00000\\
833667.33467	0.00000\\
835671.34269	0.00000\\
837675.35070	0.00000\\
839679.35872	0.00000\\
841683.36673	0.00000\\
843687.37475	0.00000\\
845691.38277	0.00000\\
847695.39078	0.00000\\
849699.39880	0.00000\\
851703.40681	0.00000\\
853707.41483	0.00000\\
855711.42285	0.00000\\
857715.43086	0.00000\\
859719.43888	0.00000\\
861723.44689	0.00000\\
863727.45491	0.00000\\
865731.46293	0.00000\\
867735.47094	0.00000\\
869739.47896	0.00000\\
871743.48697	0.00000\\
873747.49499	0.00000\\
875751.50301	0.00000\\
877755.51102	0.00000\\
879759.51904	0.00000\\
881763.52705	0.00000\\
883767.53507	0.00000\\
885771.54309	0.00000\\
887775.55110	0.00000\\
889779.55912	0.00000\\
891783.56713	0.00000\\
893787.57515	0.00000\\
895791.58317	0.00000\\
897795.59118	0.00000\\
899799.59920	0.00000\\
901803.60721	0.00000\\
903807.61523	0.00000\\
905811.62325	0.00000\\
907815.63126	0.00000\\
909819.63928	0.00000\\
911823.64729	0.00000\\
913827.65531	0.00000\\
915831.66333	0.00000\\
917835.67134	0.00000\\
919839.67936	0.00000\\
921843.68737	0.00000\\
923847.69539	0.00000\\
925851.70341	0.00000\\
927855.71142	0.00000\\
929859.71944	0.00000\\
931863.72745	0.00000\\
933867.73547	0.00000\\
935871.74349	0.00000\\
937875.75150	0.00000\\
939879.75952	0.00000\\
941883.76754	0.00000\\
943887.77555	0.00000\\
945891.78357	0.00000\\
947895.79158	0.00000\\
949899.79960	0.00000\\
951903.80762	0.00000\\
953907.81563	0.00000\\
955911.82365	0.00000\\
957915.83166	0.00000\\
959919.83968	0.00000\\
961923.84770	0.00000\\
963927.85571	0.00000\\
965931.86373	0.00000\\
967935.87174	0.00000\\
969939.87976	0.00000\\
971943.88778	0.00000\\
973947.89579	0.00000\\
975951.90381	0.00000\\
977955.91182	0.00000\\
979959.91984	0.00000\\
981963.92786	0.00000\\
983967.93587	0.00000\\
985971.94389	0.00000\\
987975.95190	0.00000\\
989979.95992	0.00000\\
991983.96794	0.00000\\
993987.97595	0.00000\\
995991.98397	0.00000\\
997995.99198	0.00000\\
1000000.00000	0.00000\\
};
\addlegendentry{Basis-\ac{CCD} (90\,\% \ac{CI})}

\addplot [color=gray, dashed, line width=1.0pt, forget plot]
  table[row sep=crcr]{%
0.00000	1.00000\\
2004.00802	0.99951\\
4008.01603	0.99845\\
6012.02405	0.99694\\
8016.03206	0.99504\\
10020.04008	0.99278\\
12024.04810	0.99018\\
14028.05611	0.98725\\
16032.06413	0.98401\\
18036.07214	0.98047\\
20040.08016	0.97663\\
22044.08818	0.97252\\
24048.09619	0.96812\\
26052.10421	0.96345\\
28056.11222	0.95852\\
30060.12024	0.95333\\
32064.12826	0.94788\\
34068.13627	0.94219\\
36072.14429	0.93625\\
38076.15230	0.93007\\
40080.16032	0.92366\\
42084.16834	0.91702\\
44088.17635	0.91015\\
46092.18437	0.90305\\
48096.19238	0.89575\\
50100.20040	0.88822\\
52104.20842	0.88049\\
54108.21643	0.87256\\
56112.22445	0.86443\\
58116.23246	0.85610\\
60120.24048	0.84759\\
62124.24850	0.83889\\
64128.25651	0.83001\\
66132.26453	0.82095\\
68136.27255	0.81173\\
70140.28056	0.80234\\
72144.28858	0.79279\\
74148.29659	0.78309\\
76152.30461	0.77325\\
78156.31263	0.76326\\
80160.32064	0.75313\\
82164.32866	0.74287\\
84168.33667	0.73249\\
86172.34469	0.72199\\
88176.35271	0.71137\\
90180.36072	0.70065\\
92184.36874	0.68984\\
94188.37675	0.67892\\
96192.38477	0.66793\\
98196.39279	0.65685\\
100200.40080	0.64570\\
102204.40882	0.63448\\
104208.41683	0.62320\\
106212.42485	0.61188\\
108216.43287	0.60050\\
110220.44088	0.58909\\
112224.44890	0.57765\\
114228.45691	0.56619\\
116232.46493	0.55471\\
118236.47295	0.54322\\
120240.48096	0.53173\\
122244.48898	0.52025\\
124248.49699	0.50878\\
126252.50501	0.49733\\
128256.51303	0.48591\\
130260.52104	0.47452\\
132264.52906	0.46318\\
134268.53707	0.45188\\
136272.54509	0.44064\\
138276.55311	0.42947\\
140280.56112	0.41837\\
142284.56914	0.40734\\
144288.57715	0.39640\\
146292.58517	0.38555\\
148296.59319	0.37480\\
150300.60120	0.36415\\
152304.60922	0.35361\\
154308.61723	0.34319\\
156312.62525	0.33289\\
158316.63327	0.32272\\
160320.64128	0.31267\\
162324.64930	0.30277\\
164328.65731	0.29301\\
166332.66533	0.28340\\
168336.67335	0.27393\\
170340.68136	0.26463\\
172344.68938	0.25548\\
174348.69739	0.24650\\
176352.70541	0.23769\\
178356.71343	0.22905\\
180360.72144	0.22058\\
182364.72946	0.21229\\
184368.73747	0.20418\\
186372.74549	0.19624\\
188376.75351	0.18850\\
190380.76152	0.18093\\
192384.76954	0.17356\\
194388.77756	0.16637\\
196392.78557	0.15936\\
198396.79359	0.15255\\
200400.80160	0.14592\\
202404.80962	0.13949\\
204408.81764	0.13324\\
206412.82565	0.12718\\
208416.83367	0.12131\\
210420.84168	0.11562\\
212424.84970	0.11012\\
214428.85772	0.10480\\
216432.86573	0.09966\\
218436.87375	0.09470\\
220440.88176	0.08992\\
222444.88978	0.08531\\
224448.89780	0.08088\\
226452.90581	0.07661\\
228456.91383	0.07251\\
230460.92184	0.06858\\
232464.92986	0.06481\\
234468.93788	0.06119\\
236472.94589	0.05773\\
238476.95391	0.05442\\
240480.96192	0.05125\\
242484.96994	0.04823\\
244488.97796	0.04535\\
246492.98597	0.04260\\
248496.99399	0.03999\\
250501.00200	0.03750\\
252505.01002	0.03513\\
254509.01804	0.03289\\
256513.02605	0.03076\\
258517.03407	0.02874\\
260521.04208	0.02684\\
262525.05010	0.02503\\
264529.05812	0.02332\\
266533.06613	0.02172\\
268537.07415	0.02020\\
270541.08216	0.01877\\
272545.09018	0.01742\\
274549.09820	0.01616\\
276553.10621	0.01497\\
278557.11423	0.01386\\
280561.12224	0.01282\\
282565.13026	0.01184\\
284569.13828	0.01093\\
286573.14629	0.01007\\
288577.15431	0.00928\\
290581.16232	0.00854\\
292585.17034	0.00785\\
294589.17836	0.00721\\
296593.18637	0.00661\\
298597.19439	0.00606\\
300601.20240	0.00554\\
302605.21042	0.00507\\
304609.21844	0.00463\\
306613.22645	0.00422\\
308617.23447	0.00385\\
310621.24248	0.00350\\
312625.25050	0.00319\\
314629.25852	0.00290\\
316633.26653	0.00263\\
318637.27455	0.00238\\
320641.28257	0.00216\\
322645.29058	0.00195\\
324649.29860	0.00176\\
326653.30661	0.00159\\
328657.31463	0.00143\\
330661.32265	0.00129\\
332665.33066	0.00116\\
334669.33868	0.00104\\
336673.34669	0.00093\\
338677.35471	0.00084\\
340681.36273	0.00075\\
342685.37074	0.00067\\
344689.37876	0.00060\\
346693.38677	0.00053\\
348697.39479	0.00047\\
350701.40281	0.00042\\
352705.41082	0.00037\\
354709.41884	0.00033\\
356713.42685	0.00029\\
358717.43487	0.00026\\
360721.44289	0.00023\\
362725.45090	0.00020\\
364729.45892	0.00018\\
366733.46693	0.00016\\
368737.47495	0.00014\\
370741.48297	0.00012\\
372745.49098	0.00011\\
374749.49900	0.00009\\
376753.50701	0.00008\\
378757.51503	0.00007\\
380761.52305	0.00006\\
382765.53106	0.00005\\
384769.53908	0.00005\\
386773.54709	0.00004\\
388777.55511	0.00004\\
390781.56313	0.00003\\
392785.57114	0.00003\\
394789.57916	0.00002\\
396793.58717	0.00002\\
398797.59519	0.00002\\
400801.60321	0.00001\\
402805.61122	0.00001\\
404809.61924	0.00001\\
406813.62725	0.00001\\
408817.63527	0.00001\\
410821.64329	0.00001\\
412825.65130	0.00001\\
414829.65932	0.00000\\
416833.66733	0.00000\\
418837.67535	0.00000\\
420841.68337	0.00000\\
422845.69138	0.00000\\
424849.69940	0.00000\\
426853.70741	0.00000\\
428857.71543	0.00000\\
430861.72345	0.00000\\
432865.73146	0.00000\\
434869.73948	0.00000\\
436873.74749	0.00000\\
438877.75551	0.00000\\
440881.76353	0.00000\\
442885.77154	0.00000\\
444889.77956	0.00000\\
446893.78758	0.00000\\
448897.79559	0.00000\\
450901.80361	0.00000\\
452905.81162	0.00000\\
454909.81964	0.00000\\
456913.82766	0.00000\\
458917.83567	0.00000\\
460921.84369	0.00000\\
462925.85170	0.00000\\
464929.85972	0.00000\\
466933.86774	0.00000\\
468937.87575	0.00000\\
470941.88377	0.00000\\
472945.89178	0.00000\\
474949.89980	0.00000\\
476953.90782	0.00000\\
478957.91583	0.00000\\
480961.92385	0.00000\\
482965.93186	0.00000\\
484969.93988	0.00000\\
486973.94790	0.00000\\
488977.95591	0.00000\\
490981.96393	0.00000\\
492985.97194	0.00000\\
494989.97996	0.00000\\
496993.98798	0.00000\\
498997.99599	0.00000\\
501002.00401	0.00000\\
503006.01202	0.00000\\
505010.02004	0.00000\\
507014.02806	0.00000\\
509018.03607	0.00000\\
511022.04409	0.00000\\
513026.05210	0.00000\\
515030.06012	0.00000\\
517034.06814	0.00000\\
519038.07615	0.00000\\
521042.08417	0.00000\\
523046.09218	0.00000\\
525050.10020	0.00000\\
527054.10822	0.00000\\
529058.11623	0.00000\\
531062.12425	0.00000\\
533066.13226	0.00000\\
535070.14028	0.00000\\
537074.14830	0.00000\\
539078.15631	0.00000\\
541082.16433	0.00000\\
543086.17234	0.00000\\
545090.18036	0.00000\\
547094.18838	0.00000\\
549098.19639	0.00000\\
551102.20441	0.00000\\
553106.21242	0.00000\\
555110.22044	0.00000\\
557114.22846	0.00000\\
559118.23647	0.00000\\
561122.24449	0.00000\\
563126.25251	0.00000\\
565130.26052	0.00000\\
567134.26854	0.00000\\
569138.27655	0.00000\\
571142.28457	0.00000\\
573146.29259	0.00000\\
575150.30060	0.00000\\
577154.30862	0.00000\\
579158.31663	0.00000\\
581162.32465	0.00000\\
583166.33267	0.00000\\
585170.34068	0.00000\\
587174.34870	0.00000\\
589178.35671	0.00000\\
591182.36473	0.00000\\
593186.37275	0.00000\\
595190.38076	0.00000\\
597194.38878	0.00000\\
599198.39679	0.00000\\
601202.40481	0.00000\\
603206.41283	0.00000\\
605210.42084	0.00000\\
607214.42886	0.00000\\
609218.43687	0.00000\\
611222.44489	0.00000\\
613226.45291	0.00000\\
615230.46092	0.00000\\
617234.46894	0.00000\\
619238.47695	0.00000\\
621242.48497	0.00000\\
623246.49299	0.00000\\
625250.50100	0.00000\\
627254.50902	0.00000\\
629258.51703	0.00000\\
631262.52505	0.00000\\
633266.53307	0.00000\\
635270.54108	0.00000\\
637274.54910	0.00000\\
639278.55711	0.00000\\
641282.56513	0.00000\\
643286.57315	0.00000\\
645290.58116	0.00000\\
647294.58918	0.00000\\
649298.59719	0.00000\\
651302.60521	0.00000\\
653306.61323	0.00000\\
655310.62124	0.00000\\
657314.62926	0.00000\\
659318.63727	0.00000\\
661322.64529	0.00000\\
663326.65331	0.00000\\
665330.66132	0.00000\\
667334.66934	0.00000\\
669338.67735	0.00000\\
671342.68537	0.00000\\
673346.69339	0.00000\\
675350.70140	0.00000\\
677354.70942	0.00000\\
679358.71743	0.00000\\
681362.72545	0.00000\\
683366.73347	0.00000\\
685370.74148	0.00000\\
687374.74950	0.00000\\
689378.75752	0.00000\\
691382.76553	0.00000\\
693386.77355	0.00000\\
695390.78156	0.00000\\
697394.78958	0.00000\\
699398.79760	0.00000\\
701402.80561	0.00000\\
703406.81363	0.00000\\
705410.82164	0.00000\\
707414.82966	0.00000\\
709418.83768	0.00000\\
711422.84569	0.00000\\
713426.85371	0.00000\\
715430.86172	0.00000\\
717434.86974	0.00000\\
719438.87776	0.00000\\
721442.88577	0.00000\\
723446.89379	0.00000\\
725450.90180	0.00000\\
727454.90982	0.00000\\
729458.91784	0.00000\\
731462.92585	0.00000\\
733466.93387	0.00000\\
735470.94188	0.00000\\
737474.94990	0.00000\\
739478.95792	0.00000\\
741482.96593	0.00000\\
743486.97395	0.00000\\
745490.98196	0.00000\\
747494.98998	0.00000\\
749498.99800	0.00000\\
751503.00601	0.00000\\
753507.01403	0.00000\\
755511.02204	0.00000\\
757515.03006	0.00000\\
759519.03808	0.00000\\
761523.04609	0.00000\\
763527.05411	0.00000\\
765531.06212	0.00000\\
767535.07014	0.00000\\
769539.07816	0.00000\\
771543.08617	0.00000\\
773547.09419	0.00000\\
775551.10220	0.00000\\
777555.11022	0.00000\\
779559.11824	0.00000\\
781563.12625	0.00000\\
783567.13427	0.00000\\
785571.14228	0.00000\\
787575.15030	0.00000\\
789579.15832	0.00000\\
791583.16633	0.00000\\
793587.17435	0.00000\\
795591.18236	0.00000\\
797595.19038	0.00000\\
799599.19840	0.00000\\
801603.20641	0.00000\\
803607.21443	0.00000\\
805611.22244	0.00000\\
807615.23046	0.00000\\
809619.23848	0.00000\\
811623.24649	0.00000\\
813627.25451	0.00000\\
815631.26253	0.00000\\
817635.27054	0.00000\\
819639.27856	0.00000\\
821643.28657	0.00000\\
823647.29459	0.00000\\
825651.30261	0.00000\\
827655.31062	0.00000\\
829659.31864	0.00000\\
831663.32665	0.00000\\
833667.33467	0.00000\\
835671.34269	0.00000\\
837675.35070	0.00000\\
839679.35872	0.00000\\
841683.36673	0.00000\\
843687.37475	0.00000\\
845691.38277	0.00000\\
847695.39078	0.00000\\
849699.39880	0.00000\\
851703.40681	0.00000\\
853707.41483	0.00000\\
855711.42285	0.00000\\
857715.43086	0.00000\\
859719.43888	0.00000\\
861723.44689	0.00000\\
863727.45491	0.00000\\
865731.46293	0.00000\\
867735.47094	0.00000\\
869739.47896	0.00000\\
871743.48697	0.00000\\
873747.49499	0.00000\\
875751.50301	0.00000\\
877755.51102	0.00000\\
879759.51904	0.00000\\
881763.52705	0.00000\\
883767.53507	0.00000\\
885771.54309	0.00000\\
887775.55110	0.00000\\
889779.55912	0.00000\\
891783.56713	0.00000\\
893787.57515	0.00000\\
895791.58317	0.00000\\
897795.59118	0.00000\\
899799.59920	0.00000\\
901803.60721	0.00000\\
903807.61523	0.00000\\
905811.62325	0.00000\\
907815.63126	0.00000\\
909819.63928	0.00000\\
911823.64729	0.00000\\
913827.65531	0.00000\\
915831.66333	0.00000\\
917835.67134	0.00000\\
919839.67936	0.00000\\
921843.68737	0.00000\\
923847.69539	0.00000\\
925851.70341	0.00000\\
927855.71142	0.00000\\
929859.71944	0.00000\\
931863.72745	0.00000\\
933867.73547	0.00000\\
935871.74349	0.00000\\
937875.75150	0.00000\\
939879.75952	0.00000\\
941883.76754	0.00000\\
943887.77555	0.00000\\
945891.78357	0.00000\\
947895.79158	0.00000\\
949899.79960	0.00000\\
951903.80762	0.00000\\
953907.81563	0.00000\\
955911.82365	0.00000\\
957915.83166	0.00000\\
959919.83968	0.00000\\
961923.84770	0.00000\\
963927.85571	0.00000\\
965931.86373	0.00000\\
967935.87174	0.00000\\
969939.87976	0.00000\\
971943.88778	0.00000\\
973947.89579	0.00000\\
975951.90381	0.00000\\
977955.91182	0.00000\\
979959.91984	0.00000\\
981963.92786	0.00000\\
983967.93587	0.00000\\
985971.94389	0.00000\\
987975.95190	0.00000\\
989979.95992	0.00000\\
991983.96794	0.00000\\
993987.97595	0.00000\\
995991.98397	0.00000\\
997995.99198	0.00000\\
1000000.00000	0.00000\\
};
\addplot [color=gray, dashed, line width=1.0pt, forget plot]
  table[row sep=crcr]{%
0.00000	1.00000\\
2004.00802	1.00000\\
4008.01603	1.00000\\
6012.02405	0.99999\\
8016.03206	0.99998\\
10020.04008	0.99995\\
12024.04810	0.99992\\
14028.05611	0.99987\\
16032.06413	0.99980\\
18036.07214	0.99971\\
20040.08016	0.99960\\
22044.08818	0.99946\\
24048.09619	0.99930\\
26052.10421	0.99911\\
28056.11222	0.99888\\
30060.12024	0.99862\\
32064.12826	0.99832\\
34068.13627	0.99798\\
36072.14429	0.99760\\
38076.15230	0.99718\\
40080.16032	0.99671\\
42084.16834	0.99619\\
44088.17635	0.99562\\
46092.18437	0.99500\\
48096.19238	0.99433\\
50100.20040	0.99359\\
52104.20842	0.99280\\
54108.21643	0.99195\\
56112.22445	0.99104\\
58116.23246	0.99007\\
60120.24048	0.98903\\
62124.24850	0.98792\\
64128.25651	0.98675\\
66132.26453	0.98551\\
68136.27255	0.98420\\
70140.28056	0.98282\\
72144.28858	0.98136\\
74148.29659	0.97984\\
76152.30461	0.97823\\
78156.31263	0.97656\\
80160.32064	0.97481\\
82164.32866	0.97298\\
84168.33667	0.97107\\
86172.34469	0.96909\\
88176.35271	0.96702\\
90180.36072	0.96488\\
92184.36874	0.96266\\
94188.37675	0.96036\\
96192.38477	0.95799\\
98196.39279	0.95553\\
100200.40080	0.95299\\
102204.40882	0.95037\\
104208.41683	0.94767\\
106212.42485	0.94490\\
108216.43287	0.94204\\
110220.44088	0.93910\\
112224.44890	0.93609\\
114228.45691	0.93300\\
116232.46493	0.92983\\
118236.47295	0.92658\\
120240.48096	0.92326\\
122244.48898	0.91986\\
124248.49699	0.91638\\
126252.50501	0.91283\\
128256.51303	0.90920\\
130260.52104	0.90551\\
132264.52906	0.90174\\
134268.53707	0.89789\\
136272.54509	0.89398\\
138276.55311	0.89000\\
140280.56112	0.88595\\
142284.56914	0.88183\\
144288.57715	0.87765\\
146292.58517	0.87340\\
148296.59319	0.86909\\
150300.60120	0.86471\\
152304.60922	0.86028\\
154308.61723	0.85578\\
156312.62525	0.85123\\
158316.63327	0.84661\\
160320.64128	0.84194\\
162324.64930	0.83722\\
164328.65731	0.83244\\
166332.66533	0.82761\\
168336.67335	0.82273\\
170340.68136	0.81780\\
172344.68938	0.81282\\
174348.69739	0.80780\\
176352.70541	0.80273\\
178356.71343	0.79762\\
180360.72144	0.79246\\
182364.72946	0.78727\\
184368.73747	0.78203\\
186372.74549	0.77676\\
188376.75351	0.77145\\
190380.76152	0.76611\\
192384.76954	0.76073\\
194388.77756	0.75532\\
196392.78557	0.74988\\
198396.79359	0.74442\\
200400.80160	0.73892\\
202404.80962	0.73340\\
204408.81764	0.72786\\
206412.82565	0.72229\\
208416.83367	0.71670\\
210420.84168	0.71109\\
212424.84970	0.70546\\
214428.85772	0.69981\\
216432.86573	0.69415\\
218436.87375	0.68847\\
220440.88176	0.68278\\
222444.88978	0.67708\\
224448.89780	0.67136\\
226452.90581	0.66564\\
228456.91383	0.65991\\
230460.92184	0.65417\\
232464.92986	0.64843\\
234468.93788	0.64268\\
236472.94589	0.63693\\
238476.95391	0.63118\\
240480.96192	0.62542\\
242484.96994	0.61967\\
244488.97796	0.61392\\
246492.98597	0.60817\\
248496.99399	0.60243\\
250501.00200	0.59669\\
252505.01002	0.59096\\
254509.01804	0.58523\\
256513.02605	0.57952\\
258517.03407	0.57381\\
260521.04208	0.56811\\
262525.05010	0.56243\\
264529.05812	0.55676\\
266533.06613	0.55110\\
268537.07415	0.54545\\
270541.08216	0.53982\\
272545.09018	0.53421\\
274549.09820	0.52861\\
276553.10621	0.52303\\
278557.11423	0.51747\\
280561.12224	0.51193\\
282565.13026	0.50641\\
284569.13828	0.50091\\
286573.14629	0.49543\\
288577.15431	0.48997\\
290581.16232	0.48454\\
292585.17034	0.47913\\
294589.17836	0.47375\\
296593.18637	0.46839\\
298597.19439	0.46306\\
300601.20240	0.45775\\
302605.21042	0.45248\\
304609.21844	0.44722\\
306613.22645	0.44200\\
308617.23447	0.43681\\
310621.24248	0.43164\\
312625.25050	0.42651\\
314629.25852	0.42141\\
316633.26653	0.41633\\
318637.27455	0.41129\\
320641.28257	0.40628\\
322645.29058	0.40130\\
324649.29860	0.39636\\
326653.30661	0.39145\\
328657.31463	0.38657\\
330661.32265	0.38172\\
332665.33066	0.37691\\
334669.33868	0.37214\\
336673.34669	0.36740\\
338677.35471	0.36269\\
340681.36273	0.35802\\
342685.37074	0.35338\\
344689.37876	0.34878\\
346693.38677	0.34422\\
348697.39479	0.33969\\
350701.40281	0.33520\\
352705.41082	0.33075\\
354709.41884	0.32633\\
356713.42685	0.32195\\
358717.43487	0.31761\\
360721.44289	0.31330\\
362725.45090	0.30904\\
364729.45892	0.30481\\
366733.46693	0.30061\\
368737.47495	0.29646\\
370741.48297	0.29234\\
372745.49098	0.28827\\
374749.49900	0.28423\\
376753.50701	0.28022\\
378757.51503	0.27626\\
380761.52305	0.27233\\
382765.53106	0.26845\\
384769.53908	0.26460\\
386773.54709	0.26078\\
388777.55511	0.25701\\
390781.56313	0.25328\\
392785.57114	0.24958\\
394789.57916	0.24592\\
396793.58717	0.24230\\
398797.59519	0.23872\\
400801.60321	0.23517\\
402805.61122	0.23166\\
404809.61924	0.22820\\
406813.62725	0.22476\\
408817.63527	0.22137\\
410821.64329	0.21801\\
412825.65130	0.21469\\
414829.65932	0.21141\\
416833.66733	0.20817\\
418837.67535	0.20496\\
420841.68337	0.20179\\
422845.69138	0.19865\\
424849.69940	0.19556\\
426853.70741	0.19249\\
428857.71543	0.18947\\
430861.72345	0.18648\\
432865.73146	0.18352\\
434869.73948	0.18061\\
436873.74749	0.17772\\
438877.75551	0.17488\\
440881.76353	0.17207\\
442885.77154	0.16929\\
444889.77956	0.16655\\
446893.78758	0.16384\\
448897.79559	0.16116\\
450901.80361	0.15852\\
452905.81162	0.15592\\
454909.81964	0.15335\\
456913.82766	0.15081\\
458917.83567	0.14830\\
460921.84369	0.14583\\
462925.85170	0.14339\\
464929.85972	0.14098\\
466933.86774	0.13860\\
468937.87575	0.13626\\
470941.88377	0.13395\\
472945.89178	0.13167\\
474949.89980	0.12942\\
476953.90782	0.12720\\
478957.91583	0.12501\\
480961.92385	0.12285\\
482965.93186	0.12073\\
484969.93988	0.11863\\
486973.94790	0.11656\\
488977.95591	0.11452\\
490981.96393	0.11251\\
492985.97194	0.11053\\
494989.97996	0.10858\\
496993.98798	0.10666\\
498997.99599	0.10476\\
501002.00401	0.10289\\
503006.01202	0.10105\\
505010.02004	0.09924\\
507014.02806	0.09745\\
509018.03607	0.09569\\
511022.04409	0.09396\\
513026.05210	0.09225\\
515030.06012	0.09057\\
517034.06814	0.08892\\
519038.07615	0.08729\\
521042.08417	0.08568\\
523046.09218	0.08410\\
525050.10020	0.08255\\
527054.10822	0.08101\\
529058.11623	0.07951\\
531062.12425	0.07802\\
533066.13226	0.07656\\
535070.14028	0.07512\\
537074.14830	0.07371\\
539078.15631	0.07232\\
541082.16433	0.07095\\
543086.17234	0.06960\\
545090.18036	0.06827\\
547094.18838	0.06697\\
549098.19639	0.06568\\
551102.20441	0.06442\\
553106.21242	0.06318\\
555110.22044	0.06196\\
557114.22846	0.06076\\
559118.23647	0.05958\\
561122.24449	0.05842\\
563126.25251	0.05728\\
565130.26052	0.05615\\
567134.26854	0.05505\\
569138.27655	0.05397\\
571142.28457	0.05290\\
573146.29259	0.05185\\
575150.30060	0.05082\\
577154.30862	0.04981\\
579158.31663	0.04882\\
581162.32465	0.04784\\
583166.33267	0.04688\\
585170.34068	0.04593\\
587174.34870	0.04501\\
589178.35671	0.04410\\
591182.36473	0.04320\\
593186.37275	0.04232\\
595190.38076	0.04146\\
597194.38878	0.04061\\
599198.39679	0.03978\\
601202.40481	0.03896\\
603206.41283	0.03816\\
605210.42084	0.03737\\
607214.42886	0.03660\\
609218.43687	0.03584\\
611222.44489	0.03509\\
613226.45291	0.03436\\
615230.46092	0.03364\\
617234.46894	0.03294\\
619238.47695	0.03224\\
621242.48497	0.03157\\
623246.49299	0.03090\\
625250.50100	0.03024\\
627254.50902	0.02960\\
629258.51703	0.02897\\
631262.52505	0.02835\\
633266.53307	0.02775\\
635270.54108	0.02715\\
637274.54910	0.02657\\
639278.55711	0.02600\\
641282.56513	0.02544\\
643286.57315	0.02489\\
645290.58116	0.02435\\
647294.58918	0.02382\\
649298.59719	0.02330\\
651302.60521	0.02279\\
653306.61323	0.02229\\
655310.62124	0.02180\\
657314.62926	0.02132\\
659318.63727	0.02085\\
661322.64529	0.02039\\
663326.65331	0.01994\\
665330.66132	0.01950\\
667334.66934	0.01907\\
669338.67735	0.01864\\
671342.68537	0.01822\\
673346.69339	0.01782\\
675350.70140	0.01742\\
677354.70942	0.01702\\
679358.71743	0.01664\\
681362.72545	0.01627\\
683366.73347	0.01590\\
685370.74148	0.01554\\
687374.74950	0.01518\\
689378.75752	0.01484\\
691382.76553	0.01450\\
693386.77355	0.01417\\
695390.78156	0.01384\\
697394.78958	0.01352\\
699398.79760	0.01321\\
701402.80561	0.01291\\
703406.81363	0.01261\\
705410.82164	0.01232\\
707414.82966	0.01203\\
709418.83768	0.01175\\
711422.84569	0.01148\\
713426.85371	0.01121\\
715430.86172	0.01095\\
717434.86974	0.01069\\
719438.87776	0.01044\\
721442.88577	0.01019\\
723446.89379	0.00995\\
725450.90180	0.00971\\
727454.90982	0.00948\\
729458.91784	0.00926\\
731462.92585	0.00904\\
733466.93387	0.00882\\
735470.94188	0.00861\\
737474.94990	0.00841\\
739478.95792	0.00820\\
741482.96593	0.00801\\
743486.97395	0.00781\\
745490.98196	0.00762\\
747494.98998	0.00744\\
749498.99800	0.00726\\
751503.00601	0.00708\\
753507.01403	0.00691\\
755511.02204	0.00674\\
757515.03006	0.00658\\
759519.03808	0.00642\\
761523.04609	0.00626\\
763527.05411	0.00611\\
765531.06212	0.00596\\
767535.07014	0.00581\\
769539.07816	0.00566\\
771543.08617	0.00552\\
773547.09419	0.00539\\
775551.10220	0.00525\\
777555.11022	0.00512\\
779559.11824	0.00499\\
781563.12625	0.00487\\
783567.13427	0.00475\\
785571.14228	0.00463\\
787575.15030	0.00451\\
789579.15832	0.00440\\
791583.16633	0.00429\\
793587.17435	0.00418\\
795591.18236	0.00407\\
797595.19038	0.00397\\
799599.19840	0.00387\\
801603.20641	0.00377\\
803607.21443	0.00367\\
805611.22244	0.00358\\
807615.23046	0.00349\\
809619.23848	0.00340\\
811623.24649	0.00331\\
813627.25451	0.00323\\
815631.26253	0.00314\\
817635.27054	0.00306\\
819639.27856	0.00298\\
821643.28657	0.00291\\
823647.29459	0.00283\\
825651.30261	0.00276\\
827655.31062	0.00268\\
829659.31864	0.00261\\
831663.32665	0.00255\\
833667.33467	0.00248\\
835671.34269	0.00241\\
837675.35070	0.00235\\
839679.35872	0.00229\\
841683.36673	0.00223\\
843687.37475	0.00217\\
845691.38277	0.00211\\
847695.39078	0.00206\\
849699.39880	0.00200\\
851703.40681	0.00195\\
853707.41483	0.00190\\
855711.42285	0.00184\\
857715.43086	0.00180\\
859719.43888	0.00175\\
861723.44689	0.00170\\
863727.45491	0.00165\\
865731.46293	0.00161\\
867735.47094	0.00157\\
869739.47896	0.00152\\
871743.48697	0.00148\\
873747.49499	0.00144\\
875751.50301	0.00140\\
877755.51102	0.00137\\
879759.51904	0.00133\\
881763.52705	0.00129\\
883767.53507	0.00126\\
885771.54309	0.00122\\
887775.55110	0.00119\\
889779.55912	0.00116\\
891783.56713	0.00112\\
893787.57515	0.00109\\
895791.58317	0.00106\\
897795.59118	0.00103\\
899799.59920	0.00100\\
901803.60721	0.00098\\
903807.61523	0.00095\\
905811.62325	0.00092\\
907815.63126	0.00090\\
909819.63928	0.00087\\
911823.64729	0.00085\\
913827.65531	0.00082\\
915831.66333	0.00080\\
917835.67134	0.00078\\
919839.67936	0.00076\\
921843.68737	0.00074\\
923847.69539	0.00071\\
925851.70341	0.00069\\
927855.71142	0.00067\\
929859.71944	0.00066\\
931863.72745	0.00064\\
933867.73547	0.00062\\
935871.74349	0.00060\\
937875.75150	0.00058\\
939879.75952	0.00057\\
941883.76754	0.00055\\
943887.77555	0.00054\\
945891.78357	0.00052\\
947895.79158	0.00051\\
949899.79960	0.00049\\
951903.80762	0.00048\\
953907.81563	0.00046\\
955911.82365	0.00045\\
957915.83166	0.00044\\
959919.83968	0.00042\\
961923.84770	0.00041\\
963927.85571	0.00040\\
965931.86373	0.00039\\
967935.87174	0.00038\\
969939.87976	0.00037\\
971943.88778	0.00035\\
973947.89579	0.00034\\
975951.90381	0.00033\\
977955.91182	0.00032\\
979959.91984	0.00031\\
981963.92786	0.00031\\
983967.93587	0.00030\\
985971.94389	0.00029\\
987975.95190	0.00028\\
989979.95992	0.00027\\
991983.96794	0.00026\\
993987.97595	0.00025\\
995991.98397	0.00025\\
997995.99198	0.00024\\
1000000.00000	0.00023\\
};

\addplot[area legend, draw=none, fill=white!30!black, fill opacity=0.35000, forget plot]
table[row sep=crcr] {%
x	y\\
0.00000	1.00000\\
2004.00802	0.99999\\
4008.01603	0.99995\\
6012.02405	0.99987\\
8016.03206	0.99977\\
10020.04008	0.99943\\
12024.04810	0.99847\\
14028.05611	0.99716\\
16032.06413	0.99538\\
18036.07214	0.99310\\
20040.08016	0.99027\\
22044.08818	0.98686\\
24048.09619	0.98285\\
26052.10421	0.97820\\
28056.11222	0.97291\\
30060.12024	0.96695\\
32064.12826	0.96032\\
34068.13627	0.95301\\
36072.14429	0.94502\\
38076.15230	0.93635\\
40080.16032	0.92700\\
42084.16834	0.91698\\
44088.17635	0.90630\\
46092.18437	0.89496\\
48096.19238	0.88300\\
50100.20040	0.87042\\
52104.20842	0.85725\\
54108.21643	0.84351\\
56112.22445	0.82922\\
58116.23246	0.81442\\
60120.24048	0.79912\\
62124.24850	0.78337\\
64128.25651	0.76719\\
66132.26453	0.75062\\
68136.27255	0.73368\\
70140.28056	0.71643\\
72144.28858	0.69889\\
74148.29659	0.68110\\
76152.30461	0.66310\\
78156.31263	0.64492\\
80160.32064	0.62661\\
82164.32866	0.60819\\
84168.33667	0.58972\\
86172.34469	0.57121\\
88176.35271	0.55271\\
90180.36072	0.53426\\
92184.36874	0.51589\\
94188.37675	0.49762\\
96192.38477	0.47950\\
98196.39279	0.46155\\
100200.40080	0.44381\\
102204.40882	0.42629\\
104208.41683	0.40903\\
106212.42485	0.39205\\
108216.43287	0.37537\\
110220.44088	0.35902\\
112224.44890	0.34301\\
114228.45691	0.32736\\
116232.46493	0.31209\\
118236.47295	0.29721\\
120240.48096	0.28273\\
122244.48898	0.26866\\
124248.49699	0.25502\\
126252.50501	0.24180\\
128256.51303	0.22902\\
130260.52104	0.21668\\
132264.52906	0.20478\\
134268.53707	0.19332\\
136272.54509	0.18230\\
138276.55311	0.17172\\
140280.56112	0.16158\\
142284.56914	0.15186\\
144288.57715	0.14258\\
146292.58517	0.13371\\
148296.59319	0.12526\\
150300.60120	0.11721\\
152304.60922	0.10956\\
154308.61723	0.10229\\
156312.62525	0.09540\\
158316.63327	0.08888\\
160320.64128	0.08271\\
162324.64930	0.07688\\
164328.65731	0.07138\\
166332.66533	0.06621\\
168336.67335	0.06134\\
170340.68136	0.05676\\
172344.68938	0.05247\\
174348.69739	0.04845\\
176352.70541	0.04469\\
178356.71343	0.04117\\
180360.72144	0.03789\\
182364.72946	0.03483\\
184368.73747	0.03198\\
186372.74549	0.02933\\
188376.75351	0.02687\\
190380.76152	0.02459\\
192384.76954	0.02248\\
194388.77756	0.02053\\
196392.78557	0.01873\\
198396.79359	0.01706\\
200400.80160	0.01553\\
202404.80962	0.01411\\
204408.81764	0.01282\\
206412.82565	0.01163\\
208416.83367	0.01053\\
210420.84168	0.00953\\
212424.84970	0.00862\\
214428.85772	0.00778\\
216432.86573	0.00702\\
218436.87375	0.00632\\
220440.88176	0.00569\\
222444.88978	0.00512\\
224448.89780	0.00460\\
226452.90581	0.00412\\
228456.91383	0.00369\\
230460.92184	0.00331\\
232464.92986	0.00296\\
234468.93788	0.00264\\
236472.94589	0.00235\\
238476.95391	0.00210\\
240480.96192	0.00187\\
242484.96994	0.00166\\
244488.97796	0.00147\\
246492.98597	0.00131\\
248496.99399	0.00116\\
250501.00200	0.00103\\
252505.01002	0.00091\\
254509.01804	0.00080\\
256513.02605	0.00071\\
258517.03407	0.00062\\
260521.04208	0.00055\\
262525.05010	0.00048\\
264529.05812	0.00042\\
266533.06613	0.00037\\
268537.07415	0.00033\\
270541.08216	0.00029\\
272545.09018	0.00025\\
274549.09820	0.00022\\
276553.10621	0.00019\\
278557.11423	0.00017\\
280561.12224	0.00014\\
282565.13026	0.00013\\
284569.13828	0.00011\\
286573.14629	0.00009\\
288577.15431	0.00008\\
290581.16232	0.00007\\
292585.17034	0.00006\\
294589.17836	0.00005\\
296593.18637	0.00005\\
298597.19439	0.00004\\
300601.20240	0.00003\\
302605.21042	0.00003\\
304609.21844	0.00003\\
306613.22645	0.00002\\
308617.23447	0.00002\\
310621.24248	0.00002\\
312625.25050	0.00001\\
314629.25852	0.00001\\
316633.26653	0.00001\\
318637.27455	0.00001\\
320641.28257	0.00001\\
322645.29058	0.00001\\
324649.29860	0.00001\\
326653.30661	0.00000\\
328657.31463	0.00000\\
330661.32265	0.00000\\
332665.33066	0.00000\\
334669.33868	0.00000\\
336673.34669	0.00000\\
338677.35471	0.00000\\
340681.36273	0.00000\\
342685.37074	0.00000\\
344689.37876	0.00000\\
346693.38677	0.00000\\
348697.39479	0.00000\\
350701.40281	0.00000\\
352705.41082	0.00000\\
354709.41884	0.00000\\
356713.42685	0.00000\\
358717.43487	0.00000\\
360721.44289	0.00000\\
362725.45090	0.00000\\
364729.45892	0.00000\\
366733.46693	0.00000\\
368737.47495	0.00000\\
370741.48297	0.00000\\
372745.49098	0.00000\\
374749.49900	0.00000\\
376753.50701	0.00000\\
378757.51503	0.00000\\
380761.52305	0.00000\\
382765.53106	0.00000\\
384769.53908	0.00000\\
386773.54709	0.00000\\
388777.55511	0.00000\\
390781.56313	0.00000\\
392785.57114	0.00000\\
394789.57916	0.00000\\
396793.58717	0.00000\\
398797.59519	0.00000\\
400801.60321	0.00000\\
402805.61122	0.00000\\
404809.61924	0.00000\\
406813.62725	0.00000\\
408817.63527	0.00000\\
410821.64329	0.00000\\
412825.65130	0.00000\\
414829.65932	0.00000\\
416833.66733	0.00000\\
418837.67535	0.00000\\
420841.68337	0.00000\\
422845.69138	0.00000\\
424849.69940	0.00000\\
426853.70741	0.00000\\
428857.71543	0.00000\\
430861.72345	0.00000\\
432865.73146	0.00000\\
434869.73948	0.00000\\
436873.74749	0.00000\\
438877.75551	0.00000\\
440881.76353	0.00000\\
442885.77154	0.00000\\
444889.77956	0.00000\\
446893.78758	0.00000\\
448897.79559	0.00000\\
450901.80361	0.00000\\
452905.81162	0.00000\\
454909.81964	0.00000\\
456913.82766	0.00000\\
458917.83567	0.00000\\
460921.84369	0.00000\\
462925.85170	0.00000\\
464929.85972	0.00000\\
466933.86774	0.00000\\
468937.87575	0.00000\\
470941.88377	0.00000\\
472945.89178	0.00000\\
474949.89980	0.00000\\
476953.90782	0.00000\\
478957.91583	0.00000\\
480961.92385	0.00000\\
482965.93186	0.00000\\
484969.93988	0.00000\\
486973.94790	0.00000\\
488977.95591	0.00000\\
490981.96393	0.00000\\
492985.97194	0.00000\\
494989.97996	0.00000\\
496993.98798	0.00000\\
498997.99599	0.00000\\
501002.00401	0.00000\\
503006.01202	0.00000\\
505010.02004	0.00000\\
507014.02806	0.00000\\
509018.03607	0.00000\\
511022.04409	0.00000\\
513026.05210	0.00000\\
515030.06012	0.00000\\
517034.06814	0.00000\\
519038.07615	0.00000\\
521042.08417	0.00000\\
523046.09218	0.00000\\
525050.10020	0.00000\\
527054.10822	0.00000\\
529058.11623	0.00000\\
531062.12425	0.00000\\
533066.13226	0.00000\\
535070.14028	0.00000\\
537074.14830	0.00000\\
539078.15631	0.00000\\
541082.16433	0.00000\\
543086.17234	0.00000\\
545090.18036	0.00000\\
547094.18838	0.00000\\
549098.19639	0.00000\\
551102.20441	0.00000\\
553106.21242	0.00000\\
555110.22044	0.00000\\
557114.22846	0.00000\\
559118.23647	0.00000\\
561122.24449	0.00000\\
563126.25251	0.00000\\
565130.26052	0.00000\\
567134.26854	0.00000\\
569138.27655	0.00000\\
571142.28457	0.00000\\
573146.29259	0.00000\\
575150.30060	0.00000\\
577154.30862	0.00000\\
579158.31663	0.00000\\
581162.32465	0.00000\\
583166.33267	0.00000\\
585170.34068	0.00000\\
587174.34870	0.00000\\
589178.35671	0.00000\\
591182.36473	0.00000\\
593186.37275	0.00000\\
595190.38076	0.00000\\
597194.38878	0.00000\\
599198.39679	0.00000\\
601202.40481	0.00000\\
603206.41283	0.00000\\
605210.42084	0.00000\\
607214.42886	0.00000\\
609218.43687	0.00000\\
611222.44489	0.00000\\
613226.45291	0.00000\\
615230.46092	0.00000\\
617234.46894	0.00000\\
619238.47695	0.00000\\
621242.48497	0.00000\\
623246.49299	0.00000\\
625250.50100	0.00000\\
627254.50902	0.00000\\
629258.51703	0.00000\\
631262.52505	0.00000\\
633266.53307	0.00000\\
635270.54108	0.00000\\
637274.54910	0.00000\\
639278.55711	0.00000\\
641282.56513	0.00000\\
643286.57315	0.00000\\
645290.58116	0.00000\\
647294.58918	0.00000\\
649298.59719	0.00000\\
651302.60521	0.00000\\
653306.61323	0.00000\\
655310.62124	0.00000\\
657314.62926	0.00000\\
659318.63727	0.00000\\
661322.64529	0.00000\\
663326.65331	0.00000\\
665330.66132	0.00000\\
667334.66934	0.00000\\
669338.67735	0.00000\\
671342.68537	0.00000\\
673346.69339	0.00000\\
675350.70140	0.00000\\
677354.70942	0.00000\\
679358.71743	0.00000\\
681362.72545	0.00000\\
683366.73347	0.00000\\
685370.74148	0.00000\\
687374.74950	0.00000\\
689378.75752	0.00000\\
691382.76553	0.00000\\
693386.77355	0.00000\\
695390.78156	0.00000\\
697394.78958	0.00000\\
699398.79760	0.00000\\
701402.80561	0.00000\\
703406.81363	0.00000\\
705410.82164	0.00000\\
707414.82966	0.00000\\
709418.83768	0.00000\\
711422.84569	0.00000\\
713426.85371	0.00000\\
715430.86172	0.00000\\
717434.86974	0.00000\\
719438.87776	0.00000\\
721442.88577	0.00000\\
723446.89379	0.00000\\
725450.90180	0.00000\\
727454.90982	0.00000\\
729458.91784	0.00000\\
731462.92585	0.00000\\
733466.93387	0.00000\\
735470.94188	0.00000\\
737474.94990	0.00000\\
739478.95792	0.00000\\
741482.96593	0.00000\\
743486.97395	0.00000\\
745490.98196	0.00000\\
747494.98998	0.00000\\
749498.99800	0.00000\\
751503.00601	0.00000\\
753507.01403	0.00000\\
755511.02204	0.00000\\
757515.03006	0.00000\\
759519.03808	0.00000\\
761523.04609	0.00000\\
763527.05411	0.00000\\
765531.06212	0.00000\\
767535.07014	0.00000\\
769539.07816	0.00000\\
771543.08617	0.00000\\
773547.09419	0.00000\\
775551.10220	0.00000\\
777555.11022	0.00000\\
779559.11824	0.00000\\
781563.12625	0.00000\\
783567.13427	0.00000\\
785571.14228	0.00000\\
787575.15030	0.00000\\
789579.15832	0.00000\\
791583.16633	0.00000\\
793587.17435	0.00000\\
795591.18236	0.00000\\
797595.19038	0.00000\\
799599.19840	0.00000\\
801603.20641	0.00000\\
803607.21443	0.00000\\
805611.22244	0.00000\\
807615.23046	0.00000\\
809619.23848	0.00000\\
811623.24649	0.00000\\
813627.25451	0.00000\\
815631.26253	0.00000\\
817635.27054	0.00000\\
819639.27856	0.00000\\
821643.28657	0.00000\\
823647.29459	0.00000\\
825651.30261	0.00000\\
827655.31062	0.00000\\
829659.31864	0.00000\\
831663.32665	0.00000\\
833667.33467	0.00000\\
835671.34269	0.00000\\
837675.35070	0.00000\\
839679.35872	0.00000\\
841683.36673	0.00000\\
843687.37475	0.00000\\
845691.38277	0.00000\\
847695.39078	0.00000\\
849699.39880	0.00000\\
851703.40681	0.00000\\
853707.41483	0.00000\\
855711.42285	0.00000\\
857715.43086	0.00000\\
859719.43888	0.00000\\
861723.44689	0.00000\\
863727.45491	0.00000\\
865731.46293	0.00000\\
867735.47094	0.00000\\
869739.47896	0.00000\\
871743.48697	0.00000\\
873747.49499	0.00000\\
875751.50301	0.00000\\
877755.51102	0.00000\\
879759.51904	0.00000\\
881763.52705	0.00000\\
883767.53507	0.00000\\
885771.54309	0.00000\\
887775.55110	0.00000\\
889779.55912	0.00000\\
891783.56713	0.00000\\
893787.57515	0.00000\\
895791.58317	0.00000\\
897795.59118	0.00000\\
899799.59920	0.00000\\
901803.60721	0.00000\\
903807.61523	0.00000\\
905811.62325	0.00000\\
907815.63126	0.00000\\
909819.63928	0.00000\\
911823.64729	0.00000\\
913827.65531	0.00000\\
915831.66333	0.00000\\
917835.67134	0.00000\\
919839.67936	0.00000\\
921843.68737	0.00000\\
923847.69539	0.00000\\
925851.70341	0.00000\\
927855.71142	0.00000\\
929859.71944	0.00000\\
931863.72745	0.00000\\
933867.73547	0.00000\\
935871.74349	0.00000\\
937875.75150	0.00000\\
939879.75952	0.00000\\
941883.76754	0.00000\\
943887.77555	0.00000\\
945891.78357	0.00000\\
947895.79158	0.00000\\
949899.79960	0.00000\\
951903.80762	0.00000\\
953907.81563	0.00000\\
955911.82365	0.00000\\
957915.83166	0.00000\\
959919.83968	0.00000\\
961923.84770	0.00000\\
963927.85571	0.00000\\
965931.86373	0.00000\\
967935.87174	0.00000\\
969939.87976	0.00000\\
971943.88778	0.00000\\
973947.89579	0.00000\\
975951.90381	0.00000\\
977955.91182	0.00000\\
979959.91984	0.00000\\
981963.92786	0.00000\\
983967.93587	0.00000\\
985971.94389	0.00000\\
987975.95190	0.00000\\
989979.95992	0.00000\\
991983.96794	0.00000\\
993987.97595	0.00000\\
995991.98397	0.00000\\
997995.99198	0.00000\\
1000000.00000	0.00000\\
1000000.00000	0.72803\\
997995.99198	0.72910\\
995991.98397	0.73016\\
993987.97595	0.73123\\
991983.96794	0.73229\\
989979.95992	0.73335\\
987975.95190	0.73441\\
985971.94389	0.73547\\
983967.93587	0.73652\\
981963.92786	0.73758\\
979959.91984	0.73863\\
977955.91182	0.73968\\
975951.90381	0.74073\\
973947.89579	0.74178\\
971943.88778	0.74282\\
969939.87976	0.74387\\
967935.87174	0.74491\\
965931.86373	0.74595\\
963927.85571	0.74699\\
961923.84770	0.74803\\
959919.83968	0.74907\\
957915.83166	0.75011\\
955911.82365	0.75114\\
953907.81563	0.75217\\
951903.80762	0.75320\\
949899.79960	0.75423\\
947895.79158	0.75526\\
945891.78357	0.75629\\
943887.77555	0.75731\\
941883.76754	0.75833\\
939879.75952	0.75936\\
937875.75150	0.76038\\
935871.74349	0.76139\\
933867.73547	0.76241\\
931863.72745	0.76342\\
929859.71944	0.76444\\
927855.71142	0.76545\\
925851.70341	0.76646\\
923847.69539	0.76747\\
921843.68737	0.76847\\
919839.67936	0.76948\\
917835.67134	0.77048\\
915831.66333	0.77148\\
913827.65531	0.77248\\
911823.64729	0.77348\\
909819.63928	0.77447\\
907815.63126	0.77547\\
905811.62325	0.77646\\
903807.61523	0.77745\\
901803.60721	0.77844\\
899799.59920	0.77942\\
897795.59118	0.78041\\
895791.58317	0.78139\\
893787.57515	0.78237\\
891783.56713	0.78335\\
889779.55912	0.78433\\
887775.55110	0.78531\\
885771.54309	0.78628\\
883767.53507	0.78725\\
881763.52705	0.78822\\
879759.51904	0.78919\\
877755.51102	0.79016\\
875751.50301	0.79113\\
873747.49499	0.79209\\
871743.48697	0.79305\\
869739.47896	0.79401\\
867735.47094	0.79497\\
865731.46293	0.79592\\
863727.45491	0.79688\\
861723.44689	0.79783\\
859719.43888	0.79878\\
857715.43086	0.79973\\
855711.42285	0.80067\\
853707.41483	0.80162\\
851703.40681	0.80256\\
849699.39880	0.80350\\
847695.39078	0.80444\\
845691.38277	0.80538\\
843687.37475	0.80631\\
841683.36673	0.80724\\
839679.35872	0.80817\\
837675.35070	0.80910\\
835671.34269	0.81003\\
833667.33467	0.81096\\
831663.32665	0.81188\\
829659.31864	0.81280\\
827655.31062	0.81372\\
825651.30261	0.81464\\
823647.29459	0.81555\\
821643.28657	0.81646\\
819639.27856	0.81737\\
817635.27054	0.81828\\
815631.26253	0.81919\\
813627.25451	0.82010\\
811623.24649	0.82100\\
809619.23848	0.82190\\
807615.23046	0.82280\\
805611.22244	0.82370\\
803607.21443	0.82459\\
801603.20641	0.82548\\
799599.19840	0.82637\\
797595.19038	0.82726\\
795591.18236	0.82815\\
793587.17435	0.82903\\
791583.16633	0.82992\\
789579.15832	0.83080\\
787575.15030	0.83167\\
785571.14228	0.83255\\
783567.13427	0.83343\\
781563.12625	0.83430\\
779559.11824	0.83517\\
777555.11022	0.83604\\
775551.10220	0.83690\\
773547.09419	0.83776\\
771543.08617	0.83863\\
769539.07816	0.83949\\
767535.07014	0.84034\\
765531.06212	0.84120\\
763527.05411	0.84205\\
761523.04609	0.84290\\
759519.03808	0.84375\\
757515.03006	0.84460\\
755511.02204	0.84544\\
753507.01403	0.84628\\
751503.00601	0.84712\\
749498.99800	0.84796\\
747494.98998	0.84880\\
745490.98196	0.84963\\
743486.97395	0.85046\\
741482.96593	0.85129\\
739478.95792	0.85212\\
737474.94990	0.85294\\
735470.94188	0.85377\\
733466.93387	0.85459\\
731462.92585	0.85541\\
729458.91784	0.85622\\
727454.90982	0.85704\\
725450.90180	0.85785\\
723446.89379	0.85866\\
721442.88577	0.85946\\
719438.87776	0.86027\\
717434.86974	0.86107\\
715430.86172	0.86187\\
713426.85371	0.86267\\
711422.84569	0.86347\\
709418.83768	0.86426\\
707414.82966	0.86505\\
705410.82164	0.86584\\
703406.81363	0.86663\\
701402.80561	0.86742\\
699398.79760	0.86820\\
697394.78958	0.86898\\
695390.78156	0.86976\\
693386.77355	0.87053\\
691382.76553	0.87131\\
689378.75752	0.87208\\
687374.74950	0.87285\\
685370.74148	0.87361\\
683366.73347	0.87438\\
681362.72545	0.87514\\
679358.71743	0.87590\\
677354.70942	0.87666\\
675350.70140	0.87741\\
673346.69339	0.87817\\
671342.68537	0.87892\\
669338.67735	0.87967\\
667334.66934	0.88041\\
665330.66132	0.88116\\
663326.65331	0.88190\\
661322.64529	0.88264\\
659318.63727	0.88337\\
657314.62926	0.88411\\
655310.62124	0.88484\\
653306.61323	0.88557\\
651302.60521	0.88630\\
649298.59719	0.88702\\
647294.58918	0.88775\\
645290.58116	0.88847\\
643286.57315	0.88919\\
641282.56513	0.88990\\
639278.55711	0.89062\\
637274.54910	0.89133\\
635270.54108	0.89204\\
633266.53307	0.89274\\
631262.52505	0.89345\\
629258.51703	0.89415\\
627254.50902	0.89485\\
625250.50100	0.89555\\
623246.49299	0.89624\\
621242.48497	0.89693\\
619238.47695	0.89762\\
617234.46894	0.89831\\
615230.46092	0.89900\\
613226.45291	0.89968\\
611222.44489	0.90036\\
609218.43687	0.90104\\
607214.42886	0.90172\\
605210.42084	0.90239\\
603206.41283	0.90306\\
601202.40481	0.90373\\
599198.39679	0.90440\\
597194.38878	0.90506\\
595190.38076	0.90572\\
593186.37275	0.90638\\
591182.36473	0.90704\\
589178.35671	0.90769\\
587174.34870	0.90835\\
585170.34068	0.90900\\
583166.33267	0.90964\\
581162.32465	0.91029\\
579158.31663	0.91093\\
577154.30862	0.91157\\
575150.30060	0.91221\\
573146.29259	0.91285\\
571142.28457	0.91348\\
569138.27655	0.91411\\
567134.26854	0.91474\\
565130.26052	0.91536\\
563126.25251	0.91599\\
561122.24449	0.91661\\
559118.23647	0.91723\\
557114.22846	0.91784\\
555110.22044	0.91846\\
553106.21242	0.91907\\
551102.20441	0.91968\\
549098.19639	0.92028\\
547094.18838	0.92089\\
545090.18036	0.92149\\
543086.17234	0.92209\\
541082.16433	0.92269\\
539078.15631	0.92328\\
537074.14830	0.92387\\
535070.14028	0.92446\\
533066.13226	0.92505\\
531062.12425	0.92564\\
529058.11623	0.92622\\
527054.10822	0.92680\\
525050.10020	0.92738\\
523046.09218	0.92795\\
521042.08417	0.92853\\
519038.07615	0.92910\\
517034.06814	0.92966\\
515030.06012	0.93023\\
513026.05210	0.93079\\
511022.04409	0.93135\\
509018.03607	0.93191\\
507014.02806	0.93247\\
505010.02004	0.93302\\
503006.01202	0.93357\\
501002.00401	0.93412\\
498997.99599	0.93467\\
496993.98798	0.93521\\
494989.97996	0.93575\\
492985.97194	0.93629\\
490981.96393	0.93683\\
488977.95591	0.93736\\
486973.94790	0.93790\\
484969.93988	0.93843\\
482965.93186	0.93895\\
480961.92385	0.93948\\
478957.91583	0.94000\\
476953.90782	0.94052\\
474949.89980	0.94104\\
472945.89178	0.94155\\
470941.88377	0.94206\\
468937.87575	0.94257\\
466933.86774	0.94308\\
464929.85972	0.94359\\
462925.85170	0.94409\\
460921.84369	0.94459\\
458917.83567	0.94509\\
456913.82766	0.94558\\
454909.81964	0.94608\\
452905.81162	0.94657\\
450901.80361	0.94706\\
448897.79559	0.94754\\
446893.78758	0.94803\\
444889.77956	0.94851\\
442885.77154	0.94899\\
440881.76353	0.94946\\
438877.75551	0.94994\\
436873.74749	0.95041\\
434869.73948	0.95088\\
432865.73146	0.95134\\
430861.72345	0.95181\\
428857.71543	0.95227\\
426853.70741	0.95273\\
424849.69940	0.95319\\
422845.69138	0.95364\\
420841.68337	0.95409\\
418837.67535	0.95454\\
416833.66733	0.95499\\
414829.65932	0.95544\\
412825.65130	0.95588\\
410821.64329	0.95632\\
408817.63527	0.95676\\
406813.62725	0.95719\\
404809.61924	0.95763\\
402805.61122	0.95806\\
400801.60321	0.95849\\
398797.59519	0.95891\\
396793.58717	0.95934\\
394789.57916	0.95976\\
392785.57114	0.96018\\
390781.56313	0.96059\\
388777.55511	0.96101\\
386773.54709	0.96142\\
384769.53908	0.96183\\
382765.53106	0.96224\\
380761.52305	0.96264\\
378757.51503	0.96304\\
376753.50701	0.96344\\
374749.49900	0.96384\\
372745.49098	0.96424\\
370741.48297	0.96463\\
368737.47495	0.96502\\
366733.46693	0.96541\\
364729.45892	0.96579\\
362725.45090	0.96618\\
360721.44289	0.96656\\
358717.43487	0.96694\\
356713.42685	0.96731\\
354709.41884	0.96769\\
352705.41082	0.96806\\
350701.40281	0.96843\\
348697.39479	0.96879\\
346693.38677	0.96916\\
344689.37876	0.96952\\
342685.37074	0.96988\\
340681.36273	0.97024\\
338677.35471	0.97059\\
336673.34669	0.97095\\
334669.33868	0.97130\\
332665.33066	0.97165\\
330661.32265	0.97199\\
328657.31463	0.97234\\
326653.30661	0.97268\\
324649.29860	0.97302\\
322645.29058	0.97335\\
320641.28257	0.97369\\
318637.27455	0.97402\\
316633.26653	0.97435\\
314629.25852	0.97468\\
312625.25050	0.97500\\
310621.24248	0.97533\\
308617.23447	0.97565\\
306613.22645	0.97597\\
304609.21844	0.97628\\
302605.21042	0.97660\\
300601.20240	0.97691\\
298597.19439	0.97722\\
296593.18637	0.97752\\
294589.17836	0.97783\\
292585.17034	0.97813\\
290581.16232	0.97843\\
288577.15431	0.97873\\
286573.14629	0.97903\\
284569.13828	0.97932\\
282565.13026	0.97961\\
280561.12224	0.97990\\
278557.11423	0.98019\\
276553.10621	0.98047\\
274549.09820	0.98075\\
272545.09018	0.98103\\
270541.08216	0.98131\\
268537.07415	0.98159\\
266533.06613	0.98186\\
264529.05812	0.98213\\
262525.05010	0.98240\\
260521.04208	0.98267\\
258517.03407	0.98293\\
256513.02605	0.98319\\
254509.01804	0.98345\\
252505.01002	0.98371\\
250501.00200	0.98397\\
248496.99399	0.98422\\
246492.98597	0.98447\\
244488.97796	0.98472\\
242484.96994	0.98497\\
240480.96192	0.98521\\
238476.95391	0.98546\\
236472.94589	0.98570\\
234468.93788	0.98593\\
232464.92986	0.98617\\
230460.92184	0.98640\\
228456.91383	0.98664\\
226452.90581	0.98687\\
224448.89780	0.98709\\
222444.88978	0.98732\\
220440.88176	0.98754\\
218436.87375	0.98776\\
216432.86573	0.98798\\
214428.85772	0.98820\\
212424.84970	0.98841\\
210420.84168	0.98862\\
208416.83367	0.98883\\
206412.82565	0.98904\\
204408.81764	0.98925\\
202404.80962	0.98945\\
200400.80160	0.98965\\
198396.79359	0.98985\\
196392.78557	0.99005\\
194388.77756	0.99025\\
192384.76954	0.99044\\
190380.76152	0.99063\\
188376.75351	0.99082\\
186372.74549	0.99101\\
184368.73747	0.99119\\
182364.72946	0.99138\\
180360.72144	0.99156\\
178356.71343	0.99174\\
176352.70541	0.99191\\
174348.69739	0.99209\\
172344.68938	0.99226\\
170340.68136	0.99243\\
168336.67335	0.99260\\
166332.66533	0.99277\\
164328.65731	0.99293\\
162324.64930	0.99309\\
160320.64128	0.99326\\
158316.63327	0.99341\\
156312.62525	0.99357\\
154308.61723	0.99373\\
152304.60922	0.99388\\
150300.60120	0.99403\\
148296.59319	0.99418\\
146292.58517	0.99432\\
144288.57715	0.99447\\
142284.56914	0.99461\\
140280.56112	0.99475\\
138276.55311	0.99489\\
136272.54509	0.99503\\
134268.53707	0.99516\\
132264.52906	0.99529\\
130260.52104	0.99542\\
128256.51303	0.99555\\
126252.50501	0.99568\\
124248.49699	0.99581\\
122244.48898	0.99593\\
120240.48096	0.99605\\
118236.47295	0.99617\\
116232.46493	0.99629\\
114228.45691	0.99640\\
112224.44890	0.99652\\
110220.44088	0.99663\\
108216.43287	0.99674\\
106212.42485	0.99684\\
104208.41683	0.99695\\
102204.40882	0.99705\\
100200.40080	0.99716\\
98196.39279	0.99726\\
96192.38477	0.99735\\
94188.37675	0.99745\\
92184.36874	0.99755\\
90180.36072	0.99764\\
88176.35271	0.99773\\
86172.34469	0.99782\\
84168.33667	0.99791\\
82164.32866	0.99799\\
80160.32064	0.99807\\
78156.31263	0.99816\\
76152.30461	0.99824\\
74148.29659	0.99831\\
72144.28858	0.99839\\
70140.28056	0.99846\\
68136.27255	0.99854\\
66132.26453	0.99861\\
64128.25651	0.99868\\
62124.24850	0.99874\\
60120.24048	0.99881\\
58116.23246	0.99887\\
56112.22445	0.99894\\
54108.21643	0.99900\\
52104.20842	0.99905\\
50100.20040	0.99911\\
48096.19238	0.99916\\
46092.18437	0.99922\\
44088.17635	0.99927\\
42084.16834	0.99932\\
40080.16032	0.99937\\
38076.15230	0.99941\\
36072.14429	0.99945\\
34068.13627	0.99950\\
32064.12826	0.99954\\
30060.12024	0.99958\\
28056.11222	0.99961\\
26052.10421	0.99965\\
24048.09619	0.99968\\
22044.08818	0.99971\\
20040.08016	0.99974\\
18036.07214	0.99976\\
16032.06413	0.99978\\
14028.05611	0.99980\\
12024.04810	0.99981\\
10020.04008	0.99976\\
8016.03206	0.99977\\
6012.02405	0.99987\\
4008.01603	0.99995\\
2004.00802	0.99999\\
0.00000	1.00000\\
}--cycle;
\addplot [color=black, line width=1.5pt]
  table[row sep=crcr]{%
0.00000	1.00000\\
2004.00802	0.99999\\
4008.01603	0.99995\\
6012.02405	0.99987\\
8016.03206	0.99977\\
10020.04008	0.99963\\
12024.04810	0.99946\\
14028.05611	0.99925\\
16032.06413	0.99900\\
18036.07214	0.99872\\
20040.08016	0.99840\\
22044.08818	0.99804\\
24048.09619	0.99765\\
26052.10421	0.99721\\
28056.11222	0.99674\\
30060.12024	0.99623\\
32064.12826	0.99568\\
34068.13627	0.99509\\
36072.14429	0.99446\\
38076.15230	0.99379\\
40080.16032	0.99309\\
42084.16834	0.99234\\
44088.17635	0.99155\\
46092.18437	0.99072\\
48096.19238	0.98985\\
50100.20040	0.98894\\
52104.20842	0.98799\\
54108.21643	0.98700\\
56112.22445	0.98597\\
58116.23246	0.98490\\
60120.24048	0.98379\\
62124.24850	0.98264\\
64128.25651	0.98144\\
66132.26453	0.98021\\
68136.27255	0.97893\\
70140.28056	0.97762\\
72144.28858	0.97626\\
74148.29659	0.97486\\
76152.30461	0.97343\\
78156.31263	0.97195\\
80160.32064	0.97043\\
82164.32866	0.96887\\
84168.33667	0.96727\\
86172.34469	0.96563\\
88176.35271	0.96395\\
90180.36072	0.96223\\
92184.36874	0.96047\\
94188.37675	0.95867\\
96192.38477	0.95683\\
98196.39279	0.95495\\
100200.40080	0.95304\\
102204.40882	0.95108\\
104208.41683	0.94908\\
106212.42485	0.94705\\
108216.43287	0.94497\\
110220.44088	0.94286\\
112224.44890	0.94071\\
114228.45691	0.93852\\
116232.46493	0.93630\\
118236.47295	0.93403\\
120240.48096	0.93173\\
122244.48898	0.92939\\
124248.49699	0.92702\\
126252.50501	0.92460\\
128256.51303	0.92215\\
130260.52104	0.91967\\
132264.52906	0.91714\\
134268.53707	0.91459\\
136272.54509	0.91199\\
138276.55311	0.90937\\
140280.56112	0.90670\\
142284.56914	0.90400\\
144288.57715	0.90127\\
146292.58517	0.89850\\
148296.59319	0.89570\\
150300.60120	0.89287\\
152304.60922	0.89000\\
154308.61723	0.88710\\
156312.62525	0.88417\\
158316.63327	0.88120\\
160320.64128	0.87821\\
162324.64930	0.87518\\
164328.65731	0.87212\\
166332.66533	0.86903\\
168336.67335	0.86591\\
170340.68136	0.86275\\
172344.68938	0.85957\\
174348.69739	0.85636\\
176352.70541	0.85312\\
178356.71343	0.84985\\
180360.72144	0.84656\\
182364.72946	0.84323\\
184368.73747	0.83988\\
186372.74549	0.83650\\
188376.75351	0.83309\\
190380.76152	0.82966\\
192384.76954	0.82620\\
194388.77756	0.82271\\
196392.78557	0.81920\\
198396.79359	0.81567\\
200400.80160	0.81211\\
202404.80962	0.80852\\
204408.81764	0.80492\\
206412.82565	0.80129\\
208416.83367	0.79763\\
210420.84168	0.79396\\
212424.84970	0.79026\\
214428.85772	0.78654\\
216432.86573	0.78280\\
218436.87375	0.77904\\
220440.88176	0.77526\\
222444.88978	0.77145\\
224448.89780	0.76763\\
226452.90581	0.76379\\
228456.91383	0.75994\\
230460.92184	0.75606\\
232464.92986	0.75217\\
234468.93788	0.74825\\
236472.94589	0.74433\\
238476.95391	0.74038\\
240480.96192	0.73642\\
242484.96994	0.73245\\
244488.97796	0.72846\\
246492.98597	0.72445\\
248496.99399	0.72043\\
250501.00200	0.71640\\
252505.01002	0.71235\\
254509.01804	0.70830\\
256513.02605	0.70422\\
258517.03407	0.70014\\
260521.04208	0.69605\\
262525.05010	0.69194\\
264529.05812	0.68783\\
266533.06613	0.68370\\
268537.07415	0.67956\\
270541.08216	0.67542\\
272545.09018	0.67127\\
274549.09820	0.66711\\
276553.10621	0.66294\\
278557.11423	0.65876\\
280561.12224	0.65457\\
282565.13026	0.65038\\
284569.13828	0.64619\\
286573.14629	0.64199\\
288577.15431	0.63778\\
290581.16232	0.63357\\
292585.17034	0.62935\\
294589.17836	0.62513\\
296593.18637	0.62091\\
298597.19439	0.61668\\
300601.20240	0.61245\\
302605.21042	0.60822\\
304609.21844	0.60399\\
306613.22645	0.59975\\
308617.23447	0.59552\\
310621.24248	0.59128\\
312625.25050	0.58704\\
314629.25852	0.58281\\
316633.26653	0.57857\\
318637.27455	0.57434\\
320641.28257	0.57011\\
322645.29058	0.56588\\
324649.29860	0.56165\\
326653.30661	0.55742\\
328657.31463	0.55320\\
330661.32265	0.54898\\
332665.33066	0.54477\\
334669.33868	0.54056\\
336673.34669	0.53635\\
338677.35471	0.53215\\
340681.36273	0.52796\\
342685.37074	0.52377\\
344689.37876	0.51959\\
346693.38677	0.51541\\
348697.39479	0.51124\\
350701.40281	0.50708\\
352705.41082	0.50292\\
354709.41884	0.49878\\
356713.42685	0.49464\\
358717.43487	0.49051\\
360721.44289	0.48639\\
362725.45090	0.48228\\
364729.45892	0.47818\\
366733.46693	0.47409\\
368737.47495	0.47000\\
370741.48297	0.46593\\
372745.49098	0.46188\\
374749.49900	0.45783\\
376753.50701	0.45379\\
378757.51503	0.44977\\
380761.52305	0.44575\\
382765.53106	0.44175\\
384769.53908	0.43777\\
386773.54709	0.43379\\
388777.55511	0.42983\\
390781.56313	0.42589\\
392785.57114	0.42195\\
394789.57916	0.41804\\
396793.58717	0.41413\\
398797.59519	0.41024\\
400801.60321	0.40637\\
402805.61122	0.40251\\
404809.61924	0.39867\\
406813.62725	0.39484\\
408817.63527	0.39103\\
410821.64329	0.38723\\
412825.65130	0.38345\\
414829.65932	0.37969\\
416833.66733	0.37595\\
418837.67535	0.37222\\
420841.68337	0.36851\\
422845.69138	0.36481\\
424849.69940	0.36114\\
426853.70741	0.35748\\
428857.71543	0.35384\\
430861.72345	0.35022\\
432865.73146	0.34662\\
434869.73948	0.34303\\
436873.74749	0.33947\\
438877.75551	0.33592\\
440881.76353	0.33240\\
442885.77154	0.32889\\
444889.77956	0.32540\\
446893.78758	0.32193\\
448897.79559	0.31849\\
450901.80361	0.31506\\
452905.81162	0.31165\\
454909.81964	0.30826\\
456913.82766	0.30490\\
458917.83567	0.30155\\
460921.84369	0.29822\\
462925.85170	0.29492\\
464929.85972	0.29164\\
466933.86774	0.28837\\
468937.87575	0.28513\\
470941.88377	0.28191\\
472945.89178	0.27871\\
474949.89980	0.27553\\
476953.90782	0.27238\\
478957.91583	0.26924\\
480961.92385	0.26613\\
482965.93186	0.26304\\
484969.93988	0.25997\\
486973.94790	0.25692\\
488977.95591	0.25390\\
490981.96393	0.25089\\
492985.97194	0.24791\\
494989.97996	0.24495\\
496993.98798	0.24202\\
498997.99599	0.23910\\
501002.00401	0.23621\\
503006.01202	0.23334\\
505010.02004	0.23049\\
507014.02806	0.22766\\
509018.03607	0.22486\\
511022.04409	0.22208\\
513026.05210	0.21932\\
515030.06012	0.21659\\
517034.06814	0.21387\\
519038.07615	0.21118\\
521042.08417	0.20852\\
523046.09218	0.20587\\
525050.10020	0.20325\\
527054.10822	0.20065\\
529058.11623	0.19807\\
531062.12425	0.19551\\
533066.13226	0.19298\\
535070.14028	0.19047\\
537074.14830	0.18798\\
539078.15631	0.18551\\
541082.16433	0.18307\\
543086.17234	0.18065\\
545090.18036	0.17825\\
547094.18838	0.17587\\
549098.19639	0.17351\\
551102.20441	0.17118\\
553106.21242	0.16887\\
555110.22044	0.16658\\
557114.22846	0.16431\\
559118.23647	0.16207\\
561122.24449	0.15985\\
563126.25251	0.15765\\
565130.26052	0.15547\\
567134.26854	0.15331\\
569138.27655	0.15117\\
571142.28457	0.14906\\
573146.29259	0.14696\\
575150.30060	0.14489\\
577154.30862	0.14284\\
579158.31663	0.14081\\
581162.32465	0.13881\\
583166.33267	0.13682\\
585170.34068	0.13485\\
587174.34870	0.13291\\
589178.35671	0.13098\\
591182.36473	0.12908\\
593186.37275	0.12720\\
595190.38076	0.12534\\
597194.38878	0.12349\\
599198.39679	0.12167\\
601202.40481	0.11987\\
603206.41283	0.11809\\
605210.42084	0.11633\\
607214.42886	0.11459\\
609218.43687	0.11287\\
611222.44489	0.11116\\
613226.45291	0.10948\\
615230.46092	0.10782\\
617234.46894	0.10617\\
619238.47695	0.10455\\
621242.48497	0.10295\\
623246.49299	0.10136\\
625250.50100	0.09979\\
627254.50902	0.09824\\
629258.51703	0.09671\\
631262.52505	0.09520\\
633266.53307	0.09371\\
635270.54108	0.09224\\
637274.54910	0.09078\\
639278.55711	0.08934\\
641282.56513	0.08792\\
643286.57315	0.08652\\
645290.58116	0.08513\\
647294.58918	0.08376\\
649298.59719	0.08241\\
651302.60521	0.08108\\
653306.61323	0.07976\\
655310.62124	0.07847\\
657314.62926	0.07718\\
659318.63727	0.07592\\
661322.64529	0.07467\\
663326.65331	0.07344\\
665330.66132	0.07222\\
667334.66934	0.07102\\
669338.67735	0.06984\\
671342.68537	0.06867\\
673346.69339	0.06752\\
675350.70140	0.06638\\
677354.70942	0.06526\\
679358.71743	0.06415\\
681362.72545	0.06306\\
683366.73347	0.06198\\
685370.74148	0.06092\\
687374.74950	0.05988\\
689378.75752	0.05885\\
691382.76553	0.05783\\
693386.77355	0.05683\\
695390.78156	0.05584\\
697394.78958	0.05487\\
699398.79760	0.05390\\
701402.80561	0.05296\\
703406.81363	0.05203\\
705410.82164	0.05111\\
707414.82966	0.05020\\
709418.83768	0.04931\\
711422.84569	0.04843\\
713426.85371	0.04756\\
715430.86172	0.04671\\
717434.86974	0.04587\\
719438.87776	0.04504\\
721442.88577	0.04422\\
723446.89379	0.04342\\
725450.90180	0.04263\\
727454.90982	0.04185\\
729458.91784	0.04108\\
731462.92585	0.04033\\
733466.93387	0.03958\\
735470.94188	0.03885\\
737474.94990	0.03813\\
739478.95792	0.03742\\
741482.96593	0.03672\\
743486.97395	0.03603\\
745490.98196	0.03536\\
747494.98998	0.03469\\
749498.99800	0.03404\\
751503.00601	0.03339\\
753507.01403	0.03276\\
755511.02204	0.03213\\
757515.03006	0.03152\\
759519.03808	0.03091\\
761523.04609	0.03032\\
763527.05411	0.02974\\
765531.06212	0.02916\\
767535.07014	0.02859\\
769539.07816	0.02804\\
771543.08617	0.02749\\
773547.09419	0.02695\\
775551.10220	0.02643\\
777555.11022	0.02591\\
779559.11824	0.02539\\
781563.12625	0.02489\\
783567.13427	0.02440\\
785571.14228	0.02391\\
787575.15030	0.02344\\
789579.15832	0.02297\\
791583.16633	0.02251\\
793587.17435	0.02205\\
795591.18236	0.02161\\
797595.19038	0.02117\\
799599.19840	0.02074\\
801603.20641	0.02032\\
803607.21443	0.01990\\
805611.22244	0.01950\\
807615.23046	0.01910\\
809619.23848	0.01870\\
811623.24649	0.01832\\
813627.25451	0.01794\\
815631.26253	0.01757\\
817635.27054	0.01720\\
819639.27856	0.01684\\
821643.28657	0.01649\\
823647.29459	0.01615\\
825651.30261	0.01581\\
827655.31062	0.01547\\
829659.31864	0.01515\\
831663.32665	0.01483\\
833667.33467	0.01451\\
835671.34269	0.01420\\
837675.35070	0.01390\\
839679.35872	0.01360\\
841683.36673	0.01331\\
843687.37475	0.01302\\
845691.38277	0.01274\\
847695.39078	0.01246\\
849699.39880	0.01219\\
851703.40681	0.01193\\
853707.41483	0.01167\\
855711.42285	0.01141\\
857715.43086	0.01116\\
859719.43888	0.01092\\
861723.44689	0.01068\\
863727.45491	0.01044\\
865731.46293	0.01021\\
867735.47094	0.00998\\
869739.47896	0.00976\\
871743.48697	0.00954\\
873747.49499	0.00933\\
875751.50301	0.00912\\
877755.51102	0.00892\\
879759.51904	0.00871\\
881763.52705	0.00852\\
883767.53507	0.00832\\
885771.54309	0.00814\\
887775.55110	0.00795\\
889779.55912	0.00777\\
891783.56713	0.00759\\
893787.57515	0.00742\\
895791.58317	0.00725\\
897795.59118	0.00708\\
899799.59920	0.00692\\
901803.60721	0.00676\\
903807.61523	0.00660\\
905811.62325	0.00645\\
907815.63126	0.00629\\
909819.63928	0.00615\\
911823.64729	0.00600\\
913827.65531	0.00586\\
915831.66333	0.00572\\
917835.67134	0.00559\\
919839.67936	0.00546\\
921843.68737	0.00533\\
923847.69539	0.00520\\
925851.70341	0.00508\\
927855.71142	0.00495\\
929859.71944	0.00484\\
931863.72745	0.00472\\
933867.73547	0.00461\\
935871.74349	0.00450\\
937875.75150	0.00439\\
939879.75952	0.00428\\
941883.76754	0.00418\\
943887.77555	0.00407\\
945891.78357	0.00397\\
947895.79158	0.00388\\
949899.79960	0.00378\\
951903.80762	0.00369\\
953907.81563	0.00360\\
955911.82365	0.00351\\
957915.83166	0.00342\\
959919.83968	0.00334\\
961923.84770	0.00325\\
963927.85571	0.00317\\
965931.86373	0.00309\\
967935.87174	0.00302\\
969939.87976	0.00294\\
971943.88778	0.00287\\
973947.89579	0.00279\\
975951.90381	0.00272\\
977955.91182	0.00266\\
979959.91984	0.00259\\
981963.92786	0.00252\\
983967.93587	0.00246\\
985971.94389	0.00239\\
987975.95190	0.00233\\
989979.95992	0.00227\\
991983.96794	0.00222\\
993987.97595	0.00216\\
995991.98397	0.00210\\
997995.99198	0.00205\\
1000000.00000	0.00199\\
};
\addlegendentry{\ac{EmALT} (90\,\% \ac{CI})}

\addplot [color=black, dashed, line width=1.0pt, forget plot]
  table[row sep=crcr]{%
0.00000	1.00000\\
2004.00802	0.99999\\
4008.01603	0.99995\\
6012.02405	0.99987\\
8016.03206	0.99977\\
10020.04008	0.99943\\
12024.04810	0.99847\\
14028.05611	0.99716\\
16032.06413	0.99538\\
18036.07214	0.99310\\
20040.08016	0.99027\\
22044.08818	0.98686\\
24048.09619	0.98285\\
26052.10421	0.97820\\
28056.11222	0.97291\\
30060.12024	0.96695\\
32064.12826	0.96032\\
34068.13627	0.95301\\
36072.14429	0.94502\\
38076.15230	0.93635\\
40080.16032	0.92700\\
42084.16834	0.91698\\
44088.17635	0.90630\\
46092.18437	0.89496\\
48096.19238	0.88300\\
50100.20040	0.87042\\
52104.20842	0.85725\\
54108.21643	0.84351\\
56112.22445	0.82922\\
58116.23246	0.81442\\
60120.24048	0.79912\\
62124.24850	0.78337\\
64128.25651	0.76719\\
66132.26453	0.75062\\
68136.27255	0.73368\\
70140.28056	0.71643\\
72144.28858	0.69889\\
74148.29659	0.68110\\
76152.30461	0.66310\\
78156.31263	0.64492\\
80160.32064	0.62661\\
82164.32866	0.60819\\
84168.33667	0.58972\\
86172.34469	0.57121\\
88176.35271	0.55271\\
90180.36072	0.53426\\
92184.36874	0.51589\\
94188.37675	0.49762\\
96192.38477	0.47950\\
98196.39279	0.46155\\
100200.40080	0.44381\\
102204.40882	0.42629\\
104208.41683	0.40903\\
106212.42485	0.39205\\
108216.43287	0.37537\\
110220.44088	0.35902\\
112224.44890	0.34301\\
114228.45691	0.32736\\
116232.46493	0.31209\\
118236.47295	0.29721\\
120240.48096	0.28273\\
122244.48898	0.26866\\
124248.49699	0.25502\\
126252.50501	0.24180\\
128256.51303	0.22902\\
130260.52104	0.21668\\
132264.52906	0.20478\\
134268.53707	0.19332\\
136272.54509	0.18230\\
138276.55311	0.17172\\
140280.56112	0.16158\\
142284.56914	0.15186\\
144288.57715	0.14258\\
146292.58517	0.13371\\
148296.59319	0.12526\\
150300.60120	0.11721\\
152304.60922	0.10956\\
154308.61723	0.10229\\
156312.62525	0.09540\\
158316.63327	0.08888\\
160320.64128	0.08271\\
162324.64930	0.07688\\
164328.65731	0.07138\\
166332.66533	0.06621\\
168336.67335	0.06134\\
170340.68136	0.05676\\
172344.68938	0.05247\\
174348.69739	0.04845\\
176352.70541	0.04469\\
178356.71343	0.04117\\
180360.72144	0.03789\\
182364.72946	0.03483\\
184368.73747	0.03198\\
186372.74549	0.02933\\
188376.75351	0.02687\\
190380.76152	0.02459\\
192384.76954	0.02248\\
194388.77756	0.02053\\
196392.78557	0.01873\\
198396.79359	0.01706\\
200400.80160	0.01553\\
202404.80962	0.01411\\
204408.81764	0.01282\\
206412.82565	0.01163\\
208416.83367	0.01053\\
210420.84168	0.00953\\
212424.84970	0.00862\\
214428.85772	0.00778\\
216432.86573	0.00702\\
218436.87375	0.00632\\
220440.88176	0.00569\\
222444.88978	0.00512\\
224448.89780	0.00460\\
226452.90581	0.00412\\
228456.91383	0.00369\\
230460.92184	0.00331\\
232464.92986	0.00296\\
234468.93788	0.00264\\
236472.94589	0.00235\\
238476.95391	0.00210\\
240480.96192	0.00187\\
242484.96994	0.00166\\
244488.97796	0.00147\\
246492.98597	0.00131\\
248496.99399	0.00116\\
250501.00200	0.00103\\
252505.01002	0.00091\\
254509.01804	0.00080\\
256513.02605	0.00071\\
258517.03407	0.00062\\
260521.04208	0.00055\\
262525.05010	0.00048\\
264529.05812	0.00042\\
266533.06613	0.00037\\
268537.07415	0.00033\\
270541.08216	0.00029\\
272545.09018	0.00025\\
274549.09820	0.00022\\
276553.10621	0.00019\\
278557.11423	0.00017\\
280561.12224	0.00014\\
282565.13026	0.00013\\
284569.13828	0.00011\\
286573.14629	0.00009\\
288577.15431	0.00008\\
290581.16232	0.00007\\
292585.17034	0.00006\\
294589.17836	0.00005\\
296593.18637	0.00005\\
298597.19439	0.00004\\
300601.20240	0.00003\\
302605.21042	0.00003\\
304609.21844	0.00003\\
306613.22645	0.00002\\
308617.23447	0.00002\\
310621.24248	0.00002\\
312625.25050	0.00001\\
314629.25852	0.00001\\
316633.26653	0.00001\\
318637.27455	0.00001\\
320641.28257	0.00001\\
322645.29058	0.00001\\
324649.29860	0.00001\\
326653.30661	0.00000\\
328657.31463	0.00000\\
330661.32265	0.00000\\
332665.33066	0.00000\\
334669.33868	0.00000\\
336673.34669	0.00000\\
338677.35471	0.00000\\
340681.36273	0.00000\\
342685.37074	0.00000\\
344689.37876	0.00000\\
346693.38677	0.00000\\
348697.39479	0.00000\\
350701.40281	0.00000\\
352705.41082	0.00000\\
354709.41884	0.00000\\
356713.42685	0.00000\\
358717.43487	0.00000\\
360721.44289	0.00000\\
362725.45090	0.00000\\
364729.45892	0.00000\\
366733.46693	0.00000\\
368737.47495	0.00000\\
370741.48297	0.00000\\
372745.49098	0.00000\\
374749.49900	0.00000\\
376753.50701	0.00000\\
378757.51503	0.00000\\
380761.52305	0.00000\\
382765.53106	0.00000\\
384769.53908	0.00000\\
386773.54709	0.00000\\
388777.55511	0.00000\\
390781.56313	0.00000\\
392785.57114	0.00000\\
394789.57916	0.00000\\
396793.58717	0.00000\\
398797.59519	0.00000\\
400801.60321	0.00000\\
402805.61122	0.00000\\
404809.61924	0.00000\\
406813.62725	0.00000\\
408817.63527	0.00000\\
410821.64329	0.00000\\
412825.65130	0.00000\\
414829.65932	0.00000\\
416833.66733	0.00000\\
418837.67535	0.00000\\
420841.68337	0.00000\\
422845.69138	0.00000\\
424849.69940	0.00000\\
426853.70741	0.00000\\
428857.71543	0.00000\\
430861.72345	0.00000\\
432865.73146	0.00000\\
434869.73948	0.00000\\
436873.74749	0.00000\\
438877.75551	0.00000\\
440881.76353	0.00000\\
442885.77154	0.00000\\
444889.77956	0.00000\\
446893.78758	0.00000\\
448897.79559	0.00000\\
450901.80361	0.00000\\
452905.81162	0.00000\\
454909.81964	0.00000\\
456913.82766	0.00000\\
458917.83567	0.00000\\
460921.84369	0.00000\\
462925.85170	0.00000\\
464929.85972	0.00000\\
466933.86774	0.00000\\
468937.87575	0.00000\\
470941.88377	0.00000\\
472945.89178	0.00000\\
474949.89980	0.00000\\
476953.90782	0.00000\\
478957.91583	0.00000\\
480961.92385	0.00000\\
482965.93186	0.00000\\
484969.93988	0.00000\\
486973.94790	0.00000\\
488977.95591	0.00000\\
490981.96393	0.00000\\
492985.97194	0.00000\\
494989.97996	0.00000\\
496993.98798	0.00000\\
498997.99599	0.00000\\
501002.00401	0.00000\\
503006.01202	0.00000\\
505010.02004	0.00000\\
507014.02806	0.00000\\
509018.03607	0.00000\\
511022.04409	0.00000\\
513026.05210	0.00000\\
515030.06012	0.00000\\
517034.06814	0.00000\\
519038.07615	0.00000\\
521042.08417	0.00000\\
523046.09218	0.00000\\
525050.10020	0.00000\\
527054.10822	0.00000\\
529058.11623	0.00000\\
531062.12425	0.00000\\
533066.13226	0.00000\\
535070.14028	0.00000\\
537074.14830	0.00000\\
539078.15631	0.00000\\
541082.16433	0.00000\\
543086.17234	0.00000\\
545090.18036	0.00000\\
547094.18838	0.00000\\
549098.19639	0.00000\\
551102.20441	0.00000\\
553106.21242	0.00000\\
555110.22044	0.00000\\
557114.22846	0.00000\\
559118.23647	0.00000\\
561122.24449	0.00000\\
563126.25251	0.00000\\
565130.26052	0.00000\\
567134.26854	0.00000\\
569138.27655	0.00000\\
571142.28457	0.00000\\
573146.29259	0.00000\\
575150.30060	0.00000\\
577154.30862	0.00000\\
579158.31663	0.00000\\
581162.32465	0.00000\\
583166.33267	0.00000\\
585170.34068	0.00000\\
587174.34870	0.00000\\
589178.35671	0.00000\\
591182.36473	0.00000\\
593186.37275	0.00000\\
595190.38076	0.00000\\
597194.38878	0.00000\\
599198.39679	0.00000\\
601202.40481	0.00000\\
603206.41283	0.00000\\
605210.42084	0.00000\\
607214.42886	0.00000\\
609218.43687	0.00000\\
611222.44489	0.00000\\
613226.45291	0.00000\\
615230.46092	0.00000\\
617234.46894	0.00000\\
619238.47695	0.00000\\
621242.48497	0.00000\\
623246.49299	0.00000\\
625250.50100	0.00000\\
627254.50902	0.00000\\
629258.51703	0.00000\\
631262.52505	0.00000\\
633266.53307	0.00000\\
635270.54108	0.00000\\
637274.54910	0.00000\\
639278.55711	0.00000\\
641282.56513	0.00000\\
643286.57315	0.00000\\
645290.58116	0.00000\\
647294.58918	0.00000\\
649298.59719	0.00000\\
651302.60521	0.00000\\
653306.61323	0.00000\\
655310.62124	0.00000\\
657314.62926	0.00000\\
659318.63727	0.00000\\
661322.64529	0.00000\\
663326.65331	0.00000\\
665330.66132	0.00000\\
667334.66934	0.00000\\
669338.67735	0.00000\\
671342.68537	0.00000\\
673346.69339	0.00000\\
675350.70140	0.00000\\
677354.70942	0.00000\\
679358.71743	0.00000\\
681362.72545	0.00000\\
683366.73347	0.00000\\
685370.74148	0.00000\\
687374.74950	0.00000\\
689378.75752	0.00000\\
691382.76553	0.00000\\
693386.77355	0.00000\\
695390.78156	0.00000\\
697394.78958	0.00000\\
699398.79760	0.00000\\
701402.80561	0.00000\\
703406.81363	0.00000\\
705410.82164	0.00000\\
707414.82966	0.00000\\
709418.83768	0.00000\\
711422.84569	0.00000\\
713426.85371	0.00000\\
715430.86172	0.00000\\
717434.86974	0.00000\\
719438.87776	0.00000\\
721442.88577	0.00000\\
723446.89379	0.00000\\
725450.90180	0.00000\\
727454.90982	0.00000\\
729458.91784	0.00000\\
731462.92585	0.00000\\
733466.93387	0.00000\\
735470.94188	0.00000\\
737474.94990	0.00000\\
739478.95792	0.00000\\
741482.96593	0.00000\\
743486.97395	0.00000\\
745490.98196	0.00000\\
747494.98998	0.00000\\
749498.99800	0.00000\\
751503.00601	0.00000\\
753507.01403	0.00000\\
755511.02204	0.00000\\
757515.03006	0.00000\\
759519.03808	0.00000\\
761523.04609	0.00000\\
763527.05411	0.00000\\
765531.06212	0.00000\\
767535.07014	0.00000\\
769539.07816	0.00000\\
771543.08617	0.00000\\
773547.09419	0.00000\\
775551.10220	0.00000\\
777555.11022	0.00000\\
779559.11824	0.00000\\
781563.12625	0.00000\\
783567.13427	0.00000\\
785571.14228	0.00000\\
787575.15030	0.00000\\
789579.15832	0.00000\\
791583.16633	0.00000\\
793587.17435	0.00000\\
795591.18236	0.00000\\
797595.19038	0.00000\\
799599.19840	0.00000\\
801603.20641	0.00000\\
803607.21443	0.00000\\
805611.22244	0.00000\\
807615.23046	0.00000\\
809619.23848	0.00000\\
811623.24649	0.00000\\
813627.25451	0.00000\\
815631.26253	0.00000\\
817635.27054	0.00000\\
819639.27856	0.00000\\
821643.28657	0.00000\\
823647.29459	0.00000\\
825651.30261	0.00000\\
827655.31062	0.00000\\
829659.31864	0.00000\\
831663.32665	0.00000\\
833667.33467	0.00000\\
835671.34269	0.00000\\
837675.35070	0.00000\\
839679.35872	0.00000\\
841683.36673	0.00000\\
843687.37475	0.00000\\
845691.38277	0.00000\\
847695.39078	0.00000\\
849699.39880	0.00000\\
851703.40681	0.00000\\
853707.41483	0.00000\\
855711.42285	0.00000\\
857715.43086	0.00000\\
859719.43888	0.00000\\
861723.44689	0.00000\\
863727.45491	0.00000\\
865731.46293	0.00000\\
867735.47094	0.00000\\
869739.47896	0.00000\\
871743.48697	0.00000\\
873747.49499	0.00000\\
875751.50301	0.00000\\
877755.51102	0.00000\\
879759.51904	0.00000\\
881763.52705	0.00000\\
883767.53507	0.00000\\
885771.54309	0.00000\\
887775.55110	0.00000\\
889779.55912	0.00000\\
891783.56713	0.00000\\
893787.57515	0.00000\\
895791.58317	0.00000\\
897795.59118	0.00000\\
899799.59920	0.00000\\
901803.60721	0.00000\\
903807.61523	0.00000\\
905811.62325	0.00000\\
907815.63126	0.00000\\
909819.63928	0.00000\\
911823.64729	0.00000\\
913827.65531	0.00000\\
915831.66333	0.00000\\
917835.67134	0.00000\\
919839.67936	0.00000\\
921843.68737	0.00000\\
923847.69539	0.00000\\
925851.70341	0.00000\\
927855.71142	0.00000\\
929859.71944	0.00000\\
931863.72745	0.00000\\
933867.73547	0.00000\\
935871.74349	0.00000\\
937875.75150	0.00000\\
939879.75952	0.00000\\
941883.76754	0.00000\\
943887.77555	0.00000\\
945891.78357	0.00000\\
947895.79158	0.00000\\
949899.79960	0.00000\\
951903.80762	0.00000\\
953907.81563	0.00000\\
955911.82365	0.00000\\
957915.83166	0.00000\\
959919.83968	0.00000\\
961923.84770	0.00000\\
963927.85571	0.00000\\
965931.86373	0.00000\\
967935.87174	0.00000\\
969939.87976	0.00000\\
971943.88778	0.00000\\
973947.89579	0.00000\\
975951.90381	0.00000\\
977955.91182	0.00000\\
979959.91984	0.00000\\
981963.92786	0.00000\\
983967.93587	0.00000\\
985971.94389	0.00000\\
987975.95190	0.00000\\
989979.95992	0.00000\\
991983.96794	0.00000\\
993987.97595	0.00000\\
995991.98397	0.00000\\
997995.99198	0.00000\\
1000000.00000	0.00000\\
};
\addplot [color=black, dashed, line width=1.0pt, forget plot]
  table[row sep=crcr]{%
0.00000	1.00000\\
2004.00802	0.99999\\
4008.01603	0.99995\\
6012.02405	0.99987\\
8016.03206	0.99977\\
10020.04008	0.99976\\
12024.04810	0.99981\\
14028.05611	0.99980\\
16032.06413	0.99978\\
18036.07214	0.99976\\
20040.08016	0.99974\\
22044.08818	0.99971\\
24048.09619	0.99968\\
26052.10421	0.99965\\
28056.11222	0.99961\\
30060.12024	0.99958\\
32064.12826	0.99954\\
34068.13627	0.99950\\
36072.14429	0.99945\\
38076.15230	0.99941\\
40080.16032	0.99937\\
42084.16834	0.99932\\
44088.17635	0.99927\\
46092.18437	0.99922\\
48096.19238	0.99916\\
50100.20040	0.99911\\
52104.20842	0.99905\\
54108.21643	0.99900\\
56112.22445	0.99894\\
58116.23246	0.99887\\
60120.24048	0.99881\\
62124.24850	0.99874\\
64128.25651	0.99868\\
66132.26453	0.99861\\
68136.27255	0.99854\\
70140.28056	0.99846\\
72144.28858	0.99839\\
74148.29659	0.99831\\
76152.30461	0.99824\\
78156.31263	0.99816\\
80160.32064	0.99807\\
82164.32866	0.99799\\
84168.33667	0.99791\\
86172.34469	0.99782\\
88176.35271	0.99773\\
90180.36072	0.99764\\
92184.36874	0.99755\\
94188.37675	0.99745\\
96192.38477	0.99735\\
98196.39279	0.99726\\
100200.40080	0.99716\\
102204.40882	0.99705\\
104208.41683	0.99695\\
106212.42485	0.99684\\
108216.43287	0.99674\\
110220.44088	0.99663\\
112224.44890	0.99652\\
114228.45691	0.99640\\
116232.46493	0.99629\\
118236.47295	0.99617\\
120240.48096	0.99605\\
122244.48898	0.99593\\
124248.49699	0.99581\\
126252.50501	0.99568\\
128256.51303	0.99555\\
130260.52104	0.99542\\
132264.52906	0.99529\\
134268.53707	0.99516\\
136272.54509	0.99503\\
138276.55311	0.99489\\
140280.56112	0.99475\\
142284.56914	0.99461\\
144288.57715	0.99447\\
146292.58517	0.99432\\
148296.59319	0.99418\\
150300.60120	0.99403\\
152304.60922	0.99388\\
154308.61723	0.99373\\
156312.62525	0.99357\\
158316.63327	0.99341\\
160320.64128	0.99326\\
162324.64930	0.99309\\
164328.65731	0.99293\\
166332.66533	0.99277\\
168336.67335	0.99260\\
170340.68136	0.99243\\
172344.68938	0.99226\\
174348.69739	0.99209\\
176352.70541	0.99191\\
178356.71343	0.99174\\
180360.72144	0.99156\\
182364.72946	0.99138\\
184368.73747	0.99119\\
186372.74549	0.99101\\
188376.75351	0.99082\\
190380.76152	0.99063\\
192384.76954	0.99044\\
194388.77756	0.99025\\
196392.78557	0.99005\\
198396.79359	0.98985\\
200400.80160	0.98965\\
202404.80962	0.98945\\
204408.81764	0.98925\\
206412.82565	0.98904\\
208416.83367	0.98883\\
210420.84168	0.98862\\
212424.84970	0.98841\\
214428.85772	0.98820\\
216432.86573	0.98798\\
218436.87375	0.98776\\
220440.88176	0.98754\\
222444.88978	0.98732\\
224448.89780	0.98709\\
226452.90581	0.98687\\
228456.91383	0.98664\\
230460.92184	0.98640\\
232464.92986	0.98617\\
234468.93788	0.98593\\
236472.94589	0.98570\\
238476.95391	0.98546\\
240480.96192	0.98521\\
242484.96994	0.98497\\
244488.97796	0.98472\\
246492.98597	0.98447\\
248496.99399	0.98422\\
250501.00200	0.98397\\
252505.01002	0.98371\\
254509.01804	0.98345\\
256513.02605	0.98319\\
258517.03407	0.98293\\
260521.04208	0.98267\\
262525.05010	0.98240\\
264529.05812	0.98213\\
266533.06613	0.98186\\
268537.07415	0.98159\\
270541.08216	0.98131\\
272545.09018	0.98103\\
274549.09820	0.98075\\
276553.10621	0.98047\\
278557.11423	0.98019\\
280561.12224	0.97990\\
282565.13026	0.97961\\
284569.13828	0.97932\\
286573.14629	0.97903\\
288577.15431	0.97873\\
290581.16232	0.97843\\
292585.17034	0.97813\\
294589.17836	0.97783\\
296593.18637	0.97752\\
298597.19439	0.97722\\
300601.20240	0.97691\\
302605.21042	0.97660\\
304609.21844	0.97628\\
306613.22645	0.97597\\
308617.23447	0.97565\\
310621.24248	0.97533\\
312625.25050	0.97500\\
314629.25852	0.97468\\
316633.26653	0.97435\\
318637.27455	0.97402\\
320641.28257	0.97369\\
322645.29058	0.97335\\
324649.29860	0.97302\\
326653.30661	0.97268\\
328657.31463	0.97234\\
330661.32265	0.97199\\
332665.33066	0.97165\\
334669.33868	0.97130\\
336673.34669	0.97095\\
338677.35471	0.97059\\
340681.36273	0.97024\\
342685.37074	0.96988\\
344689.37876	0.96952\\
346693.38677	0.96916\\
348697.39479	0.96879\\
350701.40281	0.96843\\
352705.41082	0.96806\\
354709.41884	0.96769\\
356713.42685	0.96731\\
358717.43487	0.96694\\
360721.44289	0.96656\\
362725.45090	0.96618\\
364729.45892	0.96579\\
366733.46693	0.96541\\
368737.47495	0.96502\\
370741.48297	0.96463\\
372745.49098	0.96424\\
374749.49900	0.96384\\
376753.50701	0.96344\\
378757.51503	0.96304\\
380761.52305	0.96264\\
382765.53106	0.96224\\
384769.53908	0.96183\\
386773.54709	0.96142\\
388777.55511	0.96101\\
390781.56313	0.96059\\
392785.57114	0.96018\\
394789.57916	0.95976\\
396793.58717	0.95934\\
398797.59519	0.95891\\
400801.60321	0.95849\\
402805.61122	0.95806\\
404809.61924	0.95763\\
406813.62725	0.95719\\
408817.63527	0.95676\\
410821.64329	0.95632\\
412825.65130	0.95588\\
414829.65932	0.95544\\
416833.66733	0.95499\\
418837.67535	0.95454\\
420841.68337	0.95409\\
422845.69138	0.95364\\
424849.69940	0.95319\\
426853.70741	0.95273\\
428857.71543	0.95227\\
430861.72345	0.95181\\
432865.73146	0.95134\\
434869.73948	0.95088\\
436873.74749	0.95041\\
438877.75551	0.94994\\
440881.76353	0.94946\\
442885.77154	0.94899\\
444889.77956	0.94851\\
446893.78758	0.94803\\
448897.79559	0.94754\\
450901.80361	0.94706\\
452905.81162	0.94657\\
454909.81964	0.94608\\
456913.82766	0.94558\\
458917.83567	0.94509\\
460921.84369	0.94459\\
462925.85170	0.94409\\
464929.85972	0.94359\\
466933.86774	0.94308\\
468937.87575	0.94257\\
470941.88377	0.94206\\
472945.89178	0.94155\\
474949.89980	0.94104\\
476953.90782	0.94052\\
478957.91583	0.94000\\
480961.92385	0.93948\\
482965.93186	0.93895\\
484969.93988	0.93843\\
486973.94790	0.93790\\
488977.95591	0.93736\\
490981.96393	0.93683\\
492985.97194	0.93629\\
494989.97996	0.93575\\
496993.98798	0.93521\\
498997.99599	0.93467\\
501002.00401	0.93412\\
503006.01202	0.93357\\
505010.02004	0.93302\\
507014.02806	0.93247\\
509018.03607	0.93191\\
511022.04409	0.93135\\
513026.05210	0.93079\\
515030.06012	0.93023\\
517034.06814	0.92966\\
519038.07615	0.92910\\
521042.08417	0.92853\\
523046.09218	0.92795\\
525050.10020	0.92738\\
527054.10822	0.92680\\
529058.11623	0.92622\\
531062.12425	0.92564\\
533066.13226	0.92505\\
535070.14028	0.92446\\
537074.14830	0.92387\\
539078.15631	0.92328\\
541082.16433	0.92269\\
543086.17234	0.92209\\
545090.18036	0.92149\\
547094.18838	0.92089\\
549098.19639	0.92028\\
551102.20441	0.91968\\
553106.21242	0.91907\\
555110.22044	0.91846\\
557114.22846	0.91784\\
559118.23647	0.91723\\
561122.24449	0.91661\\
563126.25251	0.91599\\
565130.26052	0.91536\\
567134.26854	0.91474\\
569138.27655	0.91411\\
571142.28457	0.91348\\
573146.29259	0.91285\\
575150.30060	0.91221\\
577154.30862	0.91157\\
579158.31663	0.91093\\
581162.32465	0.91029\\
583166.33267	0.90964\\
585170.34068	0.90900\\
587174.34870	0.90835\\
589178.35671	0.90769\\
591182.36473	0.90704\\
593186.37275	0.90638\\
595190.38076	0.90572\\
597194.38878	0.90506\\
599198.39679	0.90440\\
601202.40481	0.90373\\
603206.41283	0.90306\\
605210.42084	0.90239\\
607214.42886	0.90172\\
609218.43687	0.90104\\
611222.44489	0.90036\\
613226.45291	0.89968\\
615230.46092	0.89900\\
617234.46894	0.89831\\
619238.47695	0.89762\\
621242.48497	0.89693\\
623246.49299	0.89624\\
625250.50100	0.89555\\
627254.50902	0.89485\\
629258.51703	0.89415\\
631262.52505	0.89345\\
633266.53307	0.89274\\
635270.54108	0.89204\\
637274.54910	0.89133\\
639278.55711	0.89062\\
641282.56513	0.88990\\
643286.57315	0.88919\\
645290.58116	0.88847\\
647294.58918	0.88775\\
649298.59719	0.88702\\
651302.60521	0.88630\\
653306.61323	0.88557\\
655310.62124	0.88484\\
657314.62926	0.88411\\
659318.63727	0.88337\\
661322.64529	0.88264\\
663326.65331	0.88190\\
665330.66132	0.88116\\
667334.66934	0.88041\\
669338.67735	0.87967\\
671342.68537	0.87892\\
673346.69339	0.87817\\
675350.70140	0.87741\\
677354.70942	0.87666\\
679358.71743	0.87590\\
681362.72545	0.87514\\
683366.73347	0.87438\\
685370.74148	0.87361\\
687374.74950	0.87285\\
689378.75752	0.87208\\
691382.76553	0.87131\\
693386.77355	0.87053\\
695390.78156	0.86976\\
697394.78958	0.86898\\
699398.79760	0.86820\\
701402.80561	0.86742\\
703406.81363	0.86663\\
705410.82164	0.86584\\
707414.82966	0.86505\\
709418.83768	0.86426\\
711422.84569	0.86347\\
713426.85371	0.86267\\
715430.86172	0.86187\\
717434.86974	0.86107\\
719438.87776	0.86027\\
721442.88577	0.85946\\
723446.89379	0.85866\\
725450.90180	0.85785\\
727454.90982	0.85704\\
729458.91784	0.85622\\
731462.92585	0.85541\\
733466.93387	0.85459\\
735470.94188	0.85377\\
737474.94990	0.85294\\
739478.95792	0.85212\\
741482.96593	0.85129\\
743486.97395	0.85046\\
745490.98196	0.84963\\
747494.98998	0.84880\\
749498.99800	0.84796\\
751503.00601	0.84712\\
753507.01403	0.84628\\
755511.02204	0.84544\\
757515.03006	0.84460\\
759519.03808	0.84375\\
761523.04609	0.84290\\
763527.05411	0.84205\\
765531.06212	0.84120\\
767535.07014	0.84034\\
769539.07816	0.83949\\
771543.08617	0.83863\\
773547.09419	0.83776\\
775551.10220	0.83690\\
777555.11022	0.83604\\
779559.11824	0.83517\\
781563.12625	0.83430\\
783567.13427	0.83343\\
785571.14228	0.83255\\
787575.15030	0.83167\\
789579.15832	0.83080\\
791583.16633	0.82992\\
793587.17435	0.82903\\
795591.18236	0.82815\\
797595.19038	0.82726\\
799599.19840	0.82637\\
801603.20641	0.82548\\
803607.21443	0.82459\\
805611.22244	0.82370\\
807615.23046	0.82280\\
809619.23848	0.82190\\
811623.24649	0.82100\\
813627.25451	0.82010\\
815631.26253	0.81919\\
817635.27054	0.81828\\
819639.27856	0.81737\\
821643.28657	0.81646\\
823647.29459	0.81555\\
825651.30261	0.81464\\
827655.31062	0.81372\\
829659.31864	0.81280\\
831663.32665	0.81188\\
833667.33467	0.81096\\
835671.34269	0.81003\\
837675.35070	0.80910\\
839679.35872	0.80817\\
841683.36673	0.80724\\
843687.37475	0.80631\\
845691.38277	0.80538\\
847695.39078	0.80444\\
849699.39880	0.80350\\
851703.40681	0.80256\\
853707.41483	0.80162\\
855711.42285	0.80067\\
857715.43086	0.79973\\
859719.43888	0.79878\\
861723.44689	0.79783\\
863727.45491	0.79688\\
865731.46293	0.79592\\
867735.47094	0.79497\\
869739.47896	0.79401\\
871743.48697	0.79305\\
873747.49499	0.79209\\
875751.50301	0.79113\\
877755.51102	0.79016\\
879759.51904	0.78919\\
881763.52705	0.78822\\
883767.53507	0.78725\\
885771.54309	0.78628\\
887775.55110	0.78531\\
889779.55912	0.78433\\
891783.56713	0.78335\\
893787.57515	0.78237\\
895791.58317	0.78139\\
897795.59118	0.78041\\
899799.59920	0.77942\\
901803.60721	0.77844\\
903807.61523	0.77745\\
905811.62325	0.77646\\
907815.63126	0.77547\\
909819.63928	0.77447\\
911823.64729	0.77348\\
913827.65531	0.77248\\
915831.66333	0.77148\\
917835.67134	0.77048\\
919839.67936	0.76948\\
921843.68737	0.76847\\
923847.69539	0.76747\\
925851.70341	0.76646\\
927855.71142	0.76545\\
929859.71944	0.76444\\
931863.72745	0.76342\\
933867.73547	0.76241\\
935871.74349	0.76139\\
937875.75150	0.76038\\
939879.75952	0.75936\\
941883.76754	0.75833\\
943887.77555	0.75731\\
945891.78357	0.75629\\
947895.79158	0.75526\\
949899.79960	0.75423\\
951903.80762	0.75320\\
953907.81563	0.75217\\
955911.82365	0.75114\\
957915.83166	0.75011\\
959919.83968	0.74907\\
961923.84770	0.74803\\
963927.85571	0.74699\\
965931.86373	0.74595\\
967935.87174	0.74491\\
969939.87976	0.74387\\
971943.88778	0.74282\\
973947.89579	0.74178\\
975951.90381	0.74073\\
977955.91182	0.73968\\
979959.91984	0.73863\\
981963.92786	0.73758\\
983967.93587	0.73652\\
985971.94389	0.73547\\
987975.95190	0.73441\\
989979.95992	0.73335\\
991983.96794	0.73229\\
993987.97595	0.73123\\
995991.98397	0.73016\\
997995.99198	0.72910\\
1000000.00000	0.72803\\
};
\end{axis}
\end{tikzpicture}%
    \caption{Empirische Lebensdauerprognose am Feldpunkt. Das \ac{EmALT}-Design liefert nicht nur eine radikal verkleinerte Prädiktionsvarianz (schmales Konfidenzband), sondern prognostiziert zudem ein konservativeres, realitätsnäheres Ausfallverhalten als das ungestützte Basis-\ac{CCD}.}
    \label{fig:ma_abb6.1_reliability_ci}
\end{figure}

Die Auswertung bestätigt zunächst eindrucksvoll die theoriegeleitete Erwartung zur Varianzreduktion aus Kapitel~\ref{chap:entwurf}:
Das Basis-\ac{CCD} weist aufgrund der ungestützten Extrapolationsdistanz eine praktisch inakzeptable Prädiktionsvarianz von $\sym{SPV} = 8{,}24$ auf.
Dies führt zu einem stark ausgeprägten Konfidenzband, das bei einer beispielhaften Laufzeit von $5\times 10^5\text{ h}$ eine extreme Unschärfe von $14{,}8\,\%$ aufspannt.
Das \ac{EmALT}-Design hingegen reduziert die Varianz durch die implementierten \acp{AP} massiv: Die \ac{SPV} kollabiert auf $0{,}15$, wodurch sich die Breite des Konfidenzintervalls am Referenzpunkt auf hochpräzise $1{,}3\,\%$ verdichtet.
Die Extrapolationsunsicherheit wird somit um einen Faktor von $11{,}2$ minimiert.

Neben dieser Varianzdämpfung offenbart Abbildung~\ref{fig:ma_abb6.1_reliability_ci} jedoch ein noch kritischeres Phänomen: Die mittlere Prognosekurve (Punktschätzer) des \ac{EmALT}-Designs verläuft signifikant flacher und konservativer als jene des Basis-\ac{CCD}.
Diese drastische Abweichung wirft die essenzielle suggestive Validierungsfrage auf: Decken die \acp{AP} hier lediglich eine Präzisierung der Vorhersageunschärfe auf, oder offenbaren sie eine fundamentale Modellabweichung infolge wirtschaftlich "gut gemeinter", aber stochastisch fehlerhafter Beschleunigung (\ac{ALT})?
Haben die Ankerpunkte das Basis-Modell offiziell widerlegt?

\subsection{Confirmation Experiments nach Stevens}

Um diese Frage mathematisch zweifelsfrei zu beantworten, wird das Basis-\ac{CCD} anhand der real getesteten Ankerpunkte mittels \textit{prädiktionsintervall-basierter} Analysemethoden falsifiziert.
Hierfür fungieren die \acp{AP} als formale \textit{Confirmation Runs} \cite{Stevens.2019}.
Die Logik ist zwingend: Wenn das durch das Basis-\ac{CCD} im Hochlastbereich extrahierte Modell die physikalische Realität korrekt beschreibt, müssen die realen Ausfallzeiten der fernen Ankerpunkte innerhalb des vom \ac{CCD} tolerierten Streubandes liegen.

Dazu wird für die \ac{AP}-Koordinaten das Prädiktionsintervall (\ac{PI}) des Basis-\ac{CCD} berechnet.
Da die Ausfallzeiten Weibull-verteilt sind, definiert sich das \ac{PI} über das $5\,\%$- und $95\,\%$-Quantil der geschätzten Verteilung, auf welches additiv die epistemische Parameterunsicherheit über die $\Delta$-Methode der \ac{FIM} angewandt wird.
Fällt die reale Lebensdauer aus diesem massiv aufgeweiteten Toleranzband heraus, ist die Abweichung statistisch signifikant und das Basis-\ac{CCD} offiziell widerlegt.

Die Auswertung der realen Prüfstandsdaten liefert folgendes Resultat:
\begin{itemize}
    \item \textbf{Ankerpunkt 1 $(-3{,}15 ; -2{,}56)$:}
          Das $90\,\%$-\ac{PI} des Basis-\ac{CCD} prognostiziert für diese Koordinate eine ausgedehnte Streuung von $[24{,}8\times 10^6 ; 38\,426{,}9\times 10^6]\text{ h}$.
          Die vier realen Replikationsläufe fielen jedoch bereits bei $\approx 3{,}95$ bis $4{,}30\times 10^6\text{ h}$ aus. Sämtliche Messwerte unterschreiten das vom \ac{CCD} erlaubte Minimum massiv.
    \item \textbf{Ankerpunkt 2 $(-2{,}42 ; -2{,}17)$:}
          Das \ac{PI} des Basis-\ac{CCD} spannt einen Bereich von $[2{,}90\times 10^6 ; 542{,}3\times 10^6]\text{ h}$ auf.
          Die realen Versuchsläufe versagten bei $1{,}85\times 10^6\text{ h}$ und $1{,}91\times 10^6\text{ h}$. Auch hier liegen die Laufzeiten hochgradig signifikant außerhalb des theoretisch zulässigen Toleranzbandes.
\end{itemize}

Dieses Ergebnis erbringt den ultimativen empirischen Beweis: Das voll-orthogonale Basis-\ac{CCD} ist an den realen Daten gescheitert.
Es extrapoliert die Lebensdauer unkonservativ und überschätzt die reale Bauteilfestigkeit um ein Vielfaches.

\subsection{Ökonomische Schlussfolgerung}

Die Einbindung der \acp{AP} erzeugt logischerweise eine Verschiebung in der \ac{EoL}-Kostengleichung.
Während das Basis-\ac{CCD} mit aggregierten Gesamtkosten von $45{,}2\text{ kCU}$ extrem wirtschaftlich erscheint, beansprucht das \ac{EmALT}-Design $344{,}6\text{ kCU}$ -- wenngleich ein reines Feld-\ac{CCD} zur Erreichung vergleichbarer Varianzen überschlägig ein Vielfaches dieses Aufwands fordern würde.
Die Falsifikation durch die \textit{Confirmation Runs} beantwortet die Frage nach der Rechtfertigung dieses monetären Mehraufwands jedoch unmissverständlich.
Die Re-Allokation des Gestaltungsbudgets resultiert nicht einfach in einem "akademisch präziseren" Modell mit feineren Konfidenzgrenzen. Der Eingriff ist eine zwingende methodische Notwendigkeit.
Ohne die durch das \ac{EmALT}-Design positionierten Ankerpunkte als physikalische Frühwarnsensoren wäre ein fundamental falsches, lebensgefährlich überschätztes Zuverlässigkeitsmodell für den Lenksystem-Zahnriemen freigegeben worden.

Die adaptive Nutzung des Gestaltungsbudgets legitimiert sich somit nicht allein aus informationstheoretischer Optimierung, sondern dient vielmehr als unabdingbares methodisches Tuning und als Garantiewerkzeug, um fatale Extrapolationsfehler in der industriellen Praxis abzuwenden.